% Reproduction — Biology Upper Sixth — Cours signé OnBuch+
% Official MINESEC syllabus. Compiler avec : tectonic BIO01-reproduction.tex
\newcommand{\DOCMATIERE}{Biology}
\newcommand{\DOCNIVEAU}{Upper Sixth}
\newcommand{\DOCMODULE}{Upper Sixth — Biology}
\newcommand{\DOCLECON}{1}
\newcommand{\DOCTITRE}{Reproduction}
% OnBuch+ — préambule « cours signés » ENRICHI (compilé avec tectonic / XeLaTeX)
% Système visuel « Pop » d'OnBuch+ : fond crème (jamais blanc), bordures encre
% épaisses, ombres « dures » décalées, titres Archivo Black en capitales.
% Version MATHÉMATIQUES : graphes (pgfplots/tikz), boîtes pédagogiques
% (définition, propriété, méthode, attention, le savais-tu, exercice résolu,
% exercice, TP, bilan), page de garde et signature OnBuch+ ancrée en bas.
%
% Macros attendues dans main.tex AVANT \input{preamble} :
%   \DOCMATIERE  \DOCNIVEAU  \DOCMODULE  \DOCLECON (numéro)  \DOCTITRE
\documentclass[11pt]{article}

\usepackage{fontspec}
\usepackage[english]{babel}
\usepackage[a4paper,margin=2cm,top=3.4cm,bottom=2.8cm,headheight=1.4cm,footskip=1.3cm]{geometry}
\usepackage{amsmath,amssymb,amsfonts}
\usepackage{mathtools} % \xmapsto, \coloneqq... souvent utilisés par les modèles
\usepackage{cancel} % simplifications barrées (bilans de forces)
% \abs{z} (module, valeur absolue) et \norm{u} (norme), utilisés par les modèles
\providecommand{\abs}{}\let\abs\relax\DeclarePairedDelimiter{\abs}{\lvert}{\rvert}
\providecommand{\norm}{}\let\norm\relax\DeclarePairedDelimiter{\norm}{\lVert}{\rVert}
\usepackage[table]{xcolor} % [table] : fournit aussi \rowcolors (alternance de lignes)
\usepackage{pagecolor}
\usepackage{tikz}
\usetikzlibrary{arrows.meta,positioning,calc,shapes.geometric,shapes.misc,shapes.arrows,decorations.pathmorphing,decorations.pathreplacing,decorations.markings,patterns,shadows,mindmap,trees,fit,backgrounds,matrix,intersections,angles,quotes}
\usepackage{pgfplots}
\usepackage[version=4]{mhchem} % équations nucléaires \ce{^{235}_{92}U}
\usepackage[european,siunitx]{circuitikz} % schémas électriques
\pgfplotsset{compat=1.18}
\usepgfplotslibrary{fillbetween,groupplots} % aires entre courbes (calcul intégral)
\usepackage{enumitem}
\usepackage{fancyhdr}
% [most] charge toutes les bibliothèques tcolorbox (listings, minted, raster,
% documentation...) et gonfle inutilement la mémoire TeX sur un document long
% (~300 boîtes) : on ne charge que ce qui est réellement utilisé.
\usepackage[skins,breakable]{tcolorbox}
\usepackage{graphicx}
\usepackage{colortbl}
\usepackage{booktabs}
\usepackage{tabularx}
\usepackage{diagbox} % case d'en-tête diagonale des tableaux à double entrée
% Colonnes extensibles centrée / gauche / droite (C, L, R), souvent utilisées par les modèles
\newcolumntype{C}{>{\centering\arraybackslash}X}
\newcolumntype{L}{>{\raggedright\arraybackslash}X}
\newcolumntype{R}{>{\raggedleft\arraybackslash}X}
\usepackage{array}
\usepackage{multicol}
\usepackage{siunitx}
% parse-numbers=false : accepte aussi \SI{2,5 \times 10^{-3}}{...}
\sisetup{locale=UK,output-decimal-marker={.},group-minimum-digits=4,parse-numbers=false}
\DeclareSIUnit\ohm{\ensuremath{\symup{\Omega}}} % U+2126 absent de la police
\DeclareSIUnit\division{div} % réglages d'oscilloscope (ms/div, V/div)
\DeclareSIUnit\tr{tr} % tours (vitesse de rotation, ex. \SI{1500}{\tr\per\minute})
\usepackage{qrcode}
% \degree : commande fréquemment utilisée par les modèles pour un angle ou une
% température en dehors de \SI{}{} (non fournie par les paquets chargés).
\providecommand{\degree}{^{\circ}}
% \qed (fin de démonstration), utilisé par les modèles sans amsthm
\providecommand{\qed}{\hfill$\square$}
% \barre : double barre des valeurs interdites dans un tableau de variations (tkz-tab)
\providecommand{\barre}{\ensuremath{\|}}
% environnement proof (amsthm non chargé) utilisé par les modèles
\ifdefined\proof\else\newenvironment{proof}[1][Proof]{\par\noindent\textit{#1.}\ }{\qed\par}\fi
\usepackage{lastpage}
\usepackage{hyperref}
\hypersetup{hidelinks}

% ============================================================
% Palette « Pop » OnBuch+ — mêmes hex que la classe P (plus_ui.dart)
% ============================================================
\definecolor{poporange}{HTML}{F59321}
\definecolor{popdark}{HTML}{E07A0C}
\definecolor{popcream}{HTML}{FFF6E8}
\definecolor{popink}{HTML}{1C1714}
\definecolor{poppurple}{HTML}{7A5AE0}
\definecolor{popgreen}{HTML}{2E9E6A}
\definecolor{poppink}{HTML}{E0457B}
\definecolor{popblue}{HTML}{2F7DE1}
\definecolor{popgold}{HTML}{C98A12}
\definecolor{popmuted}{HTML}{8A7F76}
\definecolor{popline}{HTML}{EADFD2}
\definecolor{popgray}{HTML}{8A7F76}
\colorlet{popgrey}{popgray}
% Alias tolérants (le modèle nomme parfois les couleurs différemment)
\colorlet{popwhite}{white}
\colorlet{popblack}{popink}
\colorlet{popred}{poppink}
\colorlet{popyellow}{popgold}
% Teintes claires pour fonds de tableaux / zones
\colorlet{popblueL}{popblue!10!white}
\colorlet{poporangeL}{poporange!14!white}
\colorlet{popgreenL}{popgreen!12!white}
\colorlet{poppurpleL}{poppurple!12!white}
\colorlet{poppinkL}{poppink!12!white}
\colorlet{popgoldL}{popgold!14!white}

\pagecolor{popcream}

% ============================================================
% Typographie
% ============================================================
\defaultfontfeatures{Ligatures=TeX}
\setmainfont{PlusJakartaSans-Medium.ttf}[Path=../../../fonts/,BoldFont=PlusJakartaSans-Bold.ttf]
% Maths en sans-serif (Fira Math) avec chiffres et lettres droites de Plus
% Jakarta, pour que les formules se fondent dans le texte.
\usepackage[math-style=ISO,bold-style=ISO]{unicode-math}
\setmathfont{FiraMath-Regular.otf}[Path=../../../fonts/]
\setmathfont{PlusJakartaSans-Medium.ttf}[Path=../../../fonts/,range={up/num,up/{Latin,latin},\mathup}]
% (icomma retiré : décimales avec point en anglais)
% Caractères mathématiques Unicode absents des polices de texte (Plus Jakarta,
% Archivo Black) : rendus par leur équivalent mathématique, en texte comme en
% formule — sinon ils s'affichent en carré vide (ex. « Arithmétique dans ℤ »).
\usepackage{newunicodechar}
\newunicodechar{ℝ}{\ensuremath{\mathbb{R}}}
\newunicodechar{ℤ}{\ensuremath{\mathbb{Z}}}
\newunicodechar{ℂ}{\ensuremath{\mathbb{C}}}
\newunicodechar{ℕ}{\ensuremath{\mathbb{N}}}
\newunicodechar{ℚ}{\ensuremath{\mathbb{Q}}}
\newunicodechar{𝔻}{\ensuremath{\mathbb{D}}}
\newunicodechar{α}{\ensuremath{\alpha}}
\newunicodechar{β}{\ensuremath{\beta}}
\newunicodechar{γ}{\ensuremath{\gamma}}
\newunicodechar{δ}{\ensuremath{\delta}}
\newunicodechar{ε}{\ensuremath{\varepsilon}}
\newunicodechar{θ}{\ensuremath{\theta}}
\newunicodechar{λ}{\ensuremath{\lambda}}
\newunicodechar{μ}{\ensuremath{\mu}}
\newunicodechar{σ}{\ensuremath{\sigma}}
\newunicodechar{τ}{\ensuremath{\tau}}
\newunicodechar{φ}{\ensuremath{\varphi}}
\newunicodechar{ψ}{\ensuremath{\psi}}
\newunicodechar{ω}{\ensuremath{\omega}}
\newunicodechar{Σ}{\ensuremath{\Sigma}}
% majuscules produites par \MakeUppercase dans les titres (α-aminés -> Α)
\newunicodechar{Α}{\ensuremath{\alpha}}
\newunicodechar{Β}{\ensuremath{\beta}}
\newunicodechar{Γ}{\ensuremath{\Gamma}}
\newunicodechar{Δ}{\ensuremath{\Delta}}
\newunicodechar{Ε}{\ensuremath{\varepsilon}}
\newunicodechar{Θ}{\ensuremath{\Theta}}
\newunicodechar{Λ}{\ensuremath{\Lambda}}
\newunicodechar{Μ}{\ensuremath{\mu}}
\newunicodechar{Τ}{\ensuremath{\tau}}
\newunicodechar{Φ}{\ensuremath{\Phi}}
\newunicodechar{Ψ}{\ensuremath{\Psi}}
\newunicodechar{Ω}{\ensuremath{\Omega}}
\newunicodechar{↦}{\ensuremath{\mapsto}}
\newunicodechar{∈}{\ensuremath{\in}}
\newunicodechar{∉}{\ensuremath{\notin}}
\newunicodechar{∧}{\ensuremath{\wedge}}
\newunicodechar{⇔}{\ensuremath{\Leftrightarrow}}
\newunicodechar{⟺}{\ensuremath{\Longleftrightarrow}}
\newunicodechar{⇒}{\ensuremath{\Rightarrow}}
\newunicodechar{⟹}{\ensuremath{\Longrightarrow}}
\newunicodechar{∘}{\ensuremath{\circ}}
\newunicodechar{∪}{\ensuremath{\cup}}
\newunicodechar{∩}{\ensuremath{\cap}}
\newunicodechar{≡}{\ensuremath{\equiv}}
\newunicodechar{⊂}{\ensuremath{\subset}}
\newunicodechar{⊥}{\ensuremath{\perp}}
\newunicodechar{∃}{\ensuremath{\exists}}
\newunicodechar{∀}{\ensuremath{\forall}}
\newunicodechar{∛}{\ensuremath{\sqrt[3]{\,}}}
\newunicodechar{∜}{\ensuremath{\sqrt[4]{\,}}}
\newunicodechar{⃗}{\ensuremath{{}^{\to}}}
\newunicodechar{ⁿ}{\ifmmode^{n}\else\textsuperscript{\ensuremath{n}}\fi}
\newunicodechar{ˣ}{\ifmmode^{x}\else\textsuperscript{\ensuremath{x}}\fi}
\newunicodechar{ᵇ}{\ifmmode^{b}\else\textsuperscript{\ensuremath{b}}\fi}
\newunicodechar{ʳ}{\ifmmode^{r}\else\textsuperscript{\ensuremath{r}}\fi}
\newunicodechar{ᵘ}{\ifmmode^{u}\else\textsuperscript{\ensuremath{u}}\fi}
\newunicodechar{ʸ}{\ifmmode^{y}\else\textsuperscript{\ensuremath{y}}\fi}
\newunicodechar{ᵃ}{\ifmmode^{a}\else\textsuperscript{\ensuremath{a}}\fi}
\newunicodechar{ᵉ}{\ifmmode^{e}\else\textsuperscript{\ensuremath{e}}\fi}
\newunicodechar{ᵏ}{\ifmmode^{k}\else\textsuperscript{\ensuremath{k}}\fi}
\newunicodechar{ᵖ}{\ifmmode^{p}\else\textsuperscript{\ensuremath{p}}\fi}
\newunicodechar{ᴺ}{\ifmmode^{N}\else\textsuperscript{\ensuremath{N}}\fi}
\newunicodechar{ᶜ}{\ifmmode^{c}\else\textsuperscript{\ensuremath{c}}\fi}
\newunicodechar{ʰ}{\ifmmode^{h}\else\textsuperscript{\ensuremath{h}}\fi}
\newunicodechar{ᵠ}{\ifmmode^{q}\else\textsuperscript{\ensuremath{q}}\fi}
\newunicodechar{ₙ}{\ifmmode_{n}\else\textsubscript{\ensuremath{n}}\fi}
\newunicodechar{ₐ}{\ifmmode_{a}\else\textsubscript{\ensuremath{a}}\fi}
\newunicodechar{ᵢ}{\ifmmode_{i}\else\textsubscript{\ensuremath{i}}\fi}
\newunicodechar{ᵣ}{\ifmmode_{r}\else\textsubscript{\ensuremath{r}}\fi}
\newunicodechar{ₖ}{\ifmmode_{k}\else\textsubscript{\ensuremath{k}}\fi}
\newunicodechar{ᵦ}{\ifmmode_{\beta}\else\textsubscript{\ensuremath{\beta}}\fi}
\newunicodechar{ₓ}{\ifmmode_{x}\else\textsubscript{\ensuremath{x}}\fi}
\newfontfamily\popdisplay{ArchivoBlack-Regular.ttf}[Path=../../../fonts/]
\newfontfamily\poplogo{SpaceGrotesk-Bold.ttf}[Path=../../../fonts/]
\newfontfamily\popsemi{PlusJakartaSans-SemiBold.ttf}[Path=../../../fonts/]
\newfontfamily\popxbold{PlusJakartaSans-ExtraBold.ttf}[Path=../../../fonts/]
\newfontfamily\popxboldls{PlusJakartaSans-ExtraBold.ttf}[Path=../../../fonts/,LetterSpace=3.0]

\setlist[enumerate]{leftmargin=1.7em,itemsep=5pt,topsep=4pt}
\setlist[itemize]{leftmargin=1.7em,itemsep=4pt,topsep=4pt,label={\color{poporange}\textbullet}}
\setlength{\parindent}{0pt}
\setlength{\parskip}{5pt}
\renewcommand{\arraystretch}{1.25}

% pgfplots : style Pop par défaut
\pgfplotsset{
  every axis/.append style={axis line style={popink,line width=1pt},tick style={popink},
    grid style={popline,line width=0.5pt},label style={font=\small},tick label style={font=\footnotesize}},
  popcurve/.style={poporange,line width=1.8pt,no markers,smooth},
}
% Même style disponible dans un \draw TikZ simple (hors pgfplots)
\tikzset{popcurve/.style={poporange,line width=1.8pt}}
\tikzset{popgrid/.style={popline,very thin}}
\colorlet{popcurve}{poporange} % le nom du style est parfois utilisé comme couleur

% ============================================================
% Pill / squiggle
% ============================================================
\newcommand{\poppill}[3]{%
  \tikz[baseline=-0.55ex]\node[draw=popink,line width=1.2pt,rounded corners=2.6pt,%
    fill=#2,inner xsep=6.5pt,inner ysep=3pt]{%
    \popxboldls\fontsize{7}{7}\selectfont\color{#3}\MakeUppercase{#1}};%
}
\newcommand{\popsquiggle}[1]{%
  \tikz[baseline=-0.4ex]{
    \draw[#1,line width=2.6pt,line cap=round]
      plot [smooth] coordinates {(0,0.09) (0.34,0.01) (0.68,0.21) (1.02,0.01) (1.36,0.10)};
  }%
}
% Mot-clé surligné (marqueur orange sous le texte)
\newcommand{\cle}[1]{{\popxbold\color{popdark}#1}}

% ============================================================
% En-tête / pied
% ============================================================
\pagestyle{fancy}
\fancyhf{}
\renewcommand{\headrulewidth}{0pt}
\renewcommand{\footrulewidth}{0pt}
\lhead{\raisebox{-0.15ex}{\poplogo\fontsize{13}{13}\selectfont\color{popink}OnBuch\color{poporange}+}}
\rhead{\poppill{\DOCMATIERE\ \textbullet\ \DOCNIVEAU\ \textbullet\ Chapter \DOCLECON}{white}{popink}}
\lfoot{\popxbold\fontsize{7.6}{7.6}\selectfont\color{popmuted}\MakeUppercase{Signed course OnBuch+}\ \textbullet\ {\popsemi\DOCTITRE}}
\rfoot{\tikz[baseline=-0.6ex]\node[rounded corners=8.5pt,fill=popink,minimum height=17pt,minimum width=17pt,inner xsep=5pt,inner sep=0]{\popxbold\fontsize{8}{8}\selectfont\color{popcream}\thepage\,/\,\pageref*{LastPage}};}

% ============================================================
% Titres
% ============================================================
\newcommand{\chaptitre}[2]{%
  \begin{tcolorbox}[enhanced,colback=poporange,colframe=popink,boxrule=2.6pt,arc=11pt,
      drop shadow=popink,left=15pt,right=15pt,top=12pt,bottom=11pt,before skip=3pt,after skip=18pt]
    \poppill{#2}{popink}{popcream}\par\vspace{9pt}
    {\raggedright\hyphenpenalty=10000\exhyphenpenalty=10000\popdisplay\fontsize{22}{24}\selectfont\color{popcream}\MakeUppercase{#1}\par}
  \end{tcolorbox}
}
% Section numérotée : pastille orange avec le numéro + titre Archivo
\newcounter{popsec}
\newcommand{\coursec}[1]{%
  \stepcounter{popsec}\setcounter{popsub}{0}%
  \par\vspace{16pt}\noindent
  \begin{minipage}{\linewidth}
  \tikz[baseline=-0.7ex]\node[draw=popink,line width=1.6pt,fill=poporange,rounded corners=4pt,minimum size=22pt,inner sep=0]{\popdisplay\fontsize{12}{12}\selectfont\color{popcream}\thepopsec};%
  \hspace{8pt}{\popdisplay\fontsize{15}{17}\selectfont\color{popink}\MakeUppercase{#1}}\par
  \vspace{4pt}\noindent{\color{poporange}\rule{36pt}{3.6pt}}
  \end{minipage}\par\nopagebreak\vspace{8pt}
}
\newcounter{popsub}
\newcommand{\courssub}[1]{%
  \stepcounter{popsub}%
  \par\vspace{9pt}\noindent{\popxbold\fontsize{11.5}{13}\selectfont\color{popink}{\color{poporange}\thepopsec.\thepopsub}\hspace{6pt}#1}\par\nopagebreak\vspace{4pt}
}

% ============================================================
% Boîtes pédagogiques — même recette : fond blanc, bordure épaisse colorée,
% ombre dure encre, badge plein. Titre optionnel [#1].
% ============================================================
\tcbset{popbase/.style={breakable,enhanced,colback=white,boxrule=2.3pt,arc=9pt,
  shadow={3pt}{-3pt}{0pt}{popink},left=14pt,right=14pt,top=11pt,bottom=11pt,before skip=11pt,after skip=15pt,
  fontupper=\popsemi\fontsize{10.6}{14.5}\selectfont\color{popink},
  fontlower=\popsemi\fontsize{10.6}{14.5}\selectfont\color{popink}}}
% Coche dessinée (le glyphe ✓ n'existe pas dans les polices du document)
\newcommand{\popcheck}[1]{\tikz[baseline=-0.5ex]\draw[#1,line width=1.6pt,line cap=round,line join=round](-0.13,0.0)--(-0.04,-0.09)--(0.13,0.1);}
\providecommand{\coche}{\popcheck{popgreen}}
\providecommand{\resultat}[1]{\par\smallskip\noindent\fcolorbox{popgreen}{popgreenL}{\parbox{\dimexpr\linewidth-2\fboxsep-2\fboxrule\relax}{\textbf{Result.} #1}}\par\smallskip}
\newcommand{\popboxhead}[3]{\poppill{#1}{#2}{white}\ifx\relax#3\relax\else\hspace{6pt}{\popxbold\fontsize{10.6}{12}\selectfont\color{popink}#3}\fi\par\vspace{7pt}}

\newtcolorbox{aretenir}[1][]{popbase,colframe=popgreen,before upper={\popboxhead{Key points}{popgreen}{#1}}}
\newtcolorbox{exemplebox}[1][]{popbase,colframe=poporange,before upper={\popboxhead{Exemple}{poporange}{#1}}}
\newtcolorbox{definition}[1][]{popbase,colframe=popblue,before upper={\popboxhead{Definition}{popblue}{#1}}}
\newtcolorbox{propriete}[1][]{popbase,colframe=poppurple,before upper={\popboxhead{Property}{poppurple}{#1}}}
\newtcolorbox{methode}[1][]{popbase,colframe=popgold,before upper={\popboxhead{Method}{popgold}{#1}}}
\newtcolorbox{attention}[1][]{popbase,colframe=poppink,before upper={\popboxhead{Warning}{poppink}{#1}}}
\newtcolorbox{experience}[1][]{popbase,colframe=popblue,colback=popblueL,before upper={\popboxhead{Experiment}{popblue}{#1}}}
% « Le savais-tu ? » : carte violette pleine (contexte, culture, Cameroun)
\newtcolorbox{savaistu}[1][]{popbase,colback=poppurple,colframe=popink,
  fontupper=\popsemi\fontsize{10.4}{14.2}\selectfont\color{white},
  before upper={\poppill{Did you know?}{white}{poppurple}\ifx\relax#1\relax\else\hspace{6pt}{\popxbold\color{white}#1}\fi\par\vspace{7pt}}}
% Exercice résolu : énoncé en haut, \tcblower puis la solution sur fond crème
\newcounter{popexo}\newcounter{popexr}
\newtcolorbox{exoresolu}[1][]{popbase,colframe=poporange,colbacklower=popcream,
  segmentation style={popink,line width=1.2pt,dashed},
  before upper={\stepcounter{popexr}\popboxhead{Worked example \thepopexr}{poporange}{#1}},
  before lower={\poppill{Solution}{popink}{popcream}\par\vspace{6pt}}}
% Exercice (non résolu en place) — la correction est dans la partie Corrigés
\newtcolorbox{exercice}[1][]{popbase,colframe=popink,
  before upper={\stepcounter{popexo}\popboxhead{Exercise \thepopexo}{popink}{#1}}}
% Corrigé
\newtcolorbox{corrige}[1][]{popbase,colframe=popgreen,colback=popgreenL,
  before upper={\popboxhead{Solution}{popgreen}{#1}}}
% Objectifs / prérequis / bilan
\newtcolorbox{objectifs}{popbase,colframe=popink,colback=poporangeL,
  before upper={\popboxhead{Chapter objectives}{popink}{}}}
\newtcolorbox{prerequis}{popbase,colframe=popink,colback=popblueL,
  before upper={\popboxhead{Prerequisites}{popblue}{}}}
\newtcolorbox{bilan}{popbase,colframe=popink,colback=white,boxrule=2.8pt,
  before upper={\popboxhead{Fiche bilan}{poporange}{L'essentiel à connaître pour l'examen}}}
% Figure encadrée avec légende
\newtcolorbox{popfigure}[1][]{enhanced,breakable=false,colback=white,colframe=popline,boxrule=1.4pt,arc=8pt,
  left=10pt,right=10pt,top=10pt,bottom=8pt,before skip=10pt,after skip=12pt,halign=center,
  lower separated=false,halign lower=center,
  fontlower=\popsemi\fontsize{9}{11.5}\selectfont\color{popmuted},
  #1}
\newcommand{\legende}[1]{\par\vspace{6pt}{\popsemi\fontsize{9}{11.5}\selectfont\color{popmuted}#1}}

% ============================================================
% Signature OnBuch+
% ============================================================
\newcommand{\poplogoinline}[1]{{\poplogo\fontsize{#1}{#1}\selectfont\color{popink}OnBuch\color{poporange}+}}
\newcommand{\signatureonbuch}{%
  \begin{tcolorbox}[enhanced,colback=popink,colframe=popink,boxrule=2.2pt,arc=10pt,
      drop shadow=poporange,left=14pt,right=14pt,top=10pt,bottom=10pt,before skip=0pt,after skip=0pt]
    \tikz[baseline=-0.7ex]\node[circle,fill=popgreen,draw=popcream,line width=1.3pt,minimum size=22pt,inner sep=0]{\popcheck{white}};%
    \hspace{9pt}\begin{minipage}[c]{0.62\linewidth}
      {\popxbold\fontsize{12}{13}\selectfont\color{popcream}Signed\ \textbullet\ {\poplogo OnBuch\color{poporange}+}}\par\vspace{2pt}
      {\raggedright\popsemi\fontsize{8}{10.5}\selectfont\color{popline}\MakeUppercase{OnBuch+ teaching team}\\\MakeUppercase{Follows the official MINESEC syllabus}\par}
    \end{minipage}\hfill
    {\poplogo\fontsize{22}{22}\selectfont\color{popcream}OnBuch\color{poporange}+}
  \end{tcolorbox}%
}

% ============================================================
% Page de garde
% #1 titre · #2 matière · #3 niveau · #4 module · #5 n° leçon · #6 durée ·
% #7 description / objectifs (texte ou itemize)
% ============================================================
\newcommand{\pagegarde}[7]{%
  \thispagestyle{empty}
  \newgeometry{a4paper,margin=1.7cm,top=1.6cm,bottom=1.4cm}%
  \begin{tikzpicture}[remember picture,overlay]
    \fill[poporange!18] ([xshift=-3.2cm,yshift=-3.6cm]current page.north east) circle (4.6cm);
    \fill[poppurple!14] ([xshift=2.4cm,yshift=7.6cm]current page.south west) circle (3.8cm);
  \end{tikzpicture}%
  {\poplogo\fontsize{30}{30}\selectfont\color{popink}OnBuch\color{poporange}+}\hfill
  \raisebox{0.6ex}{\poppill{Signed course \textbullet\ Premium}{popink}{popcream}}\par
  \vspace{1pt}\popsquiggle{poppurple}\par
  \vspace{8pt}
  \begin{tcolorbox}[enhanced,colback=poporange,colframe=popink,boxrule=2.8pt,arc=13pt,
      drop shadow=popink,left=18pt,right=18pt,top=12pt,bottom=13pt,after skip=10pt]
    \poppill{#2 \textbullet\ #3}{popink}{popcream}\hspace{5pt}\poppill{#4}{white}{popink}\par\vspace{8pt}
    {\popdisplay\fontsize{14}{14}\selectfont\color{popink}CHAPTER #5}\par\vspace{4pt}
    {\raggedright\hyphenpenalty=10000\exhyphenpenalty=10000\popdisplay\fontsize{25}{27}\selectfont\color{popcream}\MakeUppercase{#1}\par}
    \vspace{8pt}
    {\popxbold\fontsize{9.5}{9.5}\selectfont\color{popink}\MakeUppercase{Suggested time: #6}}
  \end{tcolorbox}
  \begin{tcolorbox}[enhanced,colback=white,colframe=popink,boxrule=2.2pt,arc=10pt,
      drop shadow=popink,left=16pt,right=16pt,top=10pt,bottom=10pt,after skip=8pt]
    \poppill{In this chapter}{poppurple}{white}\par\vspace{5pt}
    \popsemi\fontsize{9.3}{12.6}\selectfont\color{popink}#7
  \end{tcolorbox}
  \noindent\begin{minipage}[t]{0.487\linewidth}
    \begin{tcolorbox}[enhanced,colback=popgreen,colframe=popink,boxrule=2.2pt,arc=10pt,
        drop shadow=popink,left=11pt,right=11pt,top=10pt,bottom=10pt,height=3.1cm,valign=center]
      \tcbox[colback=white,colframe=popink,boxrule=1.3pt,arc=5pt,boxsep=0pt,nobeforeafter,
        left=3pt,right=3pt,top=3pt,bottom=3pt,valign=center]{\qrcode[height=1.9cm]{https://play.google.com/store/apps/details?id=com.luvvix.onbuch&hl=en}}%
      \hspace{8pt}\begin{minipage}[c]{0.5\linewidth}
        {\popdisplay\fontsize{11}{12.5}\selectfont\color{white}\MakeUppercase{Download OnBuch}}\par\vspace{4pt}
        {\raggedright\popsemi\fontsize{8.4}{11}\selectfont\color{white}Free \textbullet\ Google Play\\AI tutor, past papers, results\par}
      \end{minipage}
    \end{tcolorbox}
  \end{minipage}\hfill
  \begin{minipage}[t]{0.487\linewidth}
    \begin{tcolorbox}[enhanced,colback=poppurple,colframe=popink,boxrule=2.2pt,arc=10pt,
        drop shadow=popink,left=11pt,right=11pt,top=10pt,bottom=10pt,height=3.1cm,valign=center]
      \tikz[baseline=-0.7ex]\node[circle,fill=white,draw=popink,line width=1.3pt,minimum size=1.6cm,inner sep=0]{\popdisplay\fontsize{24}{24}\selectfont\color{poppurple}+};%
      \hspace{8pt}\begin{minipage}[c]{0.6\linewidth}
        {\popdisplay\fontsize{11}{12.5}\selectfont\color{white}\MakeUppercase{Go OnBuch+}}\par\vspace{4pt}
        {\raggedright\popsemi\fontsize{8.4}{11}\selectfont\color{white}All signed courses, unlimited AI tutor, PDF sheets — OnBuch+ tab in the app\par}
      \end{minipage}
    \end{tcolorbox}
  \end{minipage}
  \vspace{8pt}
  \popcontenu
  \vfill
  \signatureonbuch
  \restoregeometry
  \newpage
}

% Rangée « Ce cours contient » (page de garde)
\newcommand{\poptile}[3]{% #1 couleur #2 gros texte #3 légende
  \begin{tcolorbox}[enhanced,colback=#1,colframe=popink,boxrule=2pt,arc=9pt,shadow={2.5pt}{-2.5pt}{0pt}{popink},
      left=6pt,right=6pt,top=9pt,bottom=9pt,halign=center,valign=center,height=1.75cm,width=\linewidth,nobeforeafter]
    {\popdisplay\fontsize{12.5}{13}\selectfont\color{white}\MakeUppercase{#2}}\par\vspace{3pt}
    {\centering\popsemi\fontsize{7.8}{9.5}\selectfont\color{white}#3\par}
  \end{tcolorbox}}
\newcommand{\popcontenu}{%
  \par\noindent{\popxboldls\fontsize{7.5}{7.5}\selectfont\color{popmuted}THIS SIGNED COURSE CONTAINS}\par\vspace{7pt}
  \noindent\begin{minipage}[t]{0.235\linewidth}\poptile{popblue}{Course}{complete, illustrated, step by step}\end{minipage}\hfill
  \begin{minipage}[t]{0.235\linewidth}\poptile{poporange}{Methods}{and worked examples}\end{minipage}\hfill
  \begin{minipage}[t]{0.235\linewidth}\poptile{poppink}{Exercises}{exam style + solutions}\end{minipage}\hfill
  \begin{minipage}[t]{0.235\linewidth}\poptile{popgreen}{Summary sheet}{the essentials to remember}\end{minipage}\par}

% Sommaire
\newtcolorbox{sommaire}{enhanced,colback=white,colframe=popink,boxrule=2.2pt,arc=10pt,
  drop shadow=popink,left=18pt,right=18pt,top=13pt,bottom=13pt,before skip=6pt,after skip=20pt,
  before upper={\poppill{Contents}{poppurple}{white}\par\vspace{9pt}\popsemi\fontsize{10.6}{14.5}\selectfont\color{popink}}}

% Signature de fin de cours — poussée tout en bas de la dernière page
\newcommand{\findecours}{%
  \par\vfill\nopagebreak
  \begin{center}\popsquiggle{poporange}\end{center}\vspace{4pt}
  \signatureonbuch
}
\AtBeginDocument{\def\llbracket{\mathopen{[\mkern-3.5mu[}}\def\rrbracket{\mathclose{]\mkern-3.5mu]}}} % intervalles d'entiers (glyphe ⟦ absent de la police)
% Glyphes absents de Fira Math : redéfinis à partir de glyphes disponibles
\AtBeginDocument{%
  \def\setminus{\mathbin{\backslash}}\let\smallsetminus\setminus
  \def\vdots{\mathord{\vbox{\baselineskip3.2pt\lineskiplimit0pt\kern4pt\hbox{.}\hbox{.}\hbox{.}}}}%
  \def\ddots{\mathinner{\mkern1mu\raise6.5pt\hbox{.}\mkern2mu\raise3.6pt\hbox{.}\mkern2mu\raise0.7pt\hbox{.}\mkern1mu}}%
  \def\triangle{\mathord{\scalebox{0.9}{$\symup{\Delta}$}}}%
  \def\nabla{\mathord{\raisebox{\depth}{\scalebox{1}[-1]{$\symup{\Delta}$}}}}%
  \def\checkmark{\popcheck{popdark}}%
  \def\star{\mathbin{\ast}}\def\bigstar{\mathord{\ast}}%
  \def\diamond{\mathbin{\scalebox{0.6}{$\Diamond$}}}%
  \def\bigcup{\mathop{\mathchoice{\raisebox{-0.3em}{\scalebox{1.7}{$\cup$}}}{\scalebox{1.25}{$\cup$}}{\cup}{\cup}}\displaylimits}%
  \def\bigcap{\mathop{\mathchoice{\raisebox{-0.3em}{\scalebox{1.7}{$\cap$}}}{\scalebox{1.25}{$\cap$}}{\cap}{\cap}}\displaylimits}%
  \def\lesseqqgtr{\mathrel{\lessgtr}}\def\bigoplus{\mathop{\oplus}\displaylimits}%
}
\providecommand{\ssi}{if and only if} % abréviation fréquente
\AtBeginDocument{\providecommand{\wideparen}{\overparen}} % arc de cercle

\begin{document}

\pagegarde{Reproduction}{Biology}{Upper Sixth}{Upper Sixth — Biology}{1}{21 periods}{This chapter explores how living organisms perpetuate their species: asexual and sexual reproduction in plants and animals, the human reproductive systems and their hormonal control, fertilization, embryonic development, birth and birth control. It ends with the study of life cycles (alternation of generations, metagenesis) and the patterns and hormonal regulation of growth.
\vspace{4pt}
\begin{multicols}{2}\small
\begin{itemize}
\item By the end of this chapter I can define reproduction and distinguish asexual from sexual reproduction with Cameroonian examples.
\item By the end of this chapter I can describe the mechanisms of asexual reproduction in plants and animals.
\item By the end of this chapter I can explain pollen and ovule formation, pollination and double fertilization in flowering plants.
\item By the end of this chapter I can describe seed formation, dispersal and the conditions and stages of germination.
\item By the end of this chapter I can label and state the functions of the parts of the male and female human reproductive systems.
\item By the end of this chapter I can explain the hormonal control of the testes and ovaries, spermatogenesis, oogenesis and the menstrual cycle.
\end{itemize}\end{multicols}}

\begin{sommaire}
\begin{multicols}{2}
\begin{enumerate}
\item Asexual Reproduction: Concepts and Mechanisms
\item Sexual Reproduction in Flowering Plants: Gamete Formation to Fertilization
\item Seeds, Dispersal and Germination
\item The Male Reproductive System and Spermatogenesis
\item The Female Reproductive System, Oogenesis and the Menstrual Cycle
\item Fertilization, Embryonic Development and Birth
\item Birth Control Methods
\item Life Cycles: Alternation of Generations and Metagenesis
\item Growth: Curves, Patterns and Hormonal Control
\item Integration activity
\item Methods and worked examples
\item Exercises
\item Solutions to the exercises
\item Summary sheet
\end{enumerate}
\end{multicols}
\end{sommaire}

\begin{objectifs}
\begin{itemize}
\item By the end of this chapter I can define reproduction and distinguish asexual from sexual reproduction with Cameroonian examples.
\item By the end of this chapter I can describe the mechanisms of asexual reproduction in plants and animals.
\item By the end of this chapter I can explain pollen and ovule formation, pollination and double fertilization in flowering plants.
\item By the end of this chapter I can describe seed formation, dispersal and the conditions and stages of germination.
\item By the end of this chapter I can label and state the functions of the parts of the male and female human reproductive systems.
\item By the end of this chapter I can explain the hormonal control of the testes and ovaries, spermatogenesis, oogenesis and the menstrual cycle.
\item By the end of this chapter I can describe fertilization, embryonic development, the birth process and methods of birth control.
\item By the end of this chapter I can compare alternation of generations and metagenesis with labelled life-cycle diagrams.
\item By the end of this chapter I can interpret growth curves, distinguish growth patterns and explain the hormonal control of growth.
\end{itemize}
\end{objectifs}

\begin{prerequis}
\begin{itemize}
\item Cell structure and the differences between mitosis and meiosis (Lower Sixth).
\item Structure of a flower and its floral parts (Lower Sixth / Form 5).
\item Basic notions of hormones and endocrine glands (Form 5 / Lower Sixth).
\item Chromosomes, genes and the haploid/diploid distinction (Lower Sixth).
\item General plant and animal tissue organisation (Lower Sixth).
\end{itemize}
\end{prerequis}

\newpage

\coursec{Asexual Reproduction: Concepts and Mechanisms}

In a village near Bafoussam, Ma Ngo plants cassava stem cuttings that grow into whole new plants without any seed, while her neighbour's maize field must be sown with seeds every season. At the Bamenda regional hospital, a midwife explains to a young couple why a pregnancy test turned positive only two weeks after fertilization. Meanwhile, a Form Five student wonders why her younger brother suddenly grew taller and developed a deeper voice at age fourteen. All these everyday Cameroonian observations are governed by one great biological theme: reproduction and growth. This chapter unravels the cellular, hormonal and developmental mechanisms behind them --- knowledge that is heavily examined at the GCE Advanced Level.

We begin with the simplest strategy of all: reproduction without sex, without a partner, and without gametes.

\courssub{Reproduction and Its Two Great Strategies}

\textbf{From observation to concept.}\par
Drop a single \emph{Amoeba} into a drop of pond water and return a few days later: you may find several. No mate was involved, no egg, no sperm. Yet the population increased. Reproduction is clearly something more general than the familiar union of male and female.

\begin{definition}[Reproduction]
\cle{Reproduction} is the biological process by which living organisms produce new individuals of their own kind, thereby ensuring the continuity of the species from one generation to the next.
\end{definition}

Biologists recognise two fundamental strategies:

\begin{itemize}
\item \textbf{\cle{Asexual reproduction}}: a single parent produces offspring \emph{without the formation or fusion of gametes}. Because all the offspring inherit their genetic material (by mitosis) from one parent only, they are \cle{genetically identical} to that parent and to one another.
\item \textbf{\cle{Sexual reproduction}}: two parents (or two gamete-producing structures) contribute gametes which fuse (\emph{fertilization}) to form a zygote. The offspring receive a new combination of genes and are therefore \emph{genetically different} from each parent.
\end{itemize}

\begin{definition}[Clone]
A \cle{clone} is a group of genetically identical individuals derived from a single parent by asexual reproduction (or, artificially, by cloning techniques).
\end{definition}

\begin{exemplebox}[Clones around us]
All the plantain suckers taken from one mother plant in a plantation near Kumba are clones of that plant. Likewise, all the bacteria in a colony grown from one cell on an agar plate form a clone.
\end{exemplebox}

\courssub{Advantages and Disadvantages of Asexual Reproduction}

Why should an organism give up the genetic mixing of sexual reproduction? The answer lies in a balance of costs and benefits:

\begin{center}
\begin{tabularx}{\linewidth}{|l|X|X|}
\hline
\rowcolor{popblueL} & \textbf{Advantages} & \textbf{Disadvantages} \\
\hline
\textbf{Speed} & Very rapid population increase (a bacterium can divide every 20 minutes) & --- \\
\hline
\textbf{No mate needed} & A single isolated individual can reproduce; useful in sparse populations or sessile organisms & --- \\
\hline
\textbf{Energy} & No energy wasted producing gametes, courtship or finding partners & --- \\
\hline
\textbf{Fidelity} & Well-adapted genotypes are preserved exactly & \textbf{No genetic variation}: offspring cannot easily adapt if the environment changes \\
\hline
\textbf{---} & --- & \textbf{Accumulation of harmful mutations} and uniform susceptibility: one disease can wipe out an entire clone \\
\hline
\end{tabularx}
\end{center}

\begin{attention}[The variation trap]
A very common exam error is to write that asexual reproduction produces ``weak'' offspring. This is wrong: clones can be perfectly vigorous. The correct point is that they show \emph{no genetic variation} (apart from rare mutations), so the whole population is equally vulnerable to a new parasite, drought or pesticide.
\end{attention}

\courssub{Mechanisms of Asexual Reproduction in Animals and Micro-organisms}

\textbf{Binary fission.}\par
In \cle{binary fission} the parent cell grows, duplicates its genetic material, and then divides into two approximately equal daughter cells, each of which becomes an independent organism. This is the mode of reproduction of \emph{Amoeba}, \emph{Paramecium}, and of bacteria.

In \emph{Amoeba}, the nucleus first divides by mitosis, then the cytoplasm constricts in the middle until two daughter amoebae separate. In bacteria, the circular chromosome is replicated, the two copies move apart, and a new transverse cell wall grows across the middle of the cell.

\textbf{Budding.}\par
In \cle{budding} a small outgrowth (\emph{bud}) forms on the parent body, grows, and eventually detaches as a new individual. The offspring is initially much smaller than the parent --- this distinguishes budding from binary fission, where the two products are roughly equal. Budding occurs in \emph{Hydra} (a freshwater cnidarian) and in yeast (\emph{Saccharomyces}), the fungus used in baking and in brewing palm wine and ``sha'a''.

\textbf{Fragmentation.}\par
In \cle{fragmentation} the body of the parent breaks into two or more pieces, and \emph{each piece regenerates} the missing parts to become a complete organism. The filamentous alga \emph{Spirogyra} reproduces this way when its filaments break; each fragment grows by cell division into a new filament. Among animals, flatworms such as planarians show remarkable fragmentation: a worm cut into pieces can regenerate a complete individual from each piece.

\textbf{Sporulation (spore formation).}\par
Many organisms produce \cle{spores}: small, usually unicellular reproductive bodies, often protected by a resistant wall, each capable of developing directly into a new individual without fertilization. Spore formation occurs in fungi (e.g. \emph{Mucor} and \emph{Rhizopus}, the black bread mould, which produces sporangia full of spores), in mosses and in ferns (sporangia on the underside of fern fronds). Spores are light and are dispersed by wind, which spreads the organism widely.

\begin{popfigure}
\begin{center}
\begin{tikzpicture}[scale=0.9, every node/.style={font=\small}]
% Binary fission
\node[font=\small\bfseries, popblue] at (2.2,3.6) {Binary fission (\emph{Amoeba})};
\draw[fill=popblueL, draw=popblue, thick] (0.6,2.4) ellipse (0.55 and 0.45);
\fill[poporange] (0.6,2.4) circle (0.12);
\draw[->, thick, popdark] (1.35,2.4) -- (1.9,2.4);
\draw[fill=popblueL, draw=popblue, thick] (2.9,2.4) ellipse (0.75 and 0.5);
\fill[poporange] (2.6,2.4) circle (0.1);
\fill[poporange] (3.2,2.4) circle (0.1);
\draw[->, thick, popdark] (3.85,2.4) -- (4.4,2.4);
\draw[fill=popblueL, draw=popblue, thick] (5.1,2.4) ellipse (0.45 and 0.4);
\fill[poporange] (5.1,2.4) circle (0.1);
\draw[fill=popblueL, draw=popblue, thick] (6.3,2.4) ellipse (0.45 and 0.4);
\fill[poporange] (6.3,2.4) circle (0.1);
% Budding
\node[font=\small\bfseries, popgreen] at (2.2,0.9) {Budding (\emph{Hydra})};
\draw[fill=popgreenL, draw=popgreen, thick] (0.6,-0.6) -- (0.6,0.3) arc (180:0:0.25) -- (1.1,-0.6) -- cycle;
\draw[popgreen, thick] (0.72,0.3) -- (0.6,0.65); \draw[popgreen, thick] (0.85,0.3) -- (0.85,0.7); \draw[popgreen, thick] (0.98,0.3) -- (1.1,0.65);
\draw[->, thick, popdark] (1.35,0) -- (1.9,0);
\draw[fill=popgreenL, draw=popgreen, thick] (2.7,-0.6) -- (2.7,0.3) arc (180:0:0.25) -- (3.2,-0.6) -- cycle;
\draw[popgreen, thick] (2.82,0.3) -- (2.7,0.65); \draw[popgreen, thick] (2.95,0.3) -- (2.95,0.7); \draw[popgreen, thick] (3.08,0.3) -- (3.2,0.65);
\draw[fill=popgreenL, draw=popgreen, thick] (3.2,-0.25) -- (3.2,0.15) arc (180:0:0.14) -- (3.48,-0.25) -- cycle;
\draw[->, thick, popdark] (3.75,0) -- (4.3,0);
\draw[fill=popgreenL, draw=popgreen, thick] (4.9,-0.6) -- (4.9,0.3) arc (180:0:0.25) -- (5.4,-0.6) -- cycle;
\draw[popgreen, thick] (5.02,0.3) -- (4.9,0.65); \draw[popgreen, thick] (5.15,0.3) -- (5.15,0.7); \draw[popgreen, thick] (5.28,0.3) -- (5.4,0.65);
\draw[fill=popgreenL, draw=popgreen, thick] (6.1,-0.6) -- (6.1,-0.25) arc (180:0:0.14) -- (6.38,-0.6) -- cycle;
\draw[popgreen, thick] (6.18,-0.25) -- (6.1,-0.05); \draw[popgreen, thick] (6.3,-0.25) -- (6.38,-0.05);
% Fragmentation
\node[font=\small\bfseries, poppurple] at (10.2,3.6) {Fragmentation (\emph{Spirogyra})};
\draw[fill=popblueL, draw=poppurple, thick, rounded corners=2pt] (8.2,2.2) rectangle (11.2,2.6);
\foreach \x in {8.95,9.7,10.45} {\draw[poppurple] (\x,2.2) -- (\x,2.6);}
\draw[->, thick, popdark] (9.7,2.05) -- (9.7,1.55);
\draw[fill=popblueL, draw=poppurple, thick, rounded corners=2pt] (8.2,0.9) rectangle (9.55,1.3);
\draw[poppurple] (8.9,0.9) -- (8.9,1.3);
\draw[fill=popblueL, draw=poppurple, thick, rounded corners=2pt] (9.9,0.9) rectangle (11.2,1.3);
\draw[poppurple] (10.55,0.9) -- (10.55,1.3);
\node[popmuted] at (9.7,0.55) {each fragment grows into a new filament};
\end{tikzpicture}
\end{center}
\legende{Three mechanisms of asexual reproduction. \textbf{Top:} binary fission in \emph{Amoeba} --- nucleus divides, then cytoplasm splits into two equal daughters. \textbf{Bottom left:} budding in \emph{Hydra} --- a small bud grows on the parent and detaches. \textbf{Right:} fragmentation in \emph{Spirogyra} --- the filament breaks and each fragment grows into a complete filament.}
\end{popfigure}

\begin{aretenir}[Key distinctions]
\begin{itemize}
\item Binary fission: parent splits into \textbf{two equal} daughters (\emph{Amoeba}, bacteria).
\item Budding: \textbf{unequal} division --- a small bud on a larger parent (\emph{Hydra}, yeast).
\item Fragmentation: body breaks into pieces, each \textbf{regenerates} a whole organism (\emph{Spirogyra}, planaria).
\item Sporulation: resistant \textbf{spores} develop directly into new individuals (fungi, mosses, ferns).
\end{itemize}
\end{aretenir}

\courssub{Vegetative Propagation in Plants}

\textbf{Natural vegetative propagation.}\par
Flowering plants possess specialised vegetative organs --- modified stems, roots or leaves --- that can detach or spread and give rise to new plants. This is \cle{vegetative propagation}: asexual reproduction using vegetative (non-reproductive) organs.

\begin{popfigure}
\begin{center}
\begin{tikzpicture}[grow=down, level distance=1.5cm,
 every node/.style={draw=popblue, fill=popblueL, rounded corners=3pt, font=\small, align=center, thick},
 level 1/.style={sibling distance=2.6cm},
 edge from parent/.style={draw=popdark, thick}]
\node[fill=poporange!30, draw=poporange, font=\small\bfseries] {Natural vegetative propagules}
 child { node[align=center]{Runner\\ \emph{(grass, strawberry)}} }
 child { node[align=center]{Rhizome\\ \emph{(ginger, canna)}} }
 child { node[align=center]{Tuber\\ \emph{(yam, Irish potato)}} }
 child { node[align=center]{Bulb\\ \emph{(onion)}} }
 child { node[align=center]{Sucker\\ \emph{(plantain, banana)}} };
\end{tikzpicture}
\end{center}
\legende{Classification of natural vegetative propagules with a familiar Cameroonian example for each. Runners are horizontal aerial stems; rhizomes are horizontal underground stems; tubers are swollen storage stems or roots; bulbs are compressed stems with fleshy leaves; suckers are shoots arising from the base of the mother plant.}
\end{popfigure}

\begin{itemize}
\item \textbf{Runners}: horizontal stems growing along the soil surface; at each node, roots and a new shoot develop (grasses, strawberry).
\item \textbf{Rhizomes}: horizontal \emph{underground} stems with nodes and scale leaves; buds on the rhizome produce new shoots (ginger, canna lily).
\item \textbf{Tubers}: swollen underground storage organs. Stem tubers such as yam and Irish potato bear ``eyes'' (buds) that sprout into new plants.
\item \textbf{Bulbs}: a very short stem surrounded by fleshy storage leaves (onion); lateral buds develop into new bulbs.
\item \textbf{Suckers}: shoots arising from buds at the base of the mother plant or on its roots (plantain, banana, pineapple).
\end{itemize}

\textbf{Artificial vegetative propagation.}\par
Farmers and horticulturists exploit the same principle deliberately:

\begin{itemize}
\item \textbf{Cuttings}: a piece of stem (cassava, hibiscus, sugar cane) or root is planted and develops adventitious roots and shoots.
\item \textbf{Grafting}: a shoot or bud (the \emph{scion}) of a desirable variety is joined to the rooted stem (the \emph{stock}) of another plant. Widely used for mango, citrus and improved cocoa varieties in the South West and Centre Regions.
\item \textbf{Layering}: a low branch is bent to the ground and covered with soil; it roots while still attached to the parent, then is separated.
\item \textbf{Tissue culture (micropropagation)}: tiny pieces of plant tissue are grown on sterile nutrient medium to produce thousands of identical plantlets --- used for disease-free plantain and banana planting material.
\end{itemize}

\begin{methode}[Demonstrating vegetative propagation with a cassava cutting]
\begin{enumerate}
\item Select a healthy mature cassava stem and cut portions about \SI{20}{cm} to \SI{25}{cm} long, each bearing at least 3--4 nodes (buds). Cut the lower end obliquely.
\item Plant the cuttings upright (or slightly slanting) in moist, well-drained soil or in a trough of moist sawdust, burying about two-thirds of the length. Keep some cuttings upside-down as a control.
\item Water regularly and observe weekly for 4--6 weeks. Record the appearance of shoots (from nodes above ground) and adventitious roots (from buried nodes).
\item Conclusion: a complete new plant has formed from a stem fragment, without seed --- evidence of vegetative propagation. Note that polarity matters: inverted cuttings root poorly.
\end{enumerate}
\end{methode}

\begin{attention}[A tuber is not a seed!]
A classic GCE trap: candidates write that yam or Irish potato ``seeds'' are planted. Wrong. A \textbf{seed} is a fertilized ovule containing an embryo --- the product of \emph{sexual} reproduction. A tuber is a swollen stem or root used for \emph{vegetative} propagation. Saying ``yam seeds'' or ``cassava reproduces by seeds'' in an essay costs marks. Similarly, do not confuse \emph{rhizome} (underground stem, with nodes) with \emph{root tuber}.
\end{attention}

\courssub{Parthenogenesis: A Special Case}

In a few animals, an \textbf{unfertilized egg} can develop into a new individual. This is \cle{parthenogenesis} (from Greek \emph{parthenos}, virgin). Although an egg cell is involved, there is no fertilization, so the process is functionally asexual.

\begin{itemize}
\item In \textbf{aphids}, females produce generation after generation of daughters from unfertilized eggs during the favourable season --- a spectacularly rapid multiplication on crop plants.
\item In \textbf{honey bees}, unfertilized eggs develop into males (drones), while fertilized eggs develop into females (workers and queens). This is exploited by beekeepers in the North West Region.
\end{itemize}

\begin{savaistu}[Vegetative propagation and the Cameroonian economy]
Cameroon's food security rests heavily on vegetative propagation. Cassava, the staple of the forest zone, is multiplied almost exclusively by stem cuttings; plantain and banana, the backbone of the Littoral and South West food economy, are propagated by suckers. Improved cocoa and robusta coffee planting material distributed to farmers is produced by grafting and rooted cuttings, so that the desirable traits of high-yielding, disease-tolerant clones are preserved exactly --- something seeds could not guarantee.
\end{savaistu}

\begin{exoresolu}[Worked exercise --- GCE style]
\textbf{Statement.} (a) Define asexual reproduction. (b) A farmer near Bafoussam multiplies her plantain plantation using suckers, while her neighbour propagates improved cocoa by grafting. State the type of reproduction involved in each case, and give \emph{two} advantages and \emph{one} disadvantage of multiplying crops this way.
\tcblower
\textbf{Solution.}\\
(a) \emph{Asexual reproduction} is the production of offspring from a single parent without the formation or fusion of gametes, so that the offspring are genetically identical to the parent (clones).\\[4pt]
(b) Both suckers and grafting are \textbf{vegetative propagation}, a form of asexual reproduction.\\
\emph{Advantages} (any two):
\begin{itemize}
\item Rapid multiplication --- many plantlets obtained in one season;
\item No need for a mate or pollinator, and no waiting for seed production;
\item The desirable qualities of the parent (yield, taste, disease tolerance) are preserved exactly in the clones.
\end{itemize}
\emph{Disadvantage} (any one):
\begin{itemize}
\item No genetic variation: all plants are equally susceptible to a new disease or pest;
\item Harmful mutations or viral infections of the parent are transmitted to all offspring.
\end{itemize}
\end{exoresolu}

\begin{exercice}[Check your understanding]
\begin{enumerate}
\item Distinguish clearly between binary fission and budding, naming one organism for each.
\item Explain why a field of cassava planted from cuttings of a single parent plant could be devastated by one viral disease, whereas a maize field grown from seed might survive.
\item Classify the following as natural or artificial vegetative propagation, and name the organ involved: (i) ginger spreading in a garden; (ii) a farmer planting cassava stems; (iii) onion bulbs multiplying; (iv) a cocoa scion joined to a stock.
\item A planarian is cut into three pieces and three complete worms are obtained. Name and define the process.
\end{enumerate}
\end{exercice}

\begin{corrige}[Exercise --- Check your understanding]
\begin{enumerate}
\item Binary fission: the parent divides into two \emph{equal} daughter cells (e.g. \emph{Amoeba}, bacteria). Budding: a small outgrowth forms on the parent and detaches as a smaller new individual (e.g. \emph{Hydra}, yeast).
\item The cassava plants are clones: they share identical genes, so none possesses resistance to the virus and all are infected. The maize plants, produced sexually, are genetically varied, so some individuals may carry resistance and survive.
\item (i) Natural --- rhizome (underground stem); (ii) artificial --- stem cutting; (iii) natural --- bulb; (iv) artificial --- grafting.
\item Fragmentation: asexual reproduction in which the parent's body breaks into pieces, each of which regenerates the missing parts to form a complete new organism.
\end{enumerate}
\end{corrige}

\begin{aretenir}[Exam tip --- GCE Advanced Level]
Whenever you are asked to \emph{describe} a type of asexual reproduction, always: (1) define it precisely, (2) describe the mechanism step by step, and (3) give a \textbf{named example} (\emph{Amoeba}, \emph{Hydra}, \emph{Spirogyra}, \emph{Mucor}, cassava, plantain...). An answer without a named example rarely earns full marks. And remember: tubers, rhizomes and cuttings are \emph{not} seeds.
\end{aretenir}

\coursec{Sexual Reproduction in Flowering Plants: Gamete Formation to Fertilization}

In the previous section, we saw that asexual reproduction is fast and economical, but it produces offspring that are genetically identical to the parent. In a changing environment --- a new parasite, a drought, a shift in climate --- such uniformity can be fatal. Nature's answer is \cle{sexual reproduction}: the mixing of genetic material from two parents so that every offspring is a unique combination. In flowering plants (angiosperms), this process is centred on a remarkable organ: the \cle{flower}. In this section we follow the whole journey, from the formation of pollen and ovules, through pollination, to the unique event of double fertilization.

\courssub{The Concept of Sexual Reproduction}

\begin{definition}[Sexual reproduction]
\cle{Sexual reproduction} is the type of reproduction involving the fusion of two haploid gametes --- a male gamete and a female gamete --- during \cle{fertilization}, to form a diploid \cle{zygote} which develops into a new individual.
\end{definition}

Because each gamete carries a different shuffled set of genes (thanks to meiosis and random fertilization), the zygote --- and therefore the offspring --- is genetically different from both parents and from its siblings. This \cle{genetic variation} is the raw material on which natural selection acts; it is the great advantage of sexual reproduction over asexual reproduction.

\begin{propriete}[Advantage of sexual reproduction]
Sexual reproduction produces \textbf{genetic variation} among offspring, increasing the chance that some individuals survive environmental changes, diseases or new predators. Its costs are that it is slower, requires two parents (or both sex organs), and uses more energy than asexual reproduction.
\end{propriete}

\courssub{Structure of a Typical Flower}

A flower is essentially a shoot modified for reproduction. Its parts are arranged in four \cle{whorls} on a swollen tip of the stalk called the \cle{receptacle}:

\begin{itemize}
\item \textbf{Sepals} (collectively the \emph{calyx}): usually green, they protect the flower in bud.
\item \textbf{Petals} (collectively the \emph{corolla}): often large, coloured and scented to attract pollinators.
\item \textbf{Stamens} (the \emph{androecium}, male part): each consists of an \cle{anther}, which produces pollen, borne on a stalk called the \cle{filament}.
\item \textbf{Carpels} (the \emph{gynoecium}, female part): each consists of a sticky \cle{stigma} (receives pollen), a \cle{style} (stalk supporting the stigma) and an \cle{ovary} containing one or more \cle{ovules}.
\end{itemize}

\begin{popfigure}
\begin{center}
\begin{tikzpicture}[scale=1.05, font=\small]
% receptacle and stalk
\draw[thick, popgreen] (0,-1.6) -- (0,-0.4);
\fill[popgreenL] (-0.9,-0.4) .. controls (-0.5,0.1) and (0.5,0.1) .. (0.9,-0.4) -- cycle;
\node[popdark] at (0,-1.85) {stalk};
% sepals
\draw[thick, popgreen] (-0.75,-0.35) -- (-1.5,0.35);
\draw[thick, popgreen] (0.75,-0.35) -- (1.5,0.35);
\node[popgreen, align=center] at (-1.75,0.5) {sepal};
% petals
\draw[very thick, poppink] (-0.6,0.0) .. controls (-1.9,1.3) and (-1.3,2.3) .. (-0.5,2.6);
\draw[very thick, poppink] (0.6,0.0) .. controls (1.9,1.3) and (1.3,2.3) .. (0.5,2.6);
\node[poppink] at (2.0,2.4) {petal};
% stamens
\draw[thick, popgold] (-0.45,0.05) -- (-0.85,1.9);
\draw[thick, popgold] (0.45,0.05) -- (0.85,1.9);
\fill[poporange] (-0.85,1.9) ellipse (0.18 and 0.32);
\fill[poporange] (0.85,1.9) ellipse (0.18 and 0.32);
\node[poporange, align=center] at (-2.3,2.1) {anther};
\draw[->, popmuted] (-2.05,2.0) -- (-1.05,1.95);
\node[popgold, align=center] at (2.3,1.35) {filament};
\draw[->, popmuted] (2.05,1.3) -- (0.95,1.5);
% carpel
\fill[popblueL] (0,0.15) ellipse (0.55 and 0.6);
\draw[thick, popblue] (0,0.15) ellipse (0.55 and 0.6);
\draw[thick, popblue] (0,0.75) -- (0,2.35);
\fill[poppurple] (0,2.5) ellipse (0.28 and 0.18);
% ovule
\draw[thick, poppurple] (0,0.1) ellipse (0.2 and 0.28);
\fill[poppurple] (0,0.1) circle (0.05);
% labels carpel parts
\node[poppurple] at (-1.9,2.75) {stigma};
\draw[->, popmuted] (-1.6,2.65) -- (-0.3,2.5);
\node[popblue] at (1.7,2.0) {style};
\draw[->, popmuted] (1.45,1.9) -- (0.08,1.6);
\node[popblue] at (1.8,0.15) {ovary};
\draw[->, popmuted] (1.55,0.15) -- (0.6,0.15);
\node[poppurple] at (-1.8,-0.15) {ovule};
\draw[->, popmuted] (-1.5,-0.1) -- (-0.25,0.05);
\node[popgreen] at (1.9,-0.55) {receptacle};
\draw[->, popmuted] (1.6,-0.5) -- (0.7,-0.35);
\end{tikzpicture}
\end{center}
\legende{Longitudinal section through a typical insect-pollinated flower, showing the four whorls: sepals, petals, stamens (anther + filament) and carpel (stigma, style, ovary with ovule).}
\end{popfigure}

\courssub{Pollen Formation: Microsporogenesis and Microgametogenesis}

Pollen grains are the male gametophytes of flowering plants. They are produced inside the \cle{anther}, which typically has four pollen sacs (microsporangia). The process occurs in two phases: \cle{microsporogenesis} (meiosis) and \cle{microgametogenesis} (mitotic divisions).

\begin{methode}[Microsporogenesis and pollen maturation]
\begin{enumerate}
\item Inside each pollen sac, diploid \cle{pollen mother cells} (microsporocytes, $2n$) are formed from the sporogenous tissue.
\item Each pollen mother cell undergoes \textbf{meiosis}, producing a tetrad of \textbf{four haploid microspores} ($n$).
\item Each microspore matures into a \cle{pollen grain} (the male gametophyte): its nucleus divides by \textbf{mitosis} to give a large \cle{tube nucleus} (vegetative nucleus) and a smaller \cle{generative nucleus}. The wall thickens into a resistant outer \emph{exine} (sporopollenin) and an inner \emph{intine} (cellulose/pectin).
\item In many species, the generative nucleus divides by a second \textbf{mitosis} \emph{after} pollination, inside the growing pollen tube, to form \textbf{two male gamete nuclei} (sperm cells). In some species this division occurs before pollen release (bicellular vs tricellular pollen).
\end{enumerate}
\end{methode}

\courssub{Ovule Formation: Megasporogenesis and Megagametogenesis}

Inside the ovary, each \cle{ovule} is attached to the placenta by a stalk (the \cle{funicle}) and is protected by one or two \cle{integuments}, leaving a small pore, the \cle{micropyle}. The female gametophyte (embryo sac) develops as follows:

\begin{methode}[Megasporogenesis and embryo sac development (Polygonum type)]
\begin{enumerate}
\item A diploid \cle{megaspore mother cell} (megasporocyte, $2n$) forms in the ovule tissue (nucellus).
\item It undergoes \textbf{meiosis} to give four haploid megaspores ($n$); \textbf{three degenerate} and only \textbf{one functional megaspore} survives (usually the chalazal one).
\item The nucleus of the functional megaspore divides by \textbf{three successive mitotic divisions} (without cytokinesis at first) to give \textbf{eight haploid nuclei}.
\item These organise into the \cle{embryo sac}, which is \textbf{7-celled and 8-nucleate}:
\begin{itemize}
\item At the micropylar end: one \cle{egg cell} flanked by two \cle{synergids} (together forming the \emph{egg apparatus}).
\item At the chalazal end: three \cle{antipodal cells}.
\item In the centre: a large \cle{central cell} containing \textbf{two polar nuclei} (which often fuse before fertilization to form a diploid secondary nucleus, $2n$).
\end{itemize}
\end{enumerate}
\end{methode}

\begin{popfigure}
\begin{center}
\begin{tikzpicture}[scale=0.92, font=\small]
% MICROSPOROGENESIS (left)
\node[popdark, font=\bfseries] at (-4.6,4.4) {Microsporogenesis (anther)};
\draw[thick, poporange] (-4.6,3.4) circle (0.5);
\node at (-4.6,3.4) {$2n$};
\node[popmuted, align=center] at (-6.3,3.4) {pollen mother\\cell};
\draw[->, thick, popblue] (-4.6,2.75) -- (-4.6,2.15) node[midway, right, popdark]{meiosis};
\foreach \x/\y in {-5.15/1.45,-4.05/1.45,-5.15/0.75,-4.05/0.75}{\draw[thick, poporange] (\x,\y) circle (0.28); \node at (\x,\y) {$n$};}
\node[popmuted, align=center] at (-6.3,1.1) {4 microspores};
\draw[->, thick, popblue] (-4.6,0.3) -- (-4.6,-0.3) node[midway, right, popdark]{mitosis};
\draw[thick, poporange] (-4.6,-1.1) circle (0.55);
\fill[poppurple] (-4.85,-1.0) circle (0.09);
\fill[popgreen] (-4.35,-1.25) circle (0.13);
\node[popmuted, align=center] at (-6.4,-1.1) {pollen grain};
\node[poppurple, align=center] at (-3.1,-0.75) {tube\\nucleus};
\node[popgreen, align=center] at (-3.1,-1.6) {generative\\nucleus};
% MEGASPOROGENESIS (right)
\node[popdark, font=\bfseries] at (4.6,4.4) {Megasporogenesis (ovule)};
\draw[thick, poppink] (4.6,3.4) circle (0.5);
\node at (4.6,3.4) {$2n$};
\node[popmuted, align=center] at (6.9,3.4) {megaspore\\mother cell};
\draw[->, thick, popblue] (4.6,2.75) -- (4.6,2.15) node[midway, right, popdark]{meiosis};
\draw[thick, poppink] (4.6,1.55) circle (0.3);
\node at (4.6,1.55) {$n$};
\draw[thick, popmuted] (5.35,1.55) circle (0.22);
\draw[thick, popmuted] (5.95,1.55) circle (0.22);
\draw[thick, popmuted] (6.55,1.55) circle (0.22);
\draw[thick, popmuted] (5.2,1.9) -- (6.7,1.2);
\node[popmuted, align=center] at (6.9,2.2) {3 degenerate};
\node[popmuted, align=center] at (2.6,1.55) {functional\\megaspore};
\draw[->, thick, popblue] (4.6,1.1) -- (4.6,0.5) node[midway, right, popdark]{3 mitoses};
% embryo sac
\draw[very thick, poppink] (4.6,-1.3) ellipse (1.15 and 1.7);
\fill[poppurple] (4.05,-0.35) circle (0.14);
\fill[popgold] (4.6,-0.55) circle (0.11);
\fill[popgold] (5.15,-0.35) circle (0.11);
\fill[popblue] (4.35,-1.3) circle (0.11);
\fill[popblue] (4.85,-1.3) circle (0.11);
\fill[popmuted] (4.05,-2.25) circle (0.1);
\fill[popmuted] (4.6,-2.4) circle (0.1);
\fill[popmuted] (5.15,-2.25) circle (0.1);
\node[poppurple] at (2.35,-0.35) {egg cell};
\draw[->, popmuted] (3.15,-0.35) -- (3.85,-0.35);
\node[popgold] at (6.9,-0.4) {synergids};
\draw[->, popmuted] (6.3,-0.4) -- (5.3,-0.4);
\node[popblue] at (6.9,-1.3) {polar nuclei};
\draw[->, popmuted] (6.25,-1.3) -- (5.0,-1.3);
\node[popmuted] at (6.9,-2.3) {antipodals};
\draw[->, popmuted] (6.3,-2.3) -- (5.3,-2.3);
\node[popmuted] at (4.6,-3.35) {embryo sac ($n$): 7 cells, 8 nuclei};
\end{tikzpicture}
\end{center}
\legende{Microsporogenesis (left) and megasporogenesis (right). Meiosis produces haploid spores; mitosis then builds the pollen grain (tube + generative nuclei) and the 7-celled, 8-nucleate embryo sac.}
\end{popfigure}

\begin{attention}[Ploidy trap]
All eight nuclei of the embryo sac are \textbf{haploid ($n$)} because they come from mitotic divisions of one haploid megaspore. The pollen mother cell and megaspore mother cell are \textbf{diploid ($2n$)}. Do not confuse the two polar nuclei ($n + n$) with a diploid nucleus --- they only become part of a \textbf{triploid ($3n$)} nucleus after fusion with a male gamete at fertilization. Note: the two polar nuclei often fuse \emph{before} fertilization to form a diploid ($2n$) central cell nucleus; the endosperm is then formed by fusion of this $2n$ nucleus with a male $n$ nucleus, still giving $3n$.
\end{attention}

\courssub{Pollination}

\begin{definition}[Pollination]
\cle{Pollination} is the transfer of pollen grains from the \textbf{anther} to the \textbf{stigma} of a flower.
\end{definition}

Two situations are possible:

\begin{itemize}
\item \textbf{Self-pollination}: pollen lands on the stigma of the \emph{same} flower (autogamy), or of another flower on the \emph{same plant} (geitonogamy).
\item \textbf{Cross-pollination}: pollen is carried to the stigma of a flower on a \emph{different plant} of the same species (xenogamy).
\end{itemize}

\begin{propriete}[Self- versus cross-pollination]
\cle{Cross-pollination} mixes genes from two different parents and therefore \textbf{increases genetic variation}, giving offspring a better chance of adapting to change. \cle{Self-pollination} produces more uniform offspring and \textbf{preserves desirable characters} (useful to farmers who want a stable crop variety), and it can occur even when pollinators are scarce.
\end{propriete}

\textbf{Agents of pollination.} Pollen may be carried by \textbf{insects} (bees, butterflies), \textbf{wind}, \textbf{water} (in some aquatic plants e.g. \emph{Vallisneria}) or \textbf{animals} such as birds (sunbirds) and bats. Each agent corresponds to a set of floral adaptations.

\begin{methode}[Comparing insect- and wind-pollinated flowers]
To distinguish the two types at the GCE, examine the following features:
\begin{center}
\begin{tabularx}{\linewidth}{|l|X|X|}
\hline
\rowcolor{popblueL} \textbf{Feature} & \textbf{Insect-pollinated} & \textbf{Wind-pollinated} \\
\hline
Petals & Large, brightly coloured & Small, dull or absent \\
\hline
Scent and nectar & Scented; nectar produced & No scent, no nectar \\
\hline
Pollen grains & Few, large, sticky or spiny & Many, small, light, smooth \\
\hline
Stigma & Small, sticky, inside flower & Large, feathery, hangs outside \\
\hline
Stamens & Enclosed in flower & Long filaments, anthers hang out \\
\hline
\end{tabularx}
\end{center}
Examples: hibiscus, flamboyant, cocoa (insect-pollinated); maize, grasses, oil palm (wind-pollinated).
\end{methode}

\textbf{Mechanisms favouring cross-pollination.} Many plants reduce self-pollination by:

\begin{itemize}
\item \textbf{Dichogamy}: anthers and stigma of the same flower mature at \emph{different times} (anthers first = \cle{protandry}; stigma first = \cle{protogyny}).
\item \textbf{Self-incompatibility}: pollen landing on the stigma of the same plant fails to germinate or grows too slowly to fertilize (genetic recognition).
\item \textbf{Unisexual flowers}: male and female flowers are separate, either on the same plant (\cle{monoecious}, e.g. maize, oil palm) or on different plants (\cle{dioecious}, e.g. pawpaw, date palm).
\item \textbf{Herkogamy}: physical separation of anthers and stigma within the same flower.
\end{itemize}

\begin{attention}[Pollination is NOT fertilization]
A very common GCE mistake: \textbf{pollination} is only the \emph{transfer of pollen to the stigma}; \textbf{fertilization} is the \emph{fusion of gamete nuclei} which happens later, after the pollen tube has grown down the style. Pollination always \textbf{precedes} fertilization.
\end{attention}

\courssub{Double Fertilization in Angiosperms}

Once a compatible pollen grain lands on the stigma, it germinates: the tube nucleus directs the growth of a \cle{pollen tube} down through the style towards an ovule. This growth is guided by chemical signals (chemotropism) from the synergids. The generative nucleus divides by mitosis into \textbf{two male gamete nuclei} (if not already divided). The pollen tube enters the ovule through the \cle{micropyle}, penetrates one synergid (which degenerates), and releases the two male nuclei into the embryo sac.

\begin{definition}[Double fertilization]
\cle{Double fertilization} is the unique event in angiosperms in which \textbf{one male nucleus ($n$) fuses with the egg cell ($n$)} to form a \textbf{diploid zygote ($2n$)}, while the \textbf{second male nucleus ($n$) fuses with the two polar nuclei ($n + n$) or the fused secondary nucleus ($2n$)} to form a \textbf{triploid endosperm nucleus ($3n$)}.
\end{definition}

The zygote develops into the \textbf{embryo}; the triploid endosperm develops into the \textbf{food store} that nourishes it; the ovule becomes the \textbf{seed} and the ovary becomes the \textbf{fruit} --- the subject of the next section. The synergids and antipodal cells degenerate after fertilization.

\begin{popfigure}
\begin{center}
\begin{tikzpicture}[scale=0.95, font=\small]
% stigma and style
\fill[poppurple] (0,3.6) ellipse (0.5 and 0.22);
\node[poppurple] at (1.6,3.6) {stigma};
\draw[thick, popblue] (-0.35,3.4) -- (-0.35,1.2);
\draw[thick, popblue] (0.35,3.4) -- (0.35,1.2);
\node[popblue] at (1.3,2.3) {style};
% pollen grain on stigma
\draw[thick, poporange] (-0.15,3.85) circle (0.22);
\node[poporange, align=center] at (-2.3,4.0) {pollen grain\\on stigma};
\draw[->, popmuted] (-1.55,3.9) -- (-0.45,3.85);
% pollen tube
\draw[very thick, popgreen] (-0.15,3.65) .. controls (-0.3,2.6) and (0.1,1.8) .. (0.0,0.9);
\node[popgreen, align=center] at (-2.3,2.2) {pollen tube\\growing down};
\draw[->, popmuted] (-1.5,2.2) -- (-0.25,2.5);
% ovary and ovule
\draw[thick, popblue] (0,0.1) ellipse (1.5 and 1.25);
\node[popblue] at (2.4,0.9) {ovary};
\draw[thick, poppink] (0.35,-0.1) ellipse (0.62 and 0.85);
\node[poppink] at (2.4,-0.5) {ovule};
% micropyle
\node[popmuted] at (-2.2,0.75) {micropyle};
\draw[->, popmuted] (-1.5,0.7) -- (-0.35,0.55);
% male nuclei in tube
\fill[popgreen] (0.0,1.15) circle (0.09);
\fill[popgreen] (0.08,0.8) circle (0.09);
\node[popgreen, align=center] at (-2.3,1.15) {two male\\nuclei ($n$)};
\draw[->, popmuted] (-1.5,1.1) -- (-0.15,1.0);
% egg and polar nuclei
\fill[poppurple] (0.2,0.45) circle (0.12);
\node[poppurple] at (2.6,0.35) {egg ($n$)};
\fill[popblue] (0.3,-0.35) circle (0.1);
\fill[popblue] (0.55,-0.35) circle (0.1);
\node[popblue] at (2.75,-0.75) {polar nuclei ($n+n$)};
% fusion arrows and products
\draw[->, very thick, poporange] (0.1,0.75) -- (0.2,0.55);
\draw[->, very thick, poporange] (0.15,0.6) .. controls (0.6,0.2) .. (0.45,-0.25);
\node[popdark, align=center] at (-3.4,-1.3) {\textbf{Fusion 1:} male ($n$) + egg ($n$)\\$\rightarrow$ \textbf{zygote ($2n$)}};
\node[popdark, align=center] at (3.4,-1.9) {\textbf{Fusion 2:} male ($n$) + polar nuclei ($2n$)\\$\rightarrow$ \textbf{endosperm ($3n$)}};
\end{tikzpicture}
\end{center}
\legende{Double fertilization. The pollen tube enters the ovule through the micropyle and discharges two male nuclei: one fuses with the egg to form the diploid zygote, the other with the two polar nuclei to form the triploid endosperm.}
\end{popfigure}

\begin{exoresolu}[Ploidy of embryo sac cells and fertilization products]
A flowering plant has a diploid number $2n = 24$. State the number of chromosomes in: (a) a pollen mother cell; (b) a microspore; (c) the egg cell; (d) a polar nucleus; (e) the zygote; (f) the endosperm nucleus.
\tcblower
Since $2n = 24$, the haploid number is $n = 12$.
\begin{itemize}
\item (a) Pollen mother cell: diploid, \textbf{24}.
\item (b) Microspore (product of meiosis): haploid, \textbf{12}.
\item (c) Egg cell: haploid, \textbf{12}.
\item (d) Each polar nucleus: haploid, \textbf{12}.
\item (e) Zygote ($n + n$): \textbf{24}.
\item (f) Endosperm nucleus ($n + n + n$): triploid, \textbf{36}.
\end{itemize}
\end{exoresolu}

\begin{attention}[GCE exam tip: know your ploidy]
Examiners frequently ask for the ploidy of \emph{every} cell in the embryo sac and its products. Memorise this summary: pollen mother cell and megaspore mother cell $= 2n$; microspores, megaspores, pollen nuclei, egg, synergids, antipodals and polar nuclei $= n$; zygote $= 2n$; endosperm $= 3n$. Note that the endosperm is \textbf{triploid}, never diploid --- this is the classic trap.
\end{attention}

\begin{exercice}[Check your understanding]
\begin{enumerate}
\item Define pollination and distinguish clearly between pollination and fertilization.
\item A maize plant and a hibiscus plant are grown side by side. For each, state the likely pollinating agent and give three floral features that support your answer.
\item Draw and label the embryo sac of a flowering plant just before fertilization, indicating the ploidy of each nucleus.
\item Explain two mechanisms by which a flower may favour cross-pollination, and state the advantage this confers.
\item Describe the fate of the two male nuclei after the pollen tube enters the embryo sac.
\end{enumerate}
\end{exercice}

\begin{corrige}[Exercise 1]
\begin{enumerate}
\item Pollination is the transfer of pollen grains from anther to stigma. Fertilization is the fusion of male and female gamete nuclei; it occurs after pollination, once the pollen tube has reached the embryo sac.
\item Maize: wind-pollinated --- dull flowers without petals or nectar, large feathery stigmas hanging outside, abundant small light pollen. Hibiscus: insect-pollinated --- large colourful petals, nectar and scent, sticky stigma and large sticky pollen grains.
\item The drawing should show an oval sac with: egg cell ($n$) and two synergids ($n$) at the micropylar end, three antipodals ($n$) at the opposite end, and two polar nuclei ($n$ each) in the central cell (often fused as a $2n$ secondary nucleus).
\item Any two of: dichogamy (anthers and stigma mature at different times, preventing self-pollination within a flower); self-incompatibility (own pollen fails to grow on the stigma); unisexual flowers (male and female flowers separated); herkogamy (physical separation). Advantage: cross-pollination mixes genes from two parents, increasing genetic variation and adaptability of the offspring.
\item The generative nucleus (if not already divided) divides by mitosis to form two male nuclei. One fuses with the egg cell to form the diploid zygote ($2n$); the other fuses with the two polar nuclei (or the fused secondary nucleus) to form the triploid endosperm nucleus ($3n$).
\end{enumerate}
\end{corrige}

\begin{aretenir}[Key points]
\begin{itemize}
\item Sexual reproduction is the fusion of haploid gametes to form a diploid zygote; it generates genetic variation.
\item The stamen (anther + filament) is male; the carpel (stigma, style, ovary) is female.
\item Microsporogenesis: pollen mother cell ($2n$) $\xrightarrow{\text{meiosis}}$ 4 microspores ($n$) $\rightarrow$ pollen grain with tube and generative nuclei (microgametogenesis).
\item Megasporogenesis: megaspore mother cell ($2n$) $\xrightarrow{\text{meiosis}}$ 1 functional megaspore ($n$) $\rightarrow$ 7-celled, 8-nucleate embryo sac (megagametogenesis).
\item Pollination = transfer of pollen to the stigma; it precedes fertilization. Cross-pollination increases variation; self-pollination preserves characters.
\item Double fertilization: male ($n$) + egg ($n$) $\rightarrow$ zygote ($2n$); male ($n$) + polar nuclei ($n+n$) $\rightarrow$ endosperm ($3n$).
\end{itemize}
\end{aretenir}

\coursec{Seeds, Dispersal and Germination}

Having explored how pollen grains reach the stigma and how double fertilization gives rise to a diploid zygote and a triploid endosperm, we now follow the remarkable transformation of the fertilized ovule into a seed and the ovary into a fruit. This section connects the reproductive act to the next generation: how the embryo is packaged with food reserves, how the fruit protects and disperses it, and how the dormant embryo springs back to life during germination.

\courssub{From Fertilized Ovule to Seed and Fruit}

Immediately after double fertilization, the flower undergoes dramatic changes. The sepals, petals, stamens, style and stigma usually wither and fall off (they are \emph{deciduous}), having completed their roles. The fate of the remaining parts is summarized in the table below.

\begin{center}
\begin{tabularx}{\linewidth}{|l|X|}\hline
\rowcolor{popblueL} \textbf{Flower part} & \textbf{Fate after fertilization} \\ \hline
Ovule & Becomes the \textbf{seed} \\ \hline
Ovary & Becomes the \textbf{fruit} (pericarp develops from ovary wall) \\ \hline
Zygote (2n) & Develops into the \textbf{embryo} (radicle, plumule, cotyledon(s)) \\ \hline
Primary endosperm nucleus (3n) & Divides to form \textbf{endosperm} (nutritive tissue) \\ \hline
Integuments of ovule & Become the \textbf{testa} (seed coat) \\ \hline
Micropyle & Remains as a small pore in the testa \\ \hline
Funiculus & Becomes the \textbf{stalk} attaching seed to fruit; leaves \textbf{hilum} (scar) when detached \\ \hline
\end{tabularx}
\end{center}

\begin{definition}[Seed]
A seed is a mature, fertilized ovule containing an embryonic plant, a food store (endosperm or cotyledons), and a protective seed coat (testa), capable of developing into a new plant under favourable conditions.
\end{definition}

\begin{definition}[Fruit]
A fruit is a mature, ripened ovary (often with accessory tissues) that encloses the seeds and aids in their protection and dispersal.
\end{definition}

\courssub{Structure of Dicot and Monocot Seeds}

Seeds vary greatly in size, shape and internal organization. The fundamental distinction lies in the number of cotyledons: \cle{dicotyledonous} (two cotyledons) and \cle{monocotyledonous} (one cotyledon). We examine two classic examples: the bean (\emph{Phaseolus vulgaris}) and the maize grain (\emph{Zea mays}).

\begin{popfigure}
\begin{tikzpicture}[scale=0.9, every node/.style={font=\small}]
% Bean seed (dicot) - longitudinal section
\begin{scope}[xshift=0cm]
\draw[popdark, thick] (0,0) ellipse (2.5 and 1.2); % testa outer
\draw[popdark, thick] (0,0) ellipse (2.3 and 1.0); % testa inner
% hilum and micropyle
\draw[popdark, thick] (-2.5,0) -- (-2.8,0) node[left] {Hilum};
\draw[popdark, thick] (2.5,0) -- (2.8,0) node[right] {Micropyle};
% cotyledons
\draw[popgreen!70!popdark, thick, fill=popgreen!20] (-2.3,0) .. controls (-1.5,0.8) and (0,0.8) .. (0,0) .. controls (0,-0.8) and (-1.5,-0.8) .. (-2.3,0);
\draw[popgreen!70!popdark, thick, fill=popgreen!20] (2.3,0) .. controls (1.5,0.8) and (0,0.8) .. (0,0) .. controls (0,-0.8) and (1.5,-0.8) .. (2.3,0);
% embryo axis
\draw[poporange, thick] (-0.2,0) -- (-0.2,0.6) node[above, text width=1.5cm, align=center] {Plumule\\(shoot apex)};
\draw[poporange, thick] (-0.2,0) -- (-0.2,-0.6) node[below, text width=1.5cm, align=center] {Radicle\\(root apex)};
% labels
\node[popdark] at (0,1.5) {\textbf{Bean seed (dicot)}};
\node[popdark] at (-2.3,1.0) {Cotyledon 1};
\node[popdark] at (2.3,1.0) {Cotyledon 2};
\node[popdark] at (0,-1.3) {Testa (seed coat)};
\end{scope}

% Maize grain (monocot) - longitudinal section
\begin{scope}[xshift=8cm]
\draw[popdark, thick] (0,0) ellipse (2.5 and 1.0); % pericarp + testa fused
\draw[popdark, thick] (0,0) ellipse (2.2 and 0.8); % inner
% embryo region
\draw[poporange, thick, fill=poporange!10] (-1.5,-0.2) rectangle (0.5,0.2);
\node[popdark] at (-0.5,0) {Embryo};
% scutellum
\draw[popgreen!70!popdark, thick, fill=popgreen!20] (-1.5,-0.2) .. controls (-1.5,-0.5) and (0.5,-0.5) .. (0.5,-0.2) -- (0.5,0.2) -- (-1.5,0.2) -- cycle;
\node[popdark, text width=2cm, align=center] at (-0.5,-0.7) {Scutellum (cotyledon)};
% endosperm
\draw[popblue!50!popdark, thick, fill=popblue!10] (0.5,-0.8) rectangle (2.2,0.8);
\node[popdark, text width=2cm, align=center] at (1.35,0) {Endosperm\\ (starchy)};
% plumule and radicle
\draw[poporange, thick] (-0.5,0) -- (-0.5,0.6) node[above, text width=1.5cm, align=center] {Plumule};
\draw[poporange, thick] (-0.5,0) -- (-0.5,-0.6) node[below, text width=1.5cm, align=center] {Radicle};
% coleoptile and coleorhiza
\draw[poppurple, thick] (-0.5,0.2) -- (-0.5,0.8) node[above] {Coleoptile};
\draw[poppurple, thick] (-0.5,-0.2) -- (-0.5,-0.8) node[below] {Coleorhiza};
% labels
\node[popdark] at (0,1.4) {\textbf{Maize grain (monocot)}};
\node[popdark] at (0,-1.3) {Pericarp + Testa (fused)};
\end{scope}
\end{tikzpicture}
\legende{Longitudinal sections of a dicot seed (bean) and a monocot grain (maize). Note the large cotyledons storing food in bean, versus the single scutellum absorbing food from the massive endosperm in maize.}
\end{popfigure}

\begin{aretenir}[Key structural differences]
\begin{itemize}
\item \textbf{Bean (dicot):} Food stored in two large \textbf{cotyledons}; endosperm absorbed during seed development. Embryo consists of radicle, plumule and two cotyledons attached to a short hypocotyl.
\item \textbf{Maize (monocot):} Food stored mainly in \textbf{endosperm}; single cotyledon modified into \textbf{scutellum} (absorptive shield). Embryo protected by \textbf{coleoptile} (sheath covering plumule) and \textbf{coleorhiza} (sheath covering radicle). Pericarp (fruit wall) fused to testa.
\end{itemize}
\end{aretenir}

\courssub{Fruit Types (Brief Overview)}

Fruits are classified primarily by the texture of the pericarp (ovary wall) and whether they split open at maturity.

\begin{center}
\begin{tabularx}{\linewidth}{|l|l|X|}\hline
\rowcolor{popblueL} \textbf{Category} & \textbf{Type} & \textbf{Description and Cameroonian examples} \\ \hline
Dry & Dehiscent (split open) & \textbf{Legume/Pod:} splits along two sutures (bean, groundnut, flamboyant). \textbf{Capsule:} splits in various ways (cotton, okra). \\ \hline
Dry & Indehiscent (do not split) & \textbf{Caryopsis (grain):} pericarp fused to testa (maize, rice, millet, guinea corn). \textbf{Achene:} single seed, free pericarp (sunflower). \textbf{Nut:} hard woody pericarp (cashew). \\ \hline
Fleshy & \multicolumn{2}{X|}{\textbf{Berry:} many seeds, fleshy throughout (tomato, guava, papaya). \textbf{Drupe:} one seed enclosed in hard stone (mango, plum, coconut). \textbf{Pome:} fleshy swollen receptacle, true fruit is core (apple — less common in Cameroon).} \\ \hline
\end{tabularx}
\end{center}

\begin{savaistu}[Cameroon context]
The \emph{Flamboyant} (\emph{Delonix regia}) pods explode violently on hot dry days, scattering seeds metres away — a classic example of \textbf{explosive dispersal}. The mango (\emph{Mangifera indica}), a drupe, relies on bats, birds and humans eating the fleshy mesocarp and discarding the stone. Maize (\emph{Zea mays}), a caryopsis, is primarily dispersed by humans (agriculture) but its light chaff can be wind-blown.
\end{savaistu}

\courssub{Seed and Fruit Dispersal}

Dispersal is the scattering of seeds and fruits away from the parent plant. It is crucial for reducing competition among siblings and with the parent for light, water and minerals, and for colonizing new habitats.

\begin{definition}[Dispersal]
Dispersal is the movement, spread or transport of seeds/fruits away from the parent plant, achieved by structural adaptations that exploit abiotic (wind, water) or biotic (animals) agents, or by the plant's own mechanical forces.
\end{definition}

\begin{propriete}[Importance of dispersal]
\begin{enumerate}
\item Reduces intraspecific competition (overcrowding leads to weak, disease-prone seedlings).
\item Enables colonization of new, potentially favourable habitats (range expansion).
\item Reduces predation and pathogen buildup near parent plant.
\item Promotes genetic mixing by separating related individuals.
\end{enumerate}
\end{propriete}

The table below links dispersal agents to structural adaptations. Memorize these pairings for the GCE.

\begin{center}
\begin{tabularx}{\linewidth}{|l|X|X|}\hline
\rowcolor{popblueL} \textbf{Agent} & \textbf{Adaptation} & \textbf{Examples (Cameroon)} \\ \hline
Wind (Anemochory) & \textbf{Winged} (samara), \textbf{plumed/hairy} (pappus), \textbf{light/dust-like}, \textbf{balloon-like} inflated calyx & Flamboyant (winged), \emph{Ceiba pentandra} (kapok — hairy), \emph{Combretum} (winged), orchids (dust seeds) \\ \hline
Water (Hydrochory) & \textbf{Buoyant}, \textbf{waterproof}, \textbf{fibrous/spongy} mesocarp/husk, \textbf{air spaces} & Coconut (fibrous husk), water lily (spongy fruit), mangrove propagules (buoyant) \\ \hline
Animals (Zoochory) & \textbf{Edible fleshy} fruit (endozoochory), \textbf{hooks/spines} (epizoochory), \textbf{sticky} mucilage & Mango, papaya, bush mango (\emph{Irvingia}) — fleshy; \emph{Desmodium} (beggar's lice), \emph{Xanthium} (cocklebur) — hooks \\ \hline
Explosive / Self (Autochory) & \textbf{Tension-building} fruit walls that snap, \textbf{coiling} pods, \textbf{explosive} dehiscence & Flamboyant, \emph{Impatiens} (balsam), \emph{Hura crepitans} (sandbox tree) \\ \hline
\end{tabularx}
\end{center}

\begin{popfigure}
\begin{tikzpicture}[scale=0.85, every node/.style={font=\small}]
% Wind
\begin{scope}[xshift=0cm]
\draw[popblue, thick, fill=popblue!10, rounded corners] (0,0) rectangle (3,3.5);
\node[popdark] at (1.5,3.2) {\textbf{Wind}};
\draw[popdark] (1.5,2.7) -- (1.5,2.3);
% samara
\draw[popgreen!60!popdark, thick] (1.5,2.0) -- (0.5,1.0) -- (1.5,0.5) -- (2.5,1.0) -- cycle;
\node[popdark] at (1.5,1.2) {Winged (samara)};
% pappus
\draw[popdark, fill=popgray!30] (1.5,0.2) circle (0.3);
\foreach \a in {0,30,...,330} \draw[popdark, thin] (1.5,0.2) -- +(\a:0.5);
\node[popdark] at (1.5,-0.4) {Plumed (pappus)};
\end{scope}

% Water
\begin{scope}[xshift=4cm]
\draw[popblue, thick, fill=popblue!10, rounded corners] (0,0) rectangle (3,3.5);
\node[popdark] at (1.5,3.2) {\textbf{Water}};
% coconut
\draw[poporange!70!popdark, thick, fill=poporange!20] (1.5,2.0) ellipse (0.8 and 0.5);
\draw[poporange!70!popdark, thick] (0.7,2.0) -- (0.5,1.5) -- (0.7,1.0);
\draw[poporange!70!popdark, thick] (2.3,2.0) -- (2.5,1.5) -- (2.3,1.0);
\node[popdark] at (1.5,1.2) {Fibrous husk};
% mangrove propagule
\draw[popgreen!60!popdark, thick] (1.5,0.5) -- (1.5,-0.2);
\draw[popgreen!60!popdark, thick, fill=popgreen!20] (1.2,0.5) rectangle (1.8,0.2);
\node[popdark] at (1.5,0.0) {Buoyant propagule};
\end{scope}

% Animal
\begin{scope}[xshift=8cm]
\draw[popblue, thick, fill=popblue!10, rounded corners] (0,0) rectangle (3,3.5);
\node[popdark] at (1.5,3.2) {\textbf{Animal}};
% fleshy fruit
\draw[poppink, thick, fill=poppink!20] (1.5,2.2) ellipse (0.7 and 0.5);
\node[popdark] at (1.5,2.2) {Fleshy};
\node[popdark] at (1.5,1.7) {Edible};
% hook
\draw[popdark, thick] (1.5,1.0) -- (1.5,0.2);
\draw[popdark, thick] (1.3,0.6) .. controls (1.0,0.6) and (1.0,0.3) .. (1.3,0.3);
\draw[popdark, thick] (1.7,0.6) .. controls (2.0,0.6) and (2.0,0.3) .. (1.7,0.3);
\node[popdark] at (1.5,0.0) {Hooked/Spiny};
\end{scope}

% Explosive
\begin{scope}[xshift=12cm]
\draw[popblue, thick, fill=popblue!10, rounded corners] (0,0) rectangle (3,3.5);
\node[popdark] at (1.5,3.2) {\textbf{Explosive}};
% pod splitting
\draw[popgreen!60!popdark, thick, fill=popgreen!20] (0.5,2.0) .. controls (0.5,2.5) and (2.5,2.5) .. (2.5,2.0) .. controls (2.5,1.5) and (0.5,1.5) .. (0.5,2.0);
\draw[popgreen!60!popdark, thick, fill=popgreen!20] (0.5,2.0) .. controls (0.5,1.5) and (2.5,1.5) .. (2.5,2.0) .. controls (2.5,2.5) and (0.5,2.5) .. (0.5,2.0);
\draw[popdark, thick, dashed] (1.5,2.0) -- (1.5,1.0);
\node[popdark] at (1.5,1.2) {Tension splits pod};
% seeds flying
\foreach \i in {1,...,5} {
  \pgfmathsetmacro{\ang}{-90+rand*30}
  \pgfmathsetmacro{\len}{0.5+rand*0.5}
  \draw[poporange, thick, ->] (1.5,1.0) -- +(\ang:\len);
}
\node[popdark] at (1.5,0.2) {Seeds ejected};
\end{scope}
\end{tikzpicture}
\legende{Structural adaptations for dispersal by wind, water, animals and explosive mechanisms.}
\end{popfigure}

\begin{attention}[GCE Exam Tip]
When asked to "state the agent of dispersal and the adaptation", give a \textbf{single sentence} linking them explicitly. Example: "The fruit is dispersed by wind because it has a winged pericarp (samara) that increases air resistance and allows it to be carried far from the parent plant." Do not just list "wind" and "wing" separately.
\end{attention}

\courssub{Germination: Definition, Conditions and Biochemistry}

\begin{definition}[Germination]
Germination is the resumption of active growth of the embryo in a seed after a period of dormancy, resulting in the emergence of the radicle and plumule and the establishment of a seedling, provided favourable environmental conditions are met.
\end{definition}

\begin{propriete}[Conditions necessary for germination]
Three external factors are \textbf{essential}:
\begin{enumerate}
\item \textbf{Water:} Rehydrates cells (imbibition), activates hydrolytic enzymes, provides medium for metabolic reactions, and ruptures the testa.
\item \textbf{Oxygen:} Required for aerobic respiration to release energy (ATP) from food reserves. Most seeds cannot germinate anaerobically (rice is an exception).
\item \textbf{Suitable temperature:} Enzymes have an optimum temperature range (typically \SI{25}{\ensuremath{{}^{\circ}\mathrm{C}}}--\SI{35}{\ensuremath{{}^{\circ}\mathrm{C}}} for tropical crops). Too low: enzymes inactive; too high: enzymes denature.
\end{enumerate}
\textbf{Internal factors:} Viable embryo, non-dormant state (dormancy broken), sufficient food reserves.
\end{propriete}

\begin{methode}[Mobilisation of food reserves during germination]
\begin{enumerate}
\item \textbf{Imbibition:} Water enters via micropyle; cells swell, testa ruptures.
\item \textbf{Hormonal signal:} Embryo produces \textbf{gibberellins} (GA) which diffuse to the aleurone layer (in cereals) or cotyledons (in dicots).
\item \textbf{Enzyme synthesis:} GA stimulates synthesis of hydrolytic enzymes: \textbf{amylases} (starch $\to$ maltose/glucose), \textbf{proteases} (proteins $\to$ amino acids), \textbf{lipases} (lipids $\to$ fatty acids + glycerol).
\item \textbf{Translocation:} Soluble products transported via vascular tissue to growing points (radicle, plumule).
\item \textbf{Utilization:} Respiration (glycolysis, Krebs cycle, oxidative phosphorylation) provides ATP and carbon skeletons for new cell synthesis (cell division in meristems, cell elongation).
\end{enumerate}
\end{methode}

\begin{attention}[Common mistake]
Do not confuse \textbf{imbibition} (physical water uptake) with \textbf{germination} (metabolic activation and growth). A dead seed can imbibe water but will not germinate. Also, light is \textbf{not} a universal requirement; some seeds need light (lettuce), some need darkness, most are indifferent.
\end{attention}

\courssub{Epigeal vs Hypogeal Germination}

The behaviour of the cotyledons during germination defines two major patterns.

\begin{center}
\begin{tabularx}{\linewidth}{|l|X|X|}\hline
\rowcolor{popblueL} \textbf{Feature} & \textbf{Epigeal germination} & \textbf{Hypogeal germination} \\ \hline
\textbf{Cotyledons} & Pulled \textbf{above ground} by elongating \textbf{hypocotyl}; become photosynthetic (green). & Remain \textbf{below ground}; \textbf{epicotyl} elongates instead. Cotyledons stay in soil, non-photosynthetic. \\ \hline
\textbf{Hypocotyl} & Elongates vigorously, forms a hook that pulls cotyledons up. & Remains short; does not elongate significantly. \\ \hline
\textbf{Epicotyl} & Short initially; elongates later after cotyledons emerge. & Elongates vigorously, pushing plumule above ground. \\ \hline
\textbf{Protection of plumule} & Cotyledons enclose plumule; may be damaged during emergence. & Plumule protected by \textbf{coleoptile} (monocots) or by remaining inside cotyledons (dicots like pea). \\ \hline
\textbf{Examples} & Bean, castor oil, sunflower, cotton, flamboyant. & Maize, rice, millet, guinea corn, pea, groundnut, coconut. \\ \hline
\end{tabularx}
\end{center}

\begin{popfigure}
\begin{tikzpicture}[scale=0.8, every node/.style={font=\small}]
% Epigeal (Bean) - 4 stages
\begin{scope}[xshift=0cm]
\node[popdark] at (2,5.5) {\textbf{Epigeal (Bean)}};
% Stage 1: Dry seed
\draw[popdark, thick] (0.5,4.5) ellipse (0.8 and 0.4);
\node at (2,4.5) {Dry seed};
% Stage 2: Radicle emergence
\draw[popdark, thick] (0.5,3.5) ellipse (0.8 and 0.4);
\draw[poporange, thick] (0.5,3.1) -- (0.5,2.5);
\node at (2,3.5) {Radicle out};
% Stage 3: Hypocotyl hook
\draw[popdark, thick] (0.5,2.5) .. controls (0.5,2.0) and (1.5,2.0) .. (1.5,1.5);
\draw[popgreen!60!popdark, thick, fill=popgreen!20] (1.5,1.5) .. controls (1.5,1.0) and (2.5,1.0) .. (2.5,1.5) .. controls (2.5,2.0) and (1.5,2.0) .. (1.5,1.5);
\node at (2,2.5) {Hook stage};
% Stage 4: Seedling
\draw[popdark, thick] (0.5,1.0) -- (0.5,0.0); % root
\draw[popdark, thick] (0.5,1.0) -- (0.5,-0.5); % hypocotyl straight
\draw[popgreen!60!popdark, thick, fill=popgreen!20] (0.5,-0.5) .. controls (0.5,-1.0) and (1.5,-1.0) .. (1.5,-0.5) .. controls (1.5,0.0) and (0.5,0.0) .. (0.5,-0.5);
\draw[popgreen!60!popdark, thick, fill=popgreen!20] (0.5,-0.5) .. controls (0.5,-1.0) and (-0.5,-1.0) .. (-0.5,-0.5) .. controls (-0.5,0.0) and (0.5,0.0) .. (0.5,-0.5);
\draw[poporange, thick] (0.5,-0.5) -- (0.5,-1.2) node[below] {Plumule};
\node at (2,0.5) {Seedling};
\end{scope}

% Hypogeal (Maize) - 4 stages
\begin{scope}[xshift=7cm]
\node[popdark] at (2,5.5) {\textbf{Hypogeal (Maize)}};
% Stage 1: Dry grain
\draw[popdark, thick] (0.5,4.5) ellipse (0.8 and 0.35);
\node at (2,4.5) {Dry grain};
% Stage 2: Radicle + coleorhiza
\draw[popdark, thick] (0.5,3.5) ellipse (0.8 and 0.35);
\draw[poporange, thick] (0.5,3.15) -- (0.5,2.5);
\draw[poppurple, thick] (0.5,3.15) -- (0.5,2.8) node[right] {Coleorhiza};
\node at (2,3.5) {Radicle out};
% Stage 3: Coleoptile emergence
\draw[popdark, thick] (0.5,2.5) ellipse (0.8 and 0.35);
\draw[poppurple, thick] (0.5,2.15) -- (0.5,1.0) node[above right] {Coleoptile};
\node at (2,2.5) {Coleoptile up};
% Stage 4: Seedling
\draw[popdark, thick] (0.5,1.0) -- (0.5,0.0); % root
\draw[poppurple, thick] (0.5,1.0) -- (0.5,-0.5); % coleoptile
\draw[poporange, thick] (0.5,-0.5) -- (0.5,-1.2); % plumule through coleoptile
\draw[popgreen!60!popdark, thick] (0.5,-0.5) -- (1.0,-0.8) node[right] {First leaf};
\node at (2,0.0) {Seedling};
% Cotyledon (scutellum) remains underground
\draw[popgreen!60!popdark, thick, fill=popgreen!20] (0.5,2.5) ellipse (0.8 and 0.35);
\node[popdark, rotate=90] at (-0.5,2.5) {Scutellum (below ground)};
\end{scope}
\end{tikzpicture}
\legende{Comparison of epigeal (bean) and hypogeal (maize) germination stages. Note the elongating hypocotyl pulling cotyledons up in bean, versus the elongating epicotyl/coleoptile pushing plumule up while cotyledon (scutellum) stays below in maize.}
\end{popfigure}

\begin{attention}[GCE Trap]
In hypogeal germination, the \textbf{cotyledon(s) remain underground}. Do not draw them above ground! In maize, the scutellum never emerges; in pea (dicot but hypogeal), the cotyledons stay below soil. The key is: \emph{which axis elongates?} Hypocotyl $\to$ epigeal; Epicotyl $\to$ hypogeal.
\end{attention}

\courssub{Experimental Investigation: Conditions for Germination}

The classic controlled experiment uses four test tubes (or boiling tubes) to isolate each condition.

\begin{methode}[Four-tube experiment to demonstrate requirements for germination]
\begin{enumerate}
\item \textbf{Set up four tubes} with cotton wool at the bottom and 10--20 viable seeds of the same species (e.g., bean) on top.
\item \textbf{Tube A (Control):} Cotton wool moistened with water; tube open to air; kept at room temperature (\SI{\sim 25}{\ensuremath{{}^{\circ}\mathrm{C}}}). \emph{All conditions present.}
\item \textbf{Tube B (No water):} Cotton wool dry; tube open; room temperature. \emph{Water absent.}
\item \textbf{Tube C (No oxygen):} Cotton wool moistened with \textbf{boiled and cooled water} (drives out dissolved $\ce{O2}$); add a layer of \textbf{liquid paraffin} or \textbf{pyrogallol + $\ce{NaOH}$} to prevent atmospheric $\ce{O2}$ diffusion; room temperature. \emph{Oxygen absent.}
\item \textbf{Tube D (Unsuitable temperature):} Cotton wool moistened; tube open; placed in a \textbf{refrigerator} (\SI{\sim 4}{\ensuremath{{}^{\circ}\mathrm{C}}}) or an incubator at \SI{45}{\ensuremath{{}^{\circ}\mathrm{C}}}. \emph{Temperature unsuitable.}
\item \textbf{Observe daily for 5--7 days.} Record which tubes show germination (radicle emergence).
\end{enumerate}
\end{methode}

\begin{exoresolu}[Interpreting germination experiment results]
\textbf{Statement:} In the four-tube experiment described above, after 5 days at room temperature, seeds in Tube A have germinated (radicles \SI{10}{mm} long). Tubes B, C and D show no radicle emergence. Seeds from Tube B are transferred to moist cotton wool at room temperature; they germinate within 2 days. Seeds from Tube C are transferred to moist cotton wool open to air; they germinate within 2 days. Seeds from Tube D are transferred to room temperature; they germinate within 2 days.

\textbf{Questions:}
\begin{enumerate}
\item State the condition being tested in each of Tubes B, C and D.
\item Explain why boiled water cooled and covered with paraffin is used in Tube C.
\item What conclusion can be drawn about the viability of seeds in Tubes B, C and D?
\item Why is it important to use seeds from the same batch and same species?
\end{enumerate}

\tcblower
\textbf{Detailed solution:}
\begin{enumerate}
\item \textbf{Tube B:} Absence of water (dry cotton wool). \textbf{Tube C:} Absence of oxygen (boiled water + paraffin excludes $\ce{O2}$). \textbf{Tube D:} Unsuitable (low) temperature (refrigerator).
\item Boiling water removes dissolved oxygen. The layer of liquid paraffin (or oil) creates a physical barrier preventing atmospheric oxygen from dissolving back into the water. This ensures anaerobic conditions.
\item The seeds in Tubes B, C and D are \textbf{viable} (alive) because they germinated once the missing condition was supplied. They were dormant due to the lack of that specific factor, not dead.
\item Using the same batch and species controls \textbf{genetic variability} and \textbf{seed age/size}, ensuring that any difference in germination is due only to the manipulated environmental factor (validity of the controlled experiment).
\end{enumerate}
\end{exoresolu}

\courssub{Worked Example: Linking Structure to Function}

\begin{exoresolu}[Seed structure and dispersal adaptation]
\textbf{Statement:} The diagram below shows a longitudinal section of a fruit of \emph{Xanthium} (cocklebur). The fruit is a dry, indehiscent achene covered in stiff, hooked spines. The seed inside has a thin testa and a large, oil-rich embryo.

\begin{enumerate}
\item Name the agent of dispersal for this fruit.
\item Explain how the structure of the fruit adapts it to this agent.
\item The seed germinates epigeally. Describe the role of the hypocotyl during germination.
\item State two environmental conditions necessary for the germination of this seed.
\end{enumerate}

\tcblower
\textbf{Detailed solution:}
\begin{enumerate}
\item \textbf{Agent:} Animal (epizoochory / external attachment to fur/feathers/clothing).
\item \textbf{Adaptation:} The fruit wall (pericarp) is modified into a \textbf{hard, woody achene} covered with \textbf{stiff, hooked spines}. These hooks catch onto the fur of passing mammals (or feathers/clothing of humans), carrying the fruit away from the parent plant. The fruit is indehiscent, so the seed remains protected inside until it falls off or the fruit decays.
\item \textbf{Role of hypocotyl:} In epigeal germination, the \textbf{hypocotyl elongates rapidly}, forming a hook that pushes through the soil. The hook straightens in response to light (phototropism), \textbf{pulling the cotyledons above ground}. The cotyledons then expand, become green and photosynthetic, providing energy for the seedling until true leaves form.
\item \textbf{Conditions:} (1) \textbf{Water} (for imbibition and enzyme activation). (2) \textbf{Oxygen} (for aerobic respiration). (3) \textbf{Suitable temperature} (approx. \SI{25}{\ensuremath{{}^{\circ}\mathrm{C}}}--\SI{30}{\ensuremath{{}^{\circ}\mathrm{C}}}). (Any two correctly stated.)
\end{enumerate}
\end{exoresolu}

\courssub{Summary of Key Points}

\begin{aretenir}[Seeds, Dispersal and Germination — Essentials for GCE]
\begin{itemize}
\item \textbf{Seed formation:} Ovule $\to$ seed; Ovary $\to$ fruit; Zygote $\to$ embryo; Endosperm nucleus $\to$ endosperm; Integuments $\to$ testa; Funiculus $\to$ stalk/hilum.
\item \textbf{Dicot seed (bean):} Testa, hilum, micropyle, two large cotyledons (food store), embryo (radicle, plumule, hypocotyl). No endosperm.
\item \textbf{Monocot grain (maize):} Pericarp+testa fused, scutellum (single cotyledon, absorptive), large starchy endosperm, embryo with coleoptile \& coleorhiza.
\item \textbf{Fruit types:} Dry dehiscent (pod), dry indehiscent (caryopsis, nut), fleshy (berry, drupe). Local examples: bean pod, maize grain, mango, flamboyant.
\item \textbf{Dispersal agents \& adaptations:} Wind — wing/hair; Water — fibrous husk/buoyancy; Animal — fleshy edible / hooks; Explosive — tension-splitting pods.
\item \textbf{Importance of dispersal:} Reduces competition, colonizes new habitats, reduces predation/disease.
\item \textbf{Germination conditions:} Water, Oxygen, Suitable Temperature. (Light not universal).
\item \textbf{Biochemistry:} Imbibition $\to$ GA synthesis $\to$ enzyme production (amylase, protease, lipase) $\to$ hydrolysis of reserves $\to$ translocation $\to$ respiration/growth.
\item \textbf{Epigeal:} Hypocotyl elongates, cotyledons pulled above ground, become photosynthetic (bean).
\item \textbf{Hypogeal:} Epicotyl elongates, cotyledons remain underground (maize, pea).
\item \textbf{Experiment:} Four tubes (Control, -Water, -Oxygen, -Temp) prove all three factors essential. Viability confirmed by recovery.
\end{itemize}
\end{aretenir}

\begin{exercice}[Practice Questions]
\begin{enumerate}
\item Draw a large, labelled diagram of the longitudinal section of a maize grain. Indicate the positions of the scutellum, coleoptile, coleorhiza, endosperm and pericarp.
\item A student sets up the four-tube germination experiment but forgets to boil the water in Tube C. What would be the likely result in Tube C, and why?
\item Distinguish clearly between the terms \emph{imbibition}, \emph{germination} and \emph{seedling establishment}.
\item The fruit of \emph{Triumfetta} (a common weed in Cameroon) has numerous hooked bristles. Name the dispersal agent and explain the advantage of this dispersal method for the species.
\item In an experiment, bean seedlings grown in the dark have long, pale hypocotyls and small, unexpanded cotyledons. Seedlings grown in light have short, green hypocotyls and large, expanded green cotyledons. Explain these observations.
\end{enumerate}
\end{exercice}

\begin{corrige}[Exercise 1]
Maize grain LS: Oval outline (pericarp+testa fused). Large central endosperm (starchy). Embryo at one side: shield-shaped scutellum pressed against endosperm. Plumule covered by coleoptile (sheath). Radicle covered by coleorhiza (sheath). Labels: Pericarp/Testa, Endosperm, Scutellum, Coleoptile, Plumule, Coleorhiza, Radicle.
\end{corrige}

\begin{corrige}[Exercise 2]
Tube C would likely show germination because unboiled water contains dissolved oxygen. The paraffin layer alone does not remove oxygen already dissolved; it only prevents further diffusion. Boiling is essential to drive out dissolved $\ce{O2}$.
\end{corrige}

\begin{corrige}[Exercise 3]
\textbf{Imbibition:} Physical absorption of water by dry seed (colloidal matrix swells), passive, occurs in dead seeds too. \textbf{Germination:} Metabolic activation of embryo, enzyme synthesis, hydrolysis of reserves, cell division/elongation, radicle emergence — requires viable seed. \textbf{Seedling establishment:} Post-germination phase where seedling becomes autotrophic (photosynthetic), true leaves form, root system develops, reliance on seed reserves ends.
\end{corrige}

\begin{corrige}[Exercise 4]
Agent: Animal (epizoochory). Advantage: Hooked bristles attach firmly to fur of passing animals (antelopes, cattle, humans), transporting fruits far from parent. Reduces competition among dense seedlings under parent plant; colonizes new areas like paths, forest edges, farmland.
\end{corrige}

\begin{corrige}[Exercise 5]
Dark-grown seedlings show \textbf{etiolation}: hypocotyl elongates rapidly (skotomorphogenesis) to reach light, but lacks chlorophyll (pale/yellow) because chlorophyll synthesis requires light (protochlorophyllide $\to$ chlorophyllide needs photons). Cotyledons remain small and unexpanded (no photosynthetic demand). Light-grown seedlings: light inhibits hypocotyl elongation (photomorphogenesis), promotes chlorophyll synthesis (green), and stimulates cotyledon expansion for photosynthesis.
\end{corrige}

\coursec{The Male Reproductive System and Spermatogenesis}

In the previous sections we followed the flowering plant from pollination to the germination of a new seedling. We now turn to reproduction in humans, and we begin with the male partner. Every day, from puberty until old age, a healthy man produces several hundred million spermatozoa — yet each ejaculate releases cells so small that about two thousand of them could fit on the head of a pin. How is such a prodigious, continuous output organised? The answer lies in a beautifully coordinated system of organs, tubes, glands and hormones, which we now study in detail.

\courssub{Overview of the Male Reproductive System}

\begin{definition}[Male reproductive system]
The \cle{male reproductive system} is the set of organs that produce male gametes (\cle{spermatozoa}), nourish and store them, and deliver them into the female reproductive tract during copulation. It comprises the \textbf{testes} (held in the \textbf{scrotum}), a system of ducts (\textbf{epididymis}, \textbf{vas deferens}, \textbf{urethra}), accessory glands (\textbf{seminal vesicles}, \textbf{prostate gland}, \textbf{Cowper's glands}) and the \textbf{penis}.
\end{definition}

\begin{popfigure}
\begin{center}
\begin{tikzpicture}[scale=1.0, font=\small]
% Bladder
\draw[fill=popblueL, draw=popblue, thick] (3.4,3.4) ellipse (0.9 and 0.75);
\node[popdark] at (3.4,3.4) {Bladder};
% Seminal vesicle
\draw[fill=poppink!40, draw=poppink, thick] (4.6,2.2) ellipse (0.55 and 0.3);
\node[popdark, align=center] at (6.3,2.2) {Seminal\\vesicle};
\draw[->, popmuted] (5.7,2.2) -- (5.15,2.2);
% Prostate
\draw[fill=popgold!40, draw=popgold, thick] (3.4,1.6) circle (0.42);
\node[popdark] at (5.6,1.55) {Prostate gland};
\draw[->, popmuted] (5.0,1.55) -- (3.85,1.58);
% Cowper's gland
\draw[fill=popgreenL, draw=popgreen, thick] (3.0,0.85) circle (0.18);
\node[popdark, align=center] at (5.6,0.8) {Cowper's\\gland};
\draw[->, popmuted] (4.9,0.8) -- (3.2,0.85);
% Vas deferens
\draw[very thick, poporange] (1.6,0.4) .. controls (1.2,1.6) and (2.2,3.2) .. (3.0,3.9) .. controls (3.6,4.3) .. (4.0,3.4) .. controls (4.2,2.8) .. (3.9,2.2) .. controls (3.7,1.9) .. (3.55,1.85);
\node[poporange, align=center] at (0.4,2.9) {Vas\\deferens};
\draw[->, popmuted] (0.75,2.6) -- (1.35,2.2);
% Urethra
\draw[very thick, poppurple] (3.4,1.2) .. controls (3.2,0.4) and (2.6,0.1) .. (2.2,-0.1) -- (5.6,-0.35);
\node[poppurple] at (6.6,-0.35) {Urethra};
\draw[->, popmuted] (6.0,-0.35) -- (5.6,-0.35);
% Penis
\draw[fill=popblueL!50, draw=popdark, thick, rounded corners=6pt] (2.0,-0.75) rectangle (5.8,0.15);
\node[popdark] at (6.9,-1.0) {Penis};
\draw[->, popmuted] (6.3,-0.85) -- (4.6,-0.72);
% Erectile tissue hint
\draw[popmuted, dashed] (2.2,0.05) -- (5.5,0.05);
% Testis and scrotum
\draw[fill=popgreenL, draw=popgreen, very thick] (1.5,-0.4) ellipse (0.55 and 0.75);
\node[popdark] at (0.1,-0.4) {Testis};
\draw[->, popmuted] (0.65,-0.4) -- (0.95,-0.4);
\draw[draw=popdark, thick] (0.85,-1.35) .. controls (0.8,-0.2) and (1.1,0.35) .. (1.5,0.38) .. controls (1.95,0.35) and (2.25,-0.3) .. (2.2,-1.35) .. controls (1.9,-1.6) and (1.15,-1.6) .. (0.85,-1.35);
\node[popdark] at (0.1,-1.5) {Scrotum};
\draw[->, popmuted] (0.6,-1.45) -- (1.1,-1.4);
% Epididymis
\draw[very thick, poppink] (1.95,0.15) .. controls (2.15,-0.4) .. (1.95,-1.0);
\node[poppink, align=center] at (3.1,-1.35) {Epididymis};
\draw[->, popmuted] (2.85,-1.25) -- (2.05,-0.7);
\end{tikzpicture}
\end{center}
\legende{Simplified sagittal section of the male reproductive system, showing the testis in the scrotum, the sperm ducts and the accessory glands.}
\end{popfigure}

\courssub{Structure and Functions of the Organs}

\textbf{The testes and the scrotum.}\par
The two \cle{testes} are the male gonads: they produce spermatozoa and secrete the male sex hormone \textbf{testosterone}. Each testis is an oval organ about \SI{4}{cm} long and \SI{2.5}{cm} wide, suspended outside the abdominal cavity in a skin pouch, the \cle{scrotum}. This external position is not accidental. Spermatogenesis is extremely temperature-sensitive and proceeds normally only at about \SI{2}{\ensuremath{{}^{\circ}\mathrm{C}}}--\SI{3}{\ensuremath{{}^{\circ}\mathrm{C}}} \emph{below} core body temperature (i.e. around \SI{34}{\ensuremath{{}^{\circ}\mathrm{C}}}--\SI{35}{\ensuremath{{}^{\circ}\mathrm{C}}} instead of \SI{37}{\ensuremath{{}^{\circ}\mathrm{C}}}). The scrotum acts as a thermostat: in cold weather its muscles (dartos and cremaster) contract and draw the testes closer to the warm body; in hot weather it relaxes so the testes hang further away and lose heat. This is why tight clothing that keeps the testes pressed against the body can temporarily reduce sperm production — a fact relevant to reproductive health advice in our climate.

\textbf{The epididymis.}\par
A long, highly coiled tube (about \SI{6}{m} if uncoiled) lying on the posterior surface of each testis. Newly formed spermatozoa are immature and cannot swim; during their journey through the epididymis, which takes \SI{10}{--}\SI{14}{days}, they \textbf{mature}, gain motility and are \textbf{stored} until ejaculation. The epididymis also absorbs excess fluid and phagocytoses defective sperm.

\textbf{The vas deferens (sperm duct).}\par
A muscular tube, about \SI{45}{cm} long, that carries spermatozoa from the epididymis up into the abdomen, looping over the ureter and the bladder to join the urethra at the ejaculatory duct. During ejaculation, waves of \textbf{peristalsis} in its thick muscular wall propel the sperm forward. Cutting and tying the vas deferens is the principle of male sterilisation (vasectomy), studied in the section on birth control.

\textbf{The accessory glands.}\par
Spermatozoa alone would form a tiny, sticky mass. Three sets of glands add fluids that together with sperm form \cle{semen} (typical volume \SI{2}{--}\SI{5}{ml} per ejaculate):
\begin{itemize}
\item the \textbf{seminal vesicles} (paired, \SI{5}{cm} long) secrete an alkaline fluid (\textasciitilde\SI{60}{\%} of semen volume) rich in \textbf{fructose} (the main energy source for sperm swimming), prostaglandins (which stimulate uterine contractions to help sperm transport), and clotting proteins;
\item the \textbf{prostate gland} (single, chestnut-sized) secretes a milky, slightly acidic fluid (\textasciitilde\SI{30}{\%} of volume) containing citrate, enzymes (including prostate-specific antigen, PSA) and zinc; it helps \textbf{neutralise the acidity} of the female vagina, which would otherwise kill the spermatozoa;
\item the \textbf{Cowper's (bulbo-urethral) glands} (paired, pea-sized) secrete a clear, viscous, alkaline mucus before ejaculation, which lubricates the urethra and neutralises traces of acidic urine.
\end{itemize}

\textbf{The urethra and the penis.}\par
The \textbf{urethra} (\SI{18}{--}\SI{20}{cm} long) is the common passage for urine and semen (never at the same time: the internal urethral sphincter closes the bladder during ejaculation). The \textbf{penis} contains three columns of \textbf{erectile tissue}: two \textbf{corpora cavernosa} dorsally and one \textbf{corpus spongiosum} ventrally (which expands to form the glans penis). During sexual arousal, parasympathetic stimulation causes arterial dilation and venous compression, filling these spaces with blood and making the organ erect so it can be inserted into the vagina to deposit semen near the cervix.

\begin{aretenir}[Functions at a glance]
\begin{center}
\begin{tabularx}{\linewidth}{|l|X|}\hline
\rowcolor{popblueL} \textbf{Organ} & \textbf{Function} \\ \hline
Testis & Produces spermatozoa and testosterone \\ \hline
Scrotum & Keeps testes \SI{2}{\ensuremath{{}^{\circ}\mathrm{C}}}--\SI{3}{\ensuremath{{}^{\circ}\mathrm{C}}} below body temperature \\ \hline
Epididymis & Maturation, storage and transport of spermatozoa \\ \hline
Vas deferens & Transports sperm by peristalsis during ejaculation \\ \hline
Seminal vesicles & Secrete fructose-rich fluid (energy), prostaglandins, clotting proteins \\ \hline
Prostate gland & Secretes fluid with enzymes, citrate, zinc; neutralises vaginal acidity \\ \hline
Cowper's glands & Lubrication; neutralise urine traces in urethra \\ \hline
Penis / urethra & Copulation and delivery of semen \\ \hline
\end{tabularx}
\end{center}
\end{aretenir}

\courssub{Microscopic Structure of the Testis}

If we cut a thin section of a testis and observe it under the microscope, we see that it is packed with hundreds of tightly coiled tubes, the \cle{seminiferous tubules} (total length \textasciitilde\SI{250}{m} per testis). It is inside these tubules that spermatozoa are manufactured. Between the tubules lie clusters of \cle{interstitial cells} (also called \textbf{Leydig cells}), which secrete \textbf{testosterone}. Blood capillaries and lymphatic vessels are also found in this interstitial tissue.

The wall of each seminiferous tubule contains two important cell types:
\begin{itemize}
\item the \textbf{germinal epithelial cells} (spermatogonia, primary/secondary spermatocytes, spermatids) at different stages of division, which give rise to spermatozoa (the youngest cells lie at the outer edge near the basal lamina, the most mature near the central cavity or lumen);
\item the large \cle{Sertoli cells} (sustentacular cells), which stretch from the basal lamina to the lumen. They \textbf{nourish and support} the developing sperm cells, phagocytose excess cytoplasm, form the \textbf{blood-testis barrier} (protecting developing sperm from immune attack), and secrete the hormone \textbf{inhibin} and \textbf{androgen-binding protein} (ABP) which concentrates testosterone in the tubule.
\end{itemize}
Mature spermatozoa are released into the lumen of the tubule (spermiation) and carried by fluid towards the \textbf{rete testis}, then to the epididymis.

\begin{popfigure}
\begin{center}
\begin{tikzpicture}[font=\small, scale=1.1]
% Seminiferous tubule cross-section
\draw[fill=popblueL!30, draw=popblue, thick] (0,0) circle (2.5);
\draw[fill=white] (0,0) circle (2.0);
% Basal lamina
\draw[popdark, thick] (0,0) circle (2.0);
% Sertoli cells
\foreach \angle/\rad in {30/1.8, 90/1.8, 150/1.8, 210/1.8, 270/1.8, 330/1.8} {
  \draw[fill=popgreen!60, draw=popgreen, thick] (\angle:\rad) ellipse (0.25 and 0.6);
  \node[popdark, font=\tiny] at (\angle:\rad) {Sertoli};
}
% Germ cells at different stages
\foreach \angle/\rad in {10/1.5, 50/1.5, 130/1.5, 170/1.5, 250/1.5, 290/1.5} {
  \draw[fill=poporange!60, draw=poporange] (\angle:\rad) circle (0.15);
}
\foreach \angle/\rad in {70/1.2, 110/1.2, 190/1.2, 230/1.2, 310/1.2, 350/1.2} {
  \draw[fill=poppink!60, draw=poppink] (\angle:\rad) circle (0.12);
}
\foreach \angle/\rad in {20/0.8, 80/0.8, 140/0.8, 200/0.8, 260/0.8, 320/0.8} {
  \draw[fill=poppurple!60, draw=poppurple] (\angle:\rad) circle (0.08);
}
% Lumen with sperm
\foreach \angle/\rad in {0/0.3, 45/0.3, 90/0.3, 135/0.3, 180/0.3, 225/0.3, 270/0.3, 315/0.3} {
  \draw[popdark] (\angle:\rad) -- ++(\angle:0.2);
}
% Interstitial tissue
\draw[fill=popgold!40, draw=popgold, thick] (3.5,1.5) ellipse (0.6 and 0.4);
\node[popdark, font=\small] at (3.5,1.5) {Leydig cells};
\draw[->, popmuted] (3.0,1.5) -- (2.5,1.2);
% Labels
\node[popdark] at (0,-2.8) {Lumen};
\node[popdark] at (0,2.3) {Basal lamina};
\node[popdark, align=center] at (-3.5,0) {Spermatogonia\\(outer)};
\node[popdark, align=center] at (-3.5,-1) {Primary/secondary\\spermatocytes};
\node[popdark, align=center] at (-3.5,-2) {Spermatids\\(inner)};
\node[popdark, align=center] at (3.5,-1) {Interstitial tissue\\with Leydig cells};
\end{tikzpicture}
\end{center}
\legende{Transverse section of a seminiferous tubule showing the arrangement of germ cells and Sertoli cells.}
\end{popfigure}

\courssub{Spermatogenesis: The Production of Spermatozoa}

\begin{definition}[Spermatogenesis]
\cle{Spermatogenesis} is the production of spermatozoa inside the seminiferous tubules of the testes, from the germinal epithelium. It begins at \textbf{puberty} (around \SI{12}{--}\SI{14}{years} in Cameroon) and continues \textbf{continuously throughout adult life}, unlike oogenesis in the female, which is cyclic and stops at the menopause. The entire process takes approximately \SI{64}{--}\SI{74}{days} in humans.
\end{definition}

The process combines \textbf{growth}, \textbf{meiosis} (which halves the chromosome number from diploid $2n$ to haploid $n$) and \textbf{differentiation (spermiogenesis)}. In humans $2n = 46$ and $n = 23$.

\begin{methode}[Following spermatogenesis step by step]
\begin{enumerate}
\item \textbf{Multiplication (mitotic phase):} diploid \textbf{spermatogonia} ($2n$) at the tubule edge (near basal lamina) divide repeatedly by \textbf{mitosis}. Some remain as stem cells (type A spermatogonia) to maintain the germ line; others (type B spermatogonia) commit to differentiation.
\item \textbf{Growth phase:} type B spermatogonia enlarge to become \textbf{primary spermatocytes} ($2n$), entering prophase I of meiosis (which lasts \textasciitilde\SI{22}{days}).
\item \textbf{Meiosis I} (reduction division): each primary spermatocyte completes meiosis I to form two haploid \textbf{secondary spermatocytes} ($n$). This division separates homologous chromosomes.
\item \textbf{Meiosis II} (equational division): each secondary spermatocyte rapidly divides (without DNA replication) to form two \textbf{spermatids} ($n$) — four spermatids per primary spermatocyte. This division separates sister chromatids.
\item \textbf{Spermiogenesis} (differentiation): each spermatid, still round and immobile, transforms into a streamlined \textbf{spermatozoon}: the nucleus condenses and elongates, the \textbf{acrosome} forms from the Golgi apparatus over the anterior nucleus, mitochondria gather into a tight helix in the \textbf{middle piece}, the \textbf{flagellum} (tail) grows from the centriole, and excess cytoplasm is shed as \textbf{residual bodies} (phagocytosed by Sertoli cells).
\end{enumerate}
The whole sequence takes roughly \SI{74}{days} and occurs continuously in waves along the tubule, so millions of spermatozoa are released every day.
\end{methode}

\begin{popfigure}
\begin{center}
\begin{tikzpicture}[font=\small, node distance=0.55cm,
box/.style={draw=popblue, fill=popblueL, rounded corners=3pt, thick, minimum height=0.85cm, align=center, text width=3.1cm}]
\node[box, align=center] (a){Spermatogonium\\($2n = 46$)};
\node[box, right=of a, align=center] (b){Primary spermatocyte\\($2n = 46$)};
\node[box, right=of b, align=center] (c){2 Secondary spermatocytes\\($n = 23$)};
\node[box, below=of c, align=center] (d){4 Spermatids\\($n = 23$)};
\node[box, left=of d, align=center] (e){4 Spermatozoa\\($n = 23$)};
\draw[->, very thick, popdark] (a) -- (b) node[midway, above, popmuted]{growth};
\draw[->, very thick, popdark] (b) -- (c) node[midway, above, popmuted]{meiosis I};
\draw[->, very thick, popdark] (c) -- (d) node[midway, right, popmuted]{meiosis II};
\draw[->, very thick, popdark] (d) -- (e) node[midway, above, popmuted]{spermiogenesis};
\draw[->, poporange, thick] (a.north) .. controls +(0,0.7) and +(0,0.7) .. node[midway, above, poporange]{mitosis (self-renewal)} (a.north);
\end{tikzpicture}
\end{center}
\legende{Flow diagram of spermatogenesis. Meiosis halves the chromosome number from $2n$ to $n$; one primary spermatocyte yields four functional spermatozoa.}
\end{popfigure}

\begin{attention}[Spermatogenesis versus oogenesis]
Do not confuse the two processes. In spermatogenesis, \textbf{one} primary spermatocyte gives \textbf{four} functional spermatozoa, and production is \textbf{continuous} from puberty. In oogenesis (next section), one primary oocyte gives only \textbf{one} functional ovum plus polar bodies, and the process is \textbf{cyclic}. Examiners love this comparison.
\end{attention}

\courssub{Structure of the Mature Spermatozoon}

A human spermatozoon is about \SI{60}{\micro\meter} long and has three distinct regions, each adapted to its function of reaching and fertilising the ovum:

\begin{popfigure}
\begin{center}
\begin{tikzpicture}[font=\small, scale=1.5]
% Head
\draw[fill=popblueL, draw=popblue, thick] (0,0) ellipse (0.6 and 0.35);
% Nucleus
\draw[fill=popdark!30, draw=popdark, thick] (0,0) ellipse (0.5 and 0.3);
% Acrosome
\draw[fill=poporange!60, draw=poporange, thick] (-0.5,0) arc (180:0:0.5 and 0.3) -- (-0.5,0) -- cycle;
% Middle piece
\draw[fill=popgreenL, draw=popgreen, thick] (-0.6,0) -- (-1.2,0.1) -- (-1.2,-0.1) -- cycle;
% Mitochondria in middle piece
\foreach \x in {-0.7,-0.85,-1.0} {
  \draw[fill=popdark!50] (\x,0) circle (0.04);
}
% Tail (flagellum)
\draw[very thick, poppurple] (-1.2,0) .. controls (-2.0,0.2) and (-3.0,-0.1) .. (-4.0,0);
% Labels
\node[popdark, align=left] at (1.2,0.4) {Head};
\node[popdark, align=left] at (1.2,0.1) {Nucleus (haploid DNA)};
\node[popdark, align=left] at (1.2,-0.2) {Acrosome (enzymes)};
\node[popdark, align=left] at (1.2,-0.5) {Middle piece};
\node[popdark, align=left] at (1.2,-0.8) {Mitochondria (ATP)};
\node[popdark, align=left] at (1.2,-1.1) {Tail (flagellum)};
\node[popdark, align=left] at (1.2,-1.4) {Propulsion};
% Dimension
\draw[<->, popmuted] (0,-0.5) -- (-4.0,-0.5);
\node[popmuted] at (-2.0,-0.7) {\SI{60}{\micro\meter}};
\end{tikzpicture}
\end{center}
\legende{Structure of a mature human spermatozoon.}
\end{popfigure}

\begin{itemize}
\item the \textbf{head} (flattened, \SI{5}{\micro\meter} long), containing the highly condensed haploid \textbf{nucleus} (paternal genetic material) capped by the \textbf{acrosome}, a modified lysosome derived from the Golgi apparatus containing hydrolytic enzymes (hyaluronidase, acrosin) that digest the corona radiata and zona pellucida of the ovum at fertilisation;
\item the \textbf{neck} (connecting piece) containing the proximal and distal centrioles;
\item the \textbf{middle piece}, packed with a helical sheath of \textbf{mitochondria} which produce the ATP that powers flagellar movement;
\item the \textbf{tail (flagellum)}, the longest part (\SI{50}{\micro\meter}), whose axoneme (9+2 microtubule arrangement) beats in a whip-like motion to propel the cell through the female tract.
\end{itemize}

\begin{methode}[Drawing and labelling a spermatozoon (GCE practical skill)]
\begin{enumerate}
\item Draw an oval head (\SI{5}{\micro\meter}); shade a large nucleus inside it and add a cap-like acrosome over the anterior half.
\item Attach a short neck and a thicker middle piece (\SI{5}{\micro\meter}) and fill it with small dots or a helical line representing mitochondria.
\item Draw a long, thin, tapering tail (\SI{50}{\micro\meter}), at least ten times the head length.
\item Label each region and state its function in one line: nucleus — paternal chromosomes; acrosome — enzymes to penetrate ovum coats; mitochondria — ATP production; tail — propulsion.
\end{enumerate}
\end{methode}

\courssub{Hormonal Control of the Testes}

Why does sperm production start at puberty and not before? Because the testes are switched on by hormones arriving through the blood, in a chain of command involving three levels: the \textbf{hypothalamus} (brain), the \textbf{anterior pituitary gland}, and the \textbf{testes} themselves. This is the \cle{hypothalamus–pituitary–testis axis}.

\begin{propriete}[The hypothalamus–pituitary–testis axis]
\begin{itemize}
\item At puberty the \textbf{hypothalamus} secretes \cle{GnRH} (gonadotrophin-releasing hormone) in \textbf{pulsatile} fashion, which stimulates the \textbf{anterior pituitary}.
\item The anterior pituitary releases two gonadotrophins: \cle{FSH} (follicle-stimulating hormone), which acts on the \textbf{Sertoli cells} and, together with testosterone, \textbf{drives spermatogenesis} (stimulates ABP production, supports spermatogonial differentiation); and \cle{LH} (luteinising hormone, also called \textbf{ICSH} — interstitial cell stimulating hormone — in the male), which stimulates the \textbf{Leydig cells} to secrete \textbf{testosterone}.
\item \textbf{Testosterone} acts locally on Sertoli cells (via ABP) to support spermatogenesis, and systemically to maintain secondary sexual characteristics and libido. It also exerts \textbf{negative feedback} on the hypothalamus (reducing GnRH pulse frequency) and anterior pituitary (reducing LH response to GnRH).
\item The \textbf{Sertoli cells} secrete \textbf{inhibin}, which specifically inhibits \textbf{FSH} secretion from the anterior pituitary (negative feedback on FSH only).
\end{itemize}
(The names of the hormones, their cellular targets, and the logic of the feedback loops are examinable; learn to reproduce the diagram from memory.)
\end{propriete}

\begin{popfigure}
\begin{center}
\begin{tikzpicture}[font=\small,
gland/.style={draw=popdark, fill=popgreenL, rounded corners=3pt, thick, minimum height=0.9cm, align=center, text width=3.2cm}]
\node[gland, align=center] (h) at (0,4.4){Hypothalamus\\secretes GnRH (pulsatile)};
\node[gland, align=center] (p) at (0,2.2){Anterior pituitary\\secretes FSH and LH (ICSH)};
\node[gland, align=center] (t) at (0,0){Testes\\Sertoli cells + Leydig cells};
\draw[->, very thick, popblue] (h) -- (p) node[midway, left, popblue]{stimulates};
\draw[->, very thick, popblue] (p) -- (t) node[midway, left, popblue]{stimulates};
\node[draw=poporange, fill=poporange!15, rounded corners=3pt, thick, align=center, text width=3.4cm] (s) at (5.4,0.9) {Testosterone\\(Leydig cells)};
\node[draw=poppink, fill=poppink!15, rounded corners=3pt, thick, align=center, text width=3.4cm] (i) at (5.4,-0.9) {Inhibin\\(Sertoli cells)};
\draw[->, thick, poporange] (t.east) -- (s.west);
\draw[->, thick, poppink] (t.east) -- (i.west);
\draw[->, thick, poporange] (s.north) .. controls (5.4,4.6) .. node[midway, above, poporange]{inhibits ($-$) GnRH} (h.east);
\draw[->, thick, poporange] (s.west) .. controls (2.5,1.5) .. node[midway, left, poporange]{inhibits ($-$) LH} (p.east);
\draw[->, thick, poppink] (i.east) .. controls (8.2,-0.9) and (8.2,2.2) .. node[midway, right, poppink]{inhibits FSH ($-$)} (p.east);
\end{tikzpicture}
\end{center}
\legende{Hormonal control of the testes. GnRH stimulates FSH and LH release; testosterone and inhibin exert negative feedback (orange/pink arrows) on the hypothalamus and pituitary.}
\end{popfigure}

\textbf{Roles of testosterone.}\par
Testosterone is far more than a sperm-production hormone. It is responsible for:
\begin{itemize}
\item the development of the \textbf{secondary sexual characteristics} at puberty: deepening of the voice, growth of beard and body hair, broadening of shoulders, muscular development, growth of the penis and testes, increased bone density;
\item the maintenance of \textbf{libido} (sexual drive) and erectile function;
\item the \textbf{maintenance of spermatogenesis}, acting together with FSH on the seminiferous tubules (via ABP);
\item anabolic effects on muscle and bone, erythropoiesis stimulation, and metabolic effects.
\end{itemize}

\begin{attention}[Common mistakes in hormonal control]
\begin{itemize}
\item FSH and LH are secreted by the \textbf{anterior pituitary}, \emph{not} by the testes. The testes are the \emph{target}, not the source.
\item \textbf{Testosterone} comes from the \textbf{Leydig (interstitial) cells}, \emph{not} from the Sertoli cells. Sertoli cells nourish sperm and secrete \textbf{inhibin} and ABP.
\item LH in the male is the same hormone as in the female; in the male it is often called \textbf{ICSH} (interstitial cell stimulating hormone) — a name that tells you exactly where it acts.
\item GnRH is secreted in \textbf{pulses}; continuous GnRH (as from an agonist implant) \emph{down-regulates} pituitary receptors and \emph{suppresses} gonadotrophin release — the basis of medical castration.
\end{itemize}
\end{attention}

\begin{exoresolu}[Explaining negative feedback (typical GCE question)]
\textbf{Statement.} A man receives regular injections of testosterone as part of a medical treatment. After several months, his sperm count falls dramatically. Explain this observation using your knowledge of the hormonal control of the testes.
\tcblower
\textbf{Solution.}
\begin{enumerate}
\item The injected testosterone raises the blood testosterone level far above normal.
\item By \textbf{negative feedback}, the high testosterone level inhibits the \textbf{hypothalamus} (reduced GnRH pulse frequency) and the \textbf{anterior pituitary} (reduced sensitivity to GnRH).
\item The pituitary therefore secretes less \textbf{LH} and less \textbf{FSH}.
\item FSH, together with locally high intratesticular testosterone (maintained by LH), is required to drive \textbf{spermatogenesis} in the seminiferous tubules; with FSH and LH lacking, sperm production slows or stops, so the sperm count falls.
\end{enumerate}
This is exactly the principle explored for a possible male contraceptive: suppressing FSH/LH by external hormones switches off sperm production reversibly.
\end{exoresolu}

\begin{savaistu}[Why so many sperm?]
A single ejaculate contains \SI{200}{--}\SI{500}{million} spermatozoa, yet only one can fertilise the ovum. Most die on the way: the vagina is acidic (pH \textasciitilde 4), many swim in the wrong direction (only \textasciitilde\SI{1}{\%} enter the cervix), and the journey to the oviduct is enormous relative to their size — comparable to a human swimming \SI{100}{km}. Producing vast numbers is the male strategy that maximises the chance that at least one spermatozoon reaches the egg. In Cameroon, as elsewhere, a sperm count below \SI{15}{million} per ml (oligospermia) is considered subfertile.
\end{savaistu}

\begin{exercice}[Check your understanding]
\begin{enumerate}
\item List, in order, the structures through which a spermatozoon passes from its production site to the outside of the body.
\item State precisely which cells secrete (a) testosterone, (b) inhibin, (c) FSH, (d) GnRH.
\item Explain why the testes are located outside the abdominal cavity.
\item A tumour destroys a man's anterior pituitary. Predict the effect on his testosterone level and sperm production, and justify your answer.
\item Describe the role of Sertoli cells in spermatogenesis.
\item What is the significance of the blood-testis barrier?
\end{enumerate}
\end{exercice}

\begin{corrige}[Exercise 1]
\begin{enumerate}
\item Seminiferous tubules $\rightarrow$ rete testis $\rightarrow$ epididymis $\rightarrow$ vas deferens $\rightarrow$ ejaculatory duct $\rightarrow$ urethra $\rightarrow$ outside (via the penis).
\item (a) Leydig (interstitial) cells; (b) Sertoli cells; (c) anterior pituitary gonadotrophs; (d) hypothalamic neurosecretory cells.
\item Spermatogenesis requires a temperature \SI{2}{\ensuremath{{}^{\circ}\mathrm{C}}}--\SI{3}{\ensuremath{{}^{\circ}\mathrm{C}}} below core body temperature; the scrotum keeps the testes cooler than the abdomen, and adjusts their position via dartos and cremaster muscles to regulate temperature. Cryptorchidism (undescended testes) leads to infertility and increased cancer risk.
\item Without the anterior pituitary there is no LH and no FSH. Without LH, the Leydig cells are not stimulated, so testosterone secretion collapses; without FSH (and testosterone), spermatogenesis stops. The man becomes sterile and loses testosterone-dependent characteristics unless given hormone replacement therapy.
\item Sertoli cells: nourish developing sperm (provide nutrients, growth factors); phagocytose residual bodies; form blood-testis barrier (tight junctions) isolating adluminal compartment from immune system; secrete inhibin (FSH feedback), ABP (concentrates testosterone), and anti-Müllerian hormone (in fetal life).
\item The blood-testis barrier prevents immune cells from recognising haploid sperm antigens as foreign (since they arise after immune tolerance is established), and creates a specialised microenvironment for meiosis and spermiogenesis.
\end{enumerate}
\end{corrige}

\begin{aretenir}[Key points of this section]
\begin{itemize}
\item The testes produce spermatozoa (seminiferous tubules) and testosterone (Leydig cells); the scrotum keeps them \SI{2}{\ensuremath{{}^{\circ}\mathrm{C}}}--\SI{3}{\ensuremath{{}^{\circ}\mathrm{C}}} below body temperature.
\item Spermatogenesis: spermatogonia ($2n$) $\xrightarrow{\text{mitosis}}$ primary spermatocytes ($2n$) $\xrightarrow{\text{meiosis I}}$ secondary spermatocytes ($n$) $\xrightarrow{\text{meiosis II}}$ spermatids ($n$) $\xrightarrow{\text{spermiogenesis}}$ spermatozoa ($n$); continuous from puberty; \SI{64}{--}\SI{74}{days} per cycle; four sperm per primary spermatocyte.
\item A spermatozoon = head (nucleus + acrosome) + neck (centrioles) + middle piece (mitochondria) + tail (flagellum, 9+2 microtubules).
\item Control: hypothalamic GnRH (pulsatile) $\rightarrow$ pituitary FSH (Sertoli cells, spermatogenesis support) and LH/ICSH (Leydig cells, testosterone); testosterone and inhibin exert negative feedback on hypothalamus/pituitary.
\end{itemize}
\end{aretenir}

Having seen how the male gametes are produced and controlled, we now turn to the female partner: the structure of the female reproductive system, the production of ova, and the remarkable monthly cycle that prepares the body for a possible pregnancy.

\coursec{The Female Reproductive System, Oogenesis and the Menstrual Cycle}

In the previous section we studied the male reproductive system and saw that spermatogenesis is a \emph{continuous} process producing millions of small, motile gametes every day from puberty until old age. The female system tells a very different story. A baby girl is born with all the potential egg cells she will ever have, already paused part-way through meiosis; from puberty onwards, roughly one of them matures and is released each month in a beautifully orchestrated hormonal cycle. Understanding this cycle --- the interplay between the brain, the pituitary gland and the ovaries --- is one of the most frequently examined topics at GCE Advanced Level. We shall proceed from anatomy, to gamete production, to the ovarian and uterine cycles and their hormonal control, ending with the hormonal changes of pregnancy.

\courssub{Organs of the Female Reproductive System}

\textbf{Observation.} During a caesarean delivery in a Cameroonian hospital, the surgeon works through the abdomen to reach the uterus, a muscular organ normally the size of a pear. That such a small organ can shelter a baby of \SI{3}{kg} or more tells us immediately that the female reproductive organs are specialised, elastic and under powerful hormonal control.

\begin{definition}[The female reproductive system]
The \cle{female reproductive system} is the set of organs that produce female gametes (ova), receive sperm, provide the site of fertilisation, and nourish and protect the developing embryo and foetus until birth.
\end{definition}

\begin{popfigure}
\begin{center}
\begin{tikzpicture}[scale=1.0, every node/.style={font=\small}]
% Uterus (body)
\draw[fill=poppink!35, draw=popink, thick] (0,3.2) .. controls (-1.3,3.2) and (-1.5,2.2) .. (-1.1,1.2)
   .. controls (-0.9,0.6) and (-0.5,0.4) .. (0,0.4)
   .. controls (0.5,0.4) and (0.9,0.6) .. (1.1,1.2)
   .. controls (1.5,2.2) and (1.3,3.2) .. (0,3.2) -- cycle;
% Endometrium
\draw[fill=poppink!70, draw=popink] (0,2.9) .. controls (-0.85,2.9) and (-1.0,2.1) .. (-0.75,1.3)
   .. controls (-0.6,0.85) and (-0.3,0.7) .. (0,0.7)
   .. controls (0.3,0.7) and (0.6,0.85) .. (0.75,1.3)
   .. controls (1.0,2.1) and (0.85,2.9) .. (0,2.9) -- cycle;
% Cervix
\draw[fill=poppink!50, draw=popink, thick] (-0.35,0.5) rectangle (0.35,-0.35);
% Vagina
\draw[fill=popblueL, draw=popink, thick] (-0.5,-0.35) -- (-0.7,-1.9) -- (0.7,-1.9) -- (0.5,-0.35) -- cycle;
% Fallopian tubes
\draw[draw=popink, thick] (-1.0,2.6) .. controls (-2.2,3.4) and (-3.2,3.2) .. (-3.9,2.5);
\draw[draw=popink, thick] (1.0,2.6) .. controls (2.2,3.4) and (3.2,3.2) .. (3.9,2.5);
% Fimbriae
\foreach \a in {2.3,2.5,2.7}{\draw[draw=popink] (-3.9,\a) -- (-4.25,\a+0.12);}
\foreach \a in {2.3,2.5,2.7}{\draw[draw=popink] (3.9,\a) -- (4.25,\a+0.12);}
% Ovaries
\draw[fill=popgold!60, draw=popink, thick] (-4.35,2.0) ellipse (0.45 and 0.65);
\draw[fill=popgold!60, draw=popink, thick] (4.35,2.0) ellipse (0.45 and 0.65);
% Follicles in right ovary
\fill[white, draw=popink] (4.25,2.25) circle (0.12);
\fill[white, draw=popink] (4.5,1.8) circle (0.15);
\fill[white, draw=popink] (4.2,1.7) circle (0.09);
% Labels
\node[align=center] (ov) at (-5.6,2.0) {\textbf{Ovary}};
\draw[->, popmuted] (ov) -- (-4.8,2.0);
\node[align=center] (ft) at (-4.6,3.6) {\textbf{Fallopian tube}\\ (oviduct)};
\draw[->, popmuted] (ft) -- (-2.6,3.1);
\node[align=center] (fim) at (5.6,3.4) {\textbf{Fimbriae}};
\draw[->, popmuted] (fim) -- (4.25,2.75);
\node[align=center] (ut) at (2.9,1.9) {\textbf{Uterus}\\ (wall: myometrium)};
\draw[->, popmuted] (ut) -- (1.15,1.6);
\node[align=center] (en) at (-2.9,1.1) {\textbf{Endometrium}\\ (lining)};
\draw[->, popmuted] (en) -- (-0.8,1.35);
\node[align=center] (cx) at (-2.4,-0.1) {\textbf{Cervix}};
\draw[->, popmuted] (cx) -- (-0.37,0.05);
\node[align=center] (vg) at (2.4,-1.2) {\textbf{Vagina}};
\draw[->, popmuted] (vg) -- (0.62,-1.2);
\end{tikzpicture}
\end{center}
\legende{Front view of the female reproductive system. The \textbf{vulva} (external genitalia: labia majora, labia minora, clitoris and the vaginal opening) is not shown in this internal view.}
\end{popfigure}

\textbf{Structure and functions of each organ.}

\begin{itemize}
\item \textbf{Ovaries} (one pair, almond-shaped, about \SI{3}{cm} long): the female gonads. They contain thousands of follicles and have two functions: \emph{production of oocytes} (oogenesis) and \emph{secretion of the hormones oestrogen and progesterone}.
\item \textbf{Fallopian tubes (oviducts)} (one pair, about \SI{10}{cm}): muscular tubes lined with ciliated cells. Their funnel-shaped ends bear finger-like projections, the \cle{fimbriae}, which sweep over the ovary at ovulation to capture the released oocyte. \emph{Fertilisation normally occurs in the upper third of the Fallopian tube.}
\item \textbf{Uterus}: a hollow, thick-walled muscular organ. Its wall has two layers of interest: the \cle{myometrium} (smooth muscle which contracts during birth) and the \cle{endometrium} (a glandular lining which thickens each month and into which the embryo implants).
\item \textbf{Cervix}: the narrow, muscular neck of the uterus. It produces mucus whose consistency varies with the cycle (thin and watery at ovulation to allow sperm passage, thick at other times).
\item \textbf{Vagina}: a muscular, elastic tube that receives the penis and semen during intercourse, serves as the passage for menstrual flow, and forms the birth canal.
\item \textbf{Vulva}: the external genital organs (labia majora and minora, clitoris), which protect the vaginal opening and play a role in sexual sensation.
\end{itemize}

\begin{attention}[Common mistake]
Do not confuse the \textbf{urethra} (part of the urinary system) with the vagina. In females, unlike males, the urinary and reproductive openings are \emph{separate}. Also note: fertilisation occurs in the \textbf{Fallopian tube}, not in the uterus; the uterus is the site of \emph{implantation} and development.
\end{attention}

\courssub{Oogenesis: Production of the Female Gamete}

\textbf{Observation.} A woman releases roughly 400 oocytes in her whole reproductive life, while a man produces hundreds of millions of sperm \emph{per day}. This enormous difference in numbers is matched by a difference in strategy: the female gamete is huge, packed with cytoplasm and food reserves, and produced in a slow, discontinuous, precisely timed process.

\begin{definition}[Oogenesis]
\cle{Oogenesis} is the process by which diploid germ cells in the ovary (oogonia) develop into a haploid female gamete. It begins before birth, is arrested for years, and is completed only if fertilisation occurs.
\end{definition}

\textbf{The stages of oogenesis, step by step.}
\begin{enumerate}
\item \textbf{Before birth (foetal life).} Oogonia multiply by mitosis, then enlarge to become \cle{primary oocytes}. Each primary oocyte \emph{begins meiosis I} but is arrested in prophase I. At birth, the ovary already contains its full stock of primary oocytes, each surrounded by a layer of cells forming a \textbf{primordial follicle}.
\item \textbf{From puberty, each month.} Under FSH stimulation, several follicles start to grow, but usually only one becomes a mature \cle{Graafian follicle}. Inside it, the primary oocyte completes \textbf{meiosis I} just before ovulation, producing a large \cle{secondary oocyte} and a tiny \cle{first polar body} (the cytoplasm is divided very unequally so that the future gamete keeps nearly all the reserves).
\item \textbf{Ovulation.} The cell released is the \textbf{secondary oocyte}, which has begun meiosis II but is \emph{arrested in metaphase II}.
\item \textbf{Only if fertilisation occurs.} The entry of a sperm triggers completion of \textbf{meiosis II}, producing the mature \cle{ovum} (whose nucleus fuses with the sperm nucleus) and a \cle{second polar body}. The polar bodies degenerate.
\end{enumerate}

\begin{attention}[The classic GCE trap]
The cell released at ovulation is a \textbf{secondary oocyte arrested in metaphase II}, \emph{not} a fully mature ovum. Meiosis II is only completed \emph{after} a sperm penetrates it. Writing ``a mature ovum is released at ovulation'' costs marks.
\end{attention}

\begin{popfigure}
\begin{center}
\begin{tikzpicture}[font=\small, node distance=0.35cm,
  stage/.style={draw=popink, fill=popblueL, rounded corners, align=center, text width=2.9cm, minimum height=0.9cm},
  stagep/.style={draw=popink, fill=poppink!40, rounded corners, align=center, text width=2.9cm, minimum height=0.9cm},
  arr/.style={->, thick, popdark}]
% Spermatogenesis column
\node[stage, align=center] (s1) at (0,0){\textbf{Spermatogonium}\\ (2n)};
\node[stage, below=of s1, align=center] (s2){\textbf{Primary spermatocyte}\\ (2n)};
\node[stage, below=of s2] (s3) {2 \textbf{secondary spermatocytes} (n)};
\node[stage, below=of s3] (s4) {4 \textbf{spermatids} (n)};
\node[stage, below=of s4] (s5) {4 \textbf{spermatozoa}};
\draw[arr] (s1) -- (s2) node[midway, right, font=\scriptsize]{growth};
\draw[arr] (s2) -- (s3) node[midway, right, font=\scriptsize]{meiosis I};
\draw[arr] (s3) -- (s4) node[midway, right, font=\scriptsize]{meiosis II};
\draw[arr] (s4) -- (s5) node[midway, right, font=\scriptsize]{differentiation};
\node[above=0.15cm of s1, font=\small\bfseries, popblue] {SPERMATOGENESIS};
% Oogenesis column
\node[stagep, align=center] (o1) at (7,0){\textbf{Oogonium} (2n)\\ mitosis before birth};
\node[stagep, below=of o1, align=center] (o2){\textbf{Primary oocyte} (2n)\\ arrested in prophase I};
\node[stagep, below=of o2] (o3) {\textbf{Secondary oocyte} (n) + first polar body};
\node[stagep, below=of o3, align=center] (o4){\textbf{Ovum} (n) + polar bodies\\ (only if fertilised)};
\draw[arr] (o1) -- (o2) node[midway, right, font=\scriptsize]{before birth};
\draw[arr] (o2) -- (o3) node[midway, right, font=\scriptsize]{meiosis I at ovulation};
\draw[arr] (o3) -- (o4) node[midway, right, font=\scriptsize]{meiosis II after fertilisation};
\node[above=0.15cm of o1, font=\small\bfseries, poppink] {OOGENESIS};
\end{tikzpicture}
\end{center}
\legende{Comparison of spermatogenesis and oogenesis. Note the unequal cytoplasmic divisions in oogenesis: one large ovum plus small polar bodies, versus four equal sperm.}
\end{popfigure}

\begin{center}
\begin{tabularx}{\linewidth}{|l|X|X|}
\hline
\rowcolor{popblueL} \textbf{Feature} & \textbf{Spermatogenesis} & \textbf{Oogenesis} \\
\hline
Timing & Continuous from puberty to old age & Begins before birth; one oocyte per month from puberty to menopause \\
\hline
Meiosis & Completed without interruption & Arrested twice (prophase I, then metaphase II); completed only after fertilisation \\
\hline
Gametes per meiosis & 4 functional spermatozoa & 1 functional ovum + 2--3 polar bodies \\
\hline
Cytoplasm division & Equal & Very unequal (ovum keeps the reserves) \\
\hline
Gamete size/motility & Small, motile & Large, non-motile \\
\hline
\end{tabularx}
\end{center}

\courssub{The Ovarian and Uterine (Menstrual) Cycles}

\textbf{Observation.} Most women of reproductive age experience menstrual bleeding roughly every 28 days. This visible event is only the tip of the iceberg: it is the external sign of two coordinated cycles --- one in the \emph{ovary} (preparing and releasing an oocyte) and one in the \emph{uterus} (preparing a bed for a possible embryo).

\begin{definition}[Menstruation]
\cle{Menstruation} is the monthly shedding of the thickened endometrium (with blood and mucus) through the vagina, occurring when fertilisation has not taken place. It marks day 1 of a new cycle.
\end{definition}

\textbf{The ovarian cycle} (average 28 days) has three key events:
\begin{enumerate}
\item \textbf{Follicular phase (days 1--13).} Under FSH, several follicles grow; one becomes dominant and matures into a \cle{Graafian follicle}, which secretes increasing amounts of \textbf{oestrogen}.
\item \textbf{Ovulation (about day 14).} The Graafian follicle ruptures and releases the secondary oocyte.
\item \textbf{Luteal phase (days 15--28).} The remains of the follicle transform into the \cle{corpus luteum} (``yellow body''), which secretes \textbf{progesterone} (and some oestrogen). If there is no pregnancy, the corpus luteum degenerates after about 10--12 days.
\end{enumerate}

\textbf{The uterine cycle} follows in parallel:
\begin{enumerate}
\item \textbf{Menstruation (days 1--5):} the old endometrium breaks down and is shed.
\item \textbf{Proliferative (repair) phase:} under \emph{oestrogen} from the growing follicles, the endometrium is rebuilt and its glands and blood vessels develop.
\item \textbf{Secretory phase:} under \emph{progesterone} from the corpus luteum, the endometrium thickens further, becomes spongy and richly vascularised --- ready for implantation.
\item \textbf{Breakdown:} if no embryo implants, the corpus luteum degenerates, progesterone falls, and the endometrium breaks down: a new menstruation begins.
\end{enumerate}

\begin{definition}[Ovulation]
\cle{Ovulation} is the release of a secondary oocyte from a mature Graafian follicle at the surface of the ovary, triggered by the mid-cycle surge of LH. It occurs at about day 14 of a 28-day cycle.
\end{definition}

\courssub{Hormonal Control of the Cycle}

The cycle is directed by a chain of command: \textbf{hypothalamus} $\rightarrow$ \textbf{pituitary} $\rightarrow$ \textbf{ovary} $\rightarrow$ \textbf{uterus}, with feedback loops returning information at every level.

\begin{enumerate}
\item The hypothalamus secretes \cle{GnRH} (gonadotrophin-releasing hormone), which stimulates the anterior pituitary.
\item The pituitary releases \cle{FSH}, which stimulates follicle growth. The growing follicles secrete \textbf{oestrogen}.
\item At \emph{low/moderate} concentration, oestrogen \emph{inhibits} FSH and LH (negative feedback), so only the dominant follicle survives.
\item When oestrogen reaches a \emph{high, sustained} concentration (mature Graafian follicle), it exerts \textbf{positive feedback}: a massive \cle{LH surge} (with a smaller FSH peak) occurs, triggering \textbf{ovulation} about 24--36 hours later.
\item LH converts the ruptured follicle into the \textbf{corpus luteum}, which secretes \textbf{progesterone} (plus oestrogen). These hormones \emph{inhibit} FSH and LH (negative feedback), preventing new follicles from maturing.
\item Without pregnancy, the corpus luteum degenerates; progesterone and oestrogen fall; the endometrium breaks down (menstruation), and the removal of inhibition allows FSH to rise again --- a new cycle begins.
\end{enumerate}

\begin{propriete}[Positive feedback of oestrogen --- the key exception]
At \cle{high, sustained concentration}, oestrogen switches from negative to \textbf{positive feedback} on the pituitary (and hypothalamus), provoking the \textbf{LH surge} that causes ovulation. This is the classic exception to the rule that reproductive hormones exert negative feedback, and examiners love to test it. (No proof is required; you must be able to \emph{explain} it and locate it on a graph.)
\end{propriete}

\begin{popfigure}
\begin{center}
\begin{tikzpicture}
\begin{axis}[
  width=14.5cm, height=8.2cm,
  xlabel={Day of cycle}, ylabel={Relative hormone level},
  xmin=0, xmax=28, ymin=0, ymax=10.5,
  xtick={0,2,4,6,8,10,12,14,16,18,20,22,24,26,28},
  ytick=\empty,
  legend style={at={(0.02,0.98)}, anchor=north west, font=\scriptsize, draw=popline},
  grid=major, grid style={popline!40},
  axis line style={popmuted},
]
% FSH
\addplot[very thick, popblue, smooth] coordinates {(0,3.2) (2,3.6) (4,3.2) (6,2.6) (8,2.2) (10,2.4) (12,3.2) (13.2,5.2) (14,3.4) (16,1.8) (18,1.5) (20,1.4) (22,1.4) (24,1.6) (26,2.2) (28,3.0)};
\addlegendentry{FSH}
% LH
\addplot[very thick, poporange, smooth] coordinates {(0,1.8) (2,2.0) (4,1.8) (6,1.6) (8,1.7) (10,2.2) (12,4.0) (13.2,9.6) (14,5.0) (15,2.4) (16,1.6) (18,1.3) (20,1.2) (22,1.2) (24,1.3) (26,1.5) (28,1.7)};
\addlegendentry{LH}
% Oestrogen
\addplot[very thick, popgreen, smooth] coordinates {(0,1.0) (2,1.2) (4,1.8) (6,2.6) (8,3.6) (10,5.2) (12,7.6) (13,8.4) (14,5.4) (15,3.6) (16,3.4) (18,3.8) (20,4.0) (22,3.8) (24,3.2) (26,1.8) (28,1.0)};
\addlegendentry{Oestrogen}
% Progesterone
\addplot[very thick, poppurple, smooth] coordinates {(0,0.6) (2,0.6) (4,0.6) (6,0.7) (8,0.7) (10,0.8) (12,1.0) (14,1.6) (16,3.4) (18,5.2) (20,6.2) (22,5.8) (24,4.4) (26,2.2) (28,0.7)};
\addlegendentry{Progesterone}
% Ovulation marker
\draw[dashed, thick, popink] (axis cs:14,0) -- (axis cs:14,10.2);
\node[font=\scriptsize\bfseries, popink, align=center] at (axis cs:14,10.3) {ovulation};
% Phase annotations
\node[font=\scriptsize, align=center, popmuted] at (axis cs:6.5,9.7) {follicular phase\\ endometrium proliferates};
\node[font=\scriptsize, align=center, popmuted] at (axis cs:21,9.7) {luteal phase\\ endometrium thickens};
\draw[<->, popmuted] (axis cs:0.3,8.9) -- (axis cs:4.7,8.9) node[midway, above, font=\scriptsize]{menses};
\end{axis}
\end{tikzpicture}
\end{center}
\legende{Hormone levels across a 28-day cycle. The oestrogen peak (green) immediately precedes and triggers the LH surge (orange), which causes ovulation (dashed line). Progesterone (purple) dominates the luteal phase; its fall at day 26--28 causes menstruation.}
\end{popfigure}

\begin{methode}[Identifying the day of ovulation on a hormone graph]
\begin{enumerate}
\item Locate the \textbf{sharpest, tallest peak}: this is the \textbf{LH surge}.
\item Check that it is immediately preceded by an \textbf{oestrogen peak} (positive feedback) and followed by a \textbf{rise in progesterone} (corpus luteum forming).
\item Ovulation occurs at, or about 24--36 hours after, the LH peak --- typically \textbf{day 14} of a 28-day cycle.
\item Justify in words: ``the LH surge, triggered by high oestrogen, causes the Graafian follicle to rupture.''
\end{enumerate}
\end{methode}

\begin{exoresolu}[Interpreting a cycle graph]
On a graph of a 28-day cycle, LH rises sharply from day 12 to a maximum at day 13, then falls; progesterone begins to rise from day 15. (a) On which day does ovulation most probably occur? (b) Explain the hormonal cause of the LH peak. (c) Predict the consequence for the endometrium if progesterone falls sharply at day 26.
\tcblower
\textbf{(a)} Ovulation occurs about 24--36 hours after the LH peak, so most probably on \textbf{day 14}.
\textbf{(b)} The mature Graafian follicle has been secreting increasing amounts of oestrogen. When oestrogen reaches a high, sustained level, it exerts \emph{positive feedback} on the pituitary (via the hypothalamus), provoking the sudden \textbf{LH surge}.
\textbf{(c)} Progesterone maintains the thickened endometrium. A sharp fall of progesterone (due to degeneration of the corpus luteum in the absence of pregnancy) removes this support, so the endometrium \textbf{breaks down and is shed}: menstruation begins, marking day 1 of the next cycle.
\end{exoresolu}

\courssub{Hormonal Changes During Pregnancy}

\textbf{Observation.} If the cycle is a monthly rehearsal for pregnancy, what happens when the rehearsal succeeds? A woman who conceives does not menstruate --- and the reason is entirely hormonal.

If fertilisation occurs, the young embryo implants in the endometrium around day 20--22 and its outer layer (the future placenta) secretes \cle{hCG} (human chorionic gonadotrophin). hCG acts like LH: it \textbf{maintains the corpus luteum}, which therefore continues to secrete \textbf{progesterone and oestrogen}. These hormones keep the endometrium intact (no menstruation) and keep FSH and LH inhibited (no new ovulation). From about the third month, the \textbf{placenta takes over} the production of progesterone and oestrogen, and the corpus luteum is no longer needed. Progesterone thus remains high throughout pregnancy --- this is the hormone detected indirectly by the absence of menstruation, and hCG itself is the hormone detected by pregnancy test kits.

\begin{attention}[GCE exam tip]
A very frequent question: \emph{``Explain why progesterone levels remain high during pregnancy.''} The full-mark answer has two stages: (1) the embryo secretes \textbf{hCG}, which maintains the \textbf{corpus luteum} so it keeps producing progesterone; (2) later, the \textbf{placenta} takes over progesterone (and oestrogen) production. Mentioning only one stage loses marks.
\end{attention}

\begin{aretenir}[Key points of this section]
\begin{itemize}
\item The ovary produces oocytes and secretes oestrogen and progesterone; fertilisation occurs in the Fallopian tube; implantation occurs in the endometrium of the uterus.
\item Oogenesis is discontinuous: primary oocytes are formed before birth; meiosis I is completed at ovulation (secondary oocyte + first polar body); meiosis II is completed only after fertilisation. One ovum per cycle.
\item The cell released at ovulation is a \textbf{secondary oocyte arrested in metaphase II}.
\item Cycle logic: FSH $\rightarrow$ follicle growth $\rightarrow$ oestrogen $\rightarrow$ (positive feedback) LH surge $\rightarrow$ ovulation $\rightarrow$ corpus luteum $\rightarrow$ progesterone $\rightarrow$ negative feedback; fall of progesterone $\rightarrow$ menstruation.
\item In pregnancy, hCG maintains the corpus luteum, then the placenta takes over hormone production.
\end{itemize}
\end{aretenir}

\begin{exercice}[Check your understanding]
\begin{enumerate}
\item Name the organ where fertilisation normally occurs and the organ where implantation occurs.
\item State precisely the stage of meiosis of the cell released at ovulation, and explain when meiosis is completed.
\item Explain why only one functional ovum results from oogenesis while spermatogenesis yields four spermatozoa.
\item On a hormone graph, which peak signals that ovulation is imminent, and what causes it?
\item A woman's progesterone level remains high three weeks after ovulation. Suggest a biological explanation.
\end{enumerate}
\end{exercice}

\begin{corrige}[Exercise 1]
\begin{enumerate}
\item Fertilisation: the \textbf{Fallopian tube} (oviduct). Implantation: the \textbf{endometrium of the uterus}.
\item A \textbf{secondary oocyte arrested in metaphase II}. Meiosis II is completed only if a sperm penetrates the oocyte (after fertilisation), producing the ovum and the second polar body.
\item In oogenesis the cytoplasm divides \emph{unequally}: almost all cytoplasm and reserves go to one cell (the ovum), the rest forming tiny polar bodies that degenerate. In spermatogenesis divisions are equal, giving four functional sperm.
\item The \textbf{LH surge}. It is caused by \emph{positive feedback} of the high, sustained oestrogen secreted by the mature Graafian follicle.
\item She is probably pregnant: the implanted embryo secretes \textbf{hCG}, which maintains the corpus luteum, so progesterone secretion continues instead of falling.
\end{enumerate}
\end{corrige}

\begin{savaistu}[Did you know?]
The hormone \textbf{hCG} is the basis of the pregnancy test kits sold in pharmacies across Cameroon: it is detectable in urine from about two weeks after fertilisation. Also, the ``fertile window'' used in natural family planning corresponds to the few days around ovulation, because the secondary oocyte survives only about 24 hours while sperm can survive up to about 3--5 days in the female tract.
\end{savaistu}

We are now equipped to follow the gametes to their meeting point: in the next section we shall study fertilisation itself, the development of the embryo, and the process of birth.

\coursec{Fertilization, Embryonic Development and Birth}

In the previous section we followed the female gamete: we saw how oogenesis produces a secondary oocyte, how ovulation releases it into the Fallopian tube, and how the menstrual cycle prepares the endometrium to receive a possible embryo. All of this elaborate preparation has a single purpose: the meeting of one sperm and one oocyte. In this section we follow that meeting and its extraordinary consequences — fertilization, the nine months of development that transform a single cell into a baby, the drama of birth, and the production of milk that nourishes the newborn. This is one of the most frequently examined topics in the GCE Advanced Level paper, so master every stage.

\courssub{From Copulation to the Meeting of Gametes}

\textbf{Copulation and ejaculation.} During sexual intercourse (copulation), the erect penis is inserted into the vagina. Rhythmic stimulation triggers the \cle{ejaculation} reflex: coordinated contractions of the epididymis, vas deferens, seminal vesicles and prostate expel semen into the urethra and out into the vagina, near the cervix. A single ejaculate (about $3$–$4\ \mathrm{cm^3}$) contains roughly $300$ million spermatozoa, yet only a few hundred will ever reach the egg. Why this enormous wastage? The female tract is a hostile obstacle course.

\textbf{The journey of sperm through the female tract.} From the vagina, sperm must:

\begin{enumerate}
\item \textbf{Cross the cervix.} Around ovulation, under the influence of oestrogen, the cervical mucus becomes thin, watery and alkaline — channels form that guide sperm through. At other times of the cycle the mucus is thick and almost impassable. Millions of sperm are stopped here.
\item \textbf{Swim through the uterus}, helped by muscular contractions of the uterine wall and by the tail (flagellum) beat powered by mitochondria in the mid-piece.
\item \textbf{Enter a Fallopian tube.} Only the tube on the side that ovulated contains an egg. Sperm reaching the wrong tube are lost.
\item \textbf{Reach the ampulla} (the wider region of the tube), where the secondary oocyte, surrounded by follicle cells (the corona radiata), awaits — fertilization usually occurs here, within about 24 hours of ovulation.
\end{enumerate}

Of the hundreds of millions deposited, fewer than $1\,000$ typically arrive near the egg, and only \emph{one} will fertilize it.

\textbf{Capacitation.} Freshly ejaculated sperm are not yet able to fertilize. During their passage through the female tract (over several hours), secretions of the uterus and tube remove cholesterol and glycoproteins from the sperm surface membrane. This change, called \cle{capacitation}, destabilises the membrane of the acrosome and makes the sperm capable of undergoing the acrosome reaction. You only need to know it briefly: capacitation is the \emph{final maturation of sperm inside the female tract}.

\courssub{Fertilization in the Fallopian Tube}

\begin{definition}[Fertilization]
\cle{Fertilization} is the fusion of a haploid sperm nucleus with a haploid egg nucleus to form a diploid \cle{zygote}. In humans it normally takes place in the ampulla of the Fallopian tube.
\end{definition}

It proceeds in a precise sequence:

\begin{enumerate}
\item \textbf{Penetration of the corona radiata.} The capacitated sperm pushes between the follicle cells surrounding the oocyte, aided by enzymes.
\item \textbf{The acrosome reaction.} On contact with the \cle{zona pellucida} (the glycoprotein coat of the egg), the acrosome — a cap-like vesicle at the sperm head — releases hydrolytic enzymes (e.g. hyaluronidase, acrosin). These digest a path through the zona pellucida.
\item \textbf{Fusion of membranes.} The membrane of the sperm head fuses with the egg membrane, and the sperm nucleus (and usually the whole sperm, minus the tail) enters the cytoplasm.
\item \textbf{The block to polyspermy.} Entry of the first sperm triggers the egg to release cortical granules whose contents alter the zona pellucida, hardening it (\emph{zona reaction}). This \cle{block to polyspermy} prevents any other sperm from entering — essential, because a zygote with three sets of chromosomes cannot develop normally.
\item \textbf{Completion of meiosis and fusion of nuclei.} The secondary oocyte, which was arrested in metaphase II, now completes meiosis II, producing the ovum and a second polar body. The haploid sperm nucleus ($n = 23$) and the haploid egg nucleus ($n = 23$) — now called \emph{pronuclei} — fuse. \textbf{Diploidy is restored}: the zygote has $2n = 46$ chromosomes, half from each parent. Sex is determined at this instant: an X-bearing sperm gives a girl (XX), a Y-bearing sperm a boy (XY).
\end{enumerate}

\begin{aretenir}[Fertilization — the essentials]
Site: ampulla of the Fallopian tube. Key events: acrosome reaction $\rightarrow$ penetration of zona pellucida $\rightarrow$ block to polyspermy $\rightarrow$ fusion of pronuclei $\rightarrow$ restoration of the diploid number ($2n = 46$).
\end{aretenir}

\courssub{Early Development: From Zygote to Implanted Blastocyst}

The zygote does not wait. As it is swept along the Fallopian tube by cilia and muscular contractions (a journey of about 4–5 days), it begins to divide.

\begin{itemize}
\item \textbf{Cleavage} (days 1–3): repeated mitotic divisions \emph{without growth} — the cells (blastomeres) become smaller at each division, so the overall size stays constant. The embryo passes through 2-cell, 4-cell, 8-cell stages.
\item \textbf{Morula} (day 3–4): a solid ball of about 16–32 cells ("mulberry"), still enclosed in the zona pellucida. It enters the uterus.
\item \textbf{Blastocyst} (day 5–6): fluid accumulates inside, forming a hollow ball. Two parts can now be distinguished: the \cle{inner cell mass} (which will form the embryo proper) and the outer \cle{trophoblast} (which will form the chorion and the fetal part of the placenta). The zona pellucida breaks down ("hatching").
\item \textbf{Implantation} (day 6–7 after fertilization): the trophoblast secretes enzymes that digest a small area of endometrium, and the blastocyst burrows into it.
\end{itemize}

\begin{definition}[Implantation]
\cle{Implantation} is the embedding of the blastocyst in the endometrium (the vascularised lining) of the uterus. It occurs about 6–7 days after fertilization and marks the true beginning of pregnancy.
\end{definition}

\begin{popfigure}
\begin{tikzpicture}[font=\small, >=stealth]
% stages along an axis
\draw[->, thick, popmuted] (0,0) -- (13.4,0) node[below left=2pt and 0pt, popmuted]{days after fertilization};
% zygote day 0
\draw[fill=popblueL, draw=popblue] (0.7,1.1) circle (0.45);
\fill[popblue] (0.7,1.1) circle (0.1);
\node[align=center] at (0.7,2.0) {Zygote\\ (day 0)};
% cleavage day 2
\draw[fill=popblueL, draw=popblue] (3.4,1.1) circle (0.45);
\foreach \a in {45,135,225,315}{\fill[popblue] ({3.4+0.2*cos(\a)},{1.1+0.2*sin(\a)}) circle (0.09);}
\node[align=center] at (3.4,2.0) {Cleavage\\ (days 1--3)};
% morula day 4
\draw[fill=popgreenL, draw=popgreen] (6.2,1.1) circle (0.45);
\foreach \a in {0,60,...,330}{\fill[popgreen] ({6.2+0.22*cos(\a)},{1.1+0.22*sin(\a)}) circle (0.06);}
\fill[popgreen] (6.2,1.1) circle (0.06);
\node[align=center] at (6.2,2.0) {Morula\\ (day 3--4)};
% blastocyst day 5
\draw[fill=white, draw=poporange, thick] (9.0,1.1) circle (0.5);
\fill[poporange] (8.75,0.95) circle (0.16);
\node[align=center] at (9.0,2.0) {Blastocyst\\ (day 5)};
% implantation day 6-7
\draw[fill=poppink!20, draw=poppink, thick] (11.6,0.35) rectangle (13.2,1.85);
\draw[fill=white, draw=poporange, thick] (12.4,1.15) circle (0.42);
\fill[poporange] (12.2,1.0) circle (0.14);
\node[align=center] at (12.4,2.35) {Implantation\\ (day 6--7)};
% tube/uterus labels
\node[popmuted, align=center] at (3.4,-0.5) {in Fallopian tube};
\node[popmuted, align=center] at (11.0,-0.5) {in uterus};
\draw[popmuted, dashed] (7.6,-0.25) -- (7.6,2.3);
\end{tikzpicture}
\legende{From zygote to implanted blastocyst. Cleavage divisions occur without growth; implantation takes place about day 6--7.}
\end{popfigure}

\begin{methode}[Sequencing exercise — from ejaculation to implantation]
A classic GCE task: arrange scrambled events in order. Reason as follows.
\begin{enumerate}
\item Identify the starting point: ejaculation into the vagina.
\item Follow the sperm: cervix $\rightarrow$ uterus $\rightarrow$ Fallopian tube (capacitation occurs during this journey).
\item Fertilization events: acrosome reaction $\rightarrow$ penetration of zona pellucida $\rightarrow$ fusion of pronuclei (zygote).
\item Follow the embryo: cleavage $\rightarrow$ morula $\rightarrow$ blastocyst $\rightarrow$ implantation (day 6–7).
\item Check the logic: anything involving the uterus (morula arrival, implantation) must come \emph{after} events in the tube.
\end{enumerate}
\end{methode}

\courssub{Extra-embryonic Membranes and the Placenta}

After implantation, the embryo develops four \cle{extra-embryonic membranes} which protect and sustain it but form no part of the baby's body:

\begin{itemize}
\item \textbf{Amnion}: a fluid-filled sac enclosing the embryo. The amniotic fluid cushions the fetus against shocks, allows movement, and keeps a constant temperature. ("The waters" that break before birth.)
\item \textbf{Chorion}: the outermost membrane, derived from the trophoblast. It develops finger-like \emph{chorionic villi} that grow into the endometrium and form the fetal part of the placenta. It also secretes \cle{hCG} (human chorionic gonadotrophin), the hormone detected by pregnancy tests, which maintains the corpus luteum during the first weeks.
\item \textbf{Allantois}: a small outgrowth whose blood vessels become the blood vessels of the umbilical cord.
\item \textbf{Yolk sac}: in humans it contains no yolk; it produces the first blood cells and the primordial germ cells.
\end{itemize}

\textbf{The placenta.} By about week 10–12, the \cle{placenta} is fully formed from the chorionic villi (fetal contribution) and the endometrium (maternal contribution). The fetus is connected to it by the \cle{umbilical cord}, containing two umbilical arteries (carrying deoxygenated blood and wastes from fetus to placenta) and one umbilical vein (carrying oxygenated, nutrient-rich blood back to the fetus). The placenta has three great functions:

\begin{enumerate}
\item \textbf{Exchange organ}: oxygen and nutrients (glucose, amino acids, fatty acids, vitamins, minerals) diffuse from maternal blood to fetal blood; carbon dioxide, urea and other wastes diffuse in the opposite direction. Some antibodies (IgG) also cross, giving the newborn passive immunity.
\item \textbf{Endocrine gland}: the placenta secretes hCG, then \cle{progesterone} and \cle{oestrogens} (taking over from the corpus luteum after about 3 months), maintaining the endometrium and suppressing further ovulation and menstruation; it also produces human placental lactogen.
\item \textbf{Selective barrier}: it blocks many (but not all) harmful agents — some viruses, the rubella virus, alcohol, nicotine and certain drugs can cross and damage the fetus.
\end{enumerate}

\begin{propriete}[The placental barrier — no mixing of bloods]
Exchange across the placenta occurs \textbf{by diffusion} (and active transport for some substances) through the thin walls of the chorionic villi. Maternal and fetal blood \textbf{do NOT mix}: they remain separated by the villus epithelium, connective tissue and capillary endothelium. This is essential because mother and fetus may have different, incompatible blood groups (e.g. Rhesus). (Statement examinable; no proof required.)
\end{propriete}

\begin{attention}[Common mistake]
Never write "the mother's blood flows into the baby" or "the two bloods mix in the placenta". Examiners penalise this heavily. Correct phrasing: \emph{substances diffuse between two separate circulations}.
\end{attention}

\begin{popfigure}
\begin{tikzpicture}[font=\small]
% uterus wall
\draw[fill=poppink!15, draw=poppink, thick] (0,0) ellipse (4.6 and 3.4);
% amniotic cavity
\draw[fill=popblueL, draw=popblue, thick] (0,0) ellipse (3.6 and 2.6);
\node[popblue] at (-2.3,1.9) {Amnion};
\node[poppink] at (3.4,2.6) {Uterus wall};
% placenta
\draw[fill=popgreenL, draw=popgreen, thick] (2.6,0.4) ellipse (1.1 and 1.5);
\node[popgreen, align=center] at (2.6,0.4) {Placenta};
% fetus
\draw[fill=white, draw=popdark, thick] (-0.9,0.9) circle (0.65); % head
\draw[fill=white, draw=popdark, thick] (-0.3,-0.6) ellipse (1.0 and 0.8); % body
\node[popdark] at (-0.3,-0.6) {Fetus};
% umbilical cord
\draw[very thick, poporange] (0.5,-0.4) .. controls (1.4,-0.2) .. (2.0,0.1);
\node[poporange, align=center] at (1.1,-1.15) {Umbilical cord};
% amniotic fluid label
\node[popblue, align=center] at (-2.4,-1.6) {Amniotic fluid};
\end{tikzpicture}
\legende{The fetus in utero: the amnion encloses the fetus in amniotic fluid; the umbilical cord links the fetus to the placenta attached to the uterus wall.}
\end{popfigure}

\courssub{Gestation: Nine Months of Development}

\cle{Gestation} (pregnancy) lasts about \textbf{38–40 weeks} (roughly 9 months) from fertilization, conventionally divided into three trimesters.

\begin{itemize}
\item \textbf{First trimester (weeks 1–12)}: the period of \cle{organogenesis} — all the major organs are laid down. The heart begins beating around week 4; the neural tube, limbs, eyes and digestive system form. This is the most delicate period: the embryo is highly sensitive to alcohol, drugs and infections (e.g. rubella), which can cause malformations. By week 12 the fetus is recognisably human, about 7–8 cm long.
\item \textbf{Second trimester (weeks 13–26)}: mainly \emph{growth} and maturation. Movements are felt by the mother ("quickening"); the fetus can suck its thumb; hair and nails appear.
\item \textbf{Third trimester (weeks 27–40)}: rapid growth and fat deposition; the lungs mature (surfactant production); the fetus usually turns head-down in preparation for birth.
\end{itemize}

\begin{savaistu}[Pregnancy care in Cameroon]
Because the first trimester is the period of organogenesis, health workers in Cameroonian antenatal clinics insist that expectant mothers avoid alcohol and self-medication, attend consultations early, and sleep under insecticide-treated mosquito nets: malaria infection during pregnancy reduces placental exchange and is a major cause of low birth weight. Protecting the placenta means protecting the baby.
\end{savaistu}

\courssub{Birth (Parturition)}

\begin{definition}[Parturition]
\cle{Parturition} is the process of giving birth: the expulsion of the fetus and then of the placenta from the uterus, driven by powerful uterine contractions under the control of \cle{oxytocin} and \cle{prostaglandins}.
\end{definition}

\textbf{Hormonal trigger.} As pregnancy reaches term, the fetal pituitary and the placenta shift the hormonal balance: oestrogen rises relative to progesterone, making the myometrium (uterine muscle) sensitive to oxytocin. The fetal head pressing on the cervix sends nerve impulses to the mother's hypothalamus, causing the posterior pituitary to release \cle{oxytocin}. Oxytocin stimulates uterine contractions; the uterus and fetal membranes also release \cle{prostaglandins}, which intensify contractions and soften (ripen) the cervix.

\begin{attention}[Oxytocin: a POSITIVE feedback — the classic trap]
Most hormonal controls you have studied (e.g. thyroxine, blood glucose) are \emph{negative} feedback. Oxytocin in labour is the famous exception: contractions push the baby's head against the cervix $\rightarrow$ more stretching $\rightarrow$ \textbf{more oxytocin} released $\rightarrow$ \textbf{stronger contractions} $\rightarrow$ and so on, until the baby is born and the stimulus stops. Do not describe it as negative feedback.
\end{attention}

\begin{popfigure}
\begin{tikzpicture}[font=\small, >=stealth, node distance=0.9cm]
\node[draw=popblue, fill=popblueL, rounded corners, align=center, text width=3.1cm] (head) {Baby's head presses on cervix};
\node[draw=popgreen, fill=popgreenL, rounded corners, align=center, text width=3.1cm, right=1.4cm of head] (pit) {Posterior pituitary releases MORE oxytocin};
\node[draw=poporange, fill=white, rounded corners, align=center, text width=3.1cm, below=1.0cm of pit] (ut) {Stronger uterine contractions (+ prostaglandins)};
\node[draw=poppink, fill=white, rounded corners, align=center, text width=3.1cm, below=1.0cm of head] (stretch) {Cervix stretched further};
\draw[->, thick, popdark] (head) -- (pit);
\draw[->, thick, popdark] (pit) -- (ut);
\draw[->, thick, popdark] (ut) -- (stretch);
\draw[->, thick, popdark] (stretch) -- (head);
\node[popdark, align=center] at (4.6,-2.6) {loop repeats until birth};
\end{tikzpicture}
\legende{Positive feedback of oxytocin during labour: each contraction reinforces the next until delivery breaks the loop.}
\end{popfigure}

\textbf{The three stages of labour:}
\begin{enumerate}
\item \textbf{Dilation of the cervix} (longest stage, up to several hours): regular contractions dilate the cervix to about 10 cm; the amnion usually ruptures ("breaking of the waters").
\item \textbf{Expulsion of the baby}: powerful contractions, aided by the mother's abdominal muscles, push the baby head-first through the cervix and vagina. The umbilical cord is clamped and cut.
\item \textbf{Delivery of the placenta (afterbirth)}: about 10–30 minutes later, further contractions detach and expel the placenta and membranes. The uterus then contracts firmly, limiting bleeding.
\end{enumerate}

\courssub{Lactation: Feeding the Newborn}

During pregnancy, high levels of oestrogen and progesterone develop the mammary glands but \emph{inhibit} milk secretion. After birth, the fall in these hormones lifts the inhibition, and two pituitary hormones take over:

\begin{itemize}
\item \textbf{Prolactin} (anterior pituitary): stimulates the alveoli of the mammary glands to \emph{produce} milk. Suckling maintains prolactin secretion — the more the baby suckles, the more milk is made.
\item \textbf{Oxytocin} (posterior pituitary): released by the suckling reflex, it causes the myoepithelial cells around the alveoli to contract, \emph{ejecting} milk into the ducts (the "let-down" reflex). Even the baby's cry can trigger it.
\end{itemize}

The first secretion, \cle{colostrum}, is a yellowish fluid rich in proteins and maternal \emph{antibodies} but low in fat; it gives the newborn crucial passive immunity during its first days. Breastfeeding is strongly recommended in Cameroon and worldwide: breast milk is sterile, at the right temperature, perfectly balanced nutritionally, transmits antibodies, and is free — a decisive advantage where clean water for formula preparation is not guaranteed.

\begin{exoresolu}[Worked exercise — functions of the placenta]
\textbf{Statement (GCE style).} (a) Define implantation. (b) State four functions of the placenta, explaining briefly how exchange occurs. (c) Explain why a Rhesus-negative mother and a Rhesus-positive fetus can coexist during a first pregnancy.
\tcblower
\textbf{Solution.} (a) Implantation is the embedding of the blastocyst in the endometrium of the uterus, about 6–7 days after fertilization.
(b) Any four of: (i) diffusion of oxygen from maternal to fetal blood; (ii) diffusion of nutrients (glucose, amino acids) to the fetus; (iii) removal of carbon dioxide and urea from fetal blood; (iv) transfer of maternal antibodies (passive immunity); (v) secretion of hormones — hCG, progesterone and oestrogens — which maintain pregnancy. Exchange occurs by diffusion (and active transport for some substances) across the thin walls of the chorionic villi, which present a large surface area bathed in maternal blood.
(c) Maternal and fetal blood do \emph{not} mix in the placenta; they are separated by the villus tissues. During a first pregnancy the mother's immune system therefore does not meet the fetal Rhesus antigens, so no anti-Rhesus antibodies are produced and the fetus is unharmed. (Sensitisation may occur at birth, when bloods can mix — a risk for later pregnancies.)
\end{exoresolu}

\begin{attention}[Exam tip (GCE)]
When asked for the functions of the placenta, always give \textbf{both} categories: the \emph{nutritive/respiratory/excretory} exchanges \emph{and} the \emph{endocrine} role (secretion of progesterone, oestrogens, hCG). Answers giving only "exchange of gases and food" earn only half the marks.
\end{attention}

\begin{exercice}[Check your understanding]
\begin{enumerate}
\item Place in correct order: morula; ejaculation; implantation; acrosome reaction; blastocyst; capacitation; fusion of pronuclei.
\item Explain why the block to polyspermy is essential for normal development.
\item Distinguish clearly between the roles of prolactin and oxytocin in lactation.
\item State the three stages of parturition and the hormone(s) responsible for the contractions.
\end{enumerate}
\end{exercice}

\begin{corrige}[Exercise n°1]
\begin{enumerate}
\item Ejaculation $\rightarrow$ capacitation $\rightarrow$ acrosome reaction $\rightarrow$ fusion of pronuclei $\rightarrow$ morula $\rightarrow$ blastocyst $\rightarrow$ implantation.
\item If several sperm entered, the zygote would receive extra sets of chromosomes (triploidy or more); such abnormal chromosome numbers prevent normal development, so the embryo would die.
\item Prolactin (anterior pituitary) stimulates milk \emph{production} by the alveoli; oxytocin (posterior pituitary) causes contraction of myoepithelial cells and milk \emph{ejection} (let-down) in response to suckling.
\item (i) Dilation of the cervix; (ii) expulsion of the baby; (iii) delivery of the placenta (afterbirth). Contractions are driven by oxytocin (positive feedback) and prostaglandins.
\end{enumerate}
\end{corrige}

\begin{aretenir}[Section summary]
Fertilization occurs in the Fallopian tube: acrosome reaction, block to polyspermy, fusion of pronuclei restoring $2n = 46$. The zygote undergoes cleavage $\rightarrow$ morula $\rightarrow$ blastocyst, which implants in the endometrium around day 6–7. The extra-embryonic membranes (amnion, chorion, allantois, yolk sac) protect and nourish the embryo; the placenta exchanges gases, nutrients and wastes by diffusion — the two bloods never mix — and secretes hCG, progesterone and oestrogens. Gestation lasts 38–40 weeks, with organogenesis in the first trimester. Birth is driven by oxytocin (positive feedback) and prostaglandins through three stages: dilation, expulsion, afterbirth. Lactation depends on prolactin (production) and oxytocin (ejection); colostrum gives the newborn its first antibodies.
\end{aretenir}

\coursec{Birth Control Methods}

Having studied the intricate processes of fertilisation, embryonic development and birth, we now turn to a crucial application of reproductive biology: the deliberate regulation of fertility. Understanding \cle{contraception} is not merely a medical topic; it is a cornerstone of public health, family welfare and socio-economic planning, particularly in our Cameroonian context where maternal health and population dynamics are national priorities.

\courssub{Concept and Rationale for Birth Control}

\begin{definition}[Contraception]
\cle{Contraception} (or birth control) is the deliberate use of artificial or natural methods to prevent \cle{fertilisation} (the fusion of sperm and ovum) or, in some cases, to prevent \cle{implantation} of the blastocyst in the uterine endometrium. It allows individuals and couples to anticipate and attain their desired number of children and the spacing and timing of their births.
\end{definition}

\begin{savaistu}[Why Birth Control Matters in Cameroon]
In Cameroon, the rationale for birth control is multi-dimensional:
\begin{itemize}
    \item \textbf{Maternal and Child Health:} Spacing pregnancies (at least 24 months) reduces the risk of maternal anaemia, haemorrhage, and low birth weight babies. It lowers maternal mortality (MMR) and infant mortality (IMR).
    \item \textbf{Family Welfare:} Planned families can invest more in nutrition, education (school fees, books) and healthcare per child.
    \item \textbf{Socio-economic Development:} A controlled population growth rate eases pressure on infrastructure (schools, hospitals, roads), food security and youth employment.
    \item \textbf{Women's Empowerment:} Control over reproduction enables girls to stay in school longer and women to participate fully in the workforce.
\end{itemize}
The \emph{National Population Policy} of Cameroon explicitly promotes family planning as a right and a development tool.
\end{savaistu}

\courssub{Classification of Contraceptive Methods}

Contraceptive methods are classified by their \cle{mechanism of action} and \cle{permanence}. The following chart summarises the landscape.

\begin{popfigure}
\centering
\begin{tikzpicture}[
    node distance=1.2cm and 1.5cm,
    main/.style={draw=popblue, thick, rounded corners, fill=popblueL, text width=3.2cm, align=center, minimum height=1.2cm},
    sub/.style={draw=popdark, fill=white, rounded corners, text width=2.8cm, align=center, minimum height=1cm, font=\small},
    arrow/.style={-Stealth, thick, popdark},
    labelstyle/.style={font=\footnotesize, text width=3cm, align=left, anchor=west}
]
% Title
\node[font=\bfseries\large, text=popblue] (title) at (0,5.5) {Classification of Contraceptive Methods};

% Main Categories
\node[main] (nat) at (-7, 3.5) {Natural Methods};
\node[main] (bar) at (-2, 3.5) {Barrier Methods};
\node[main] (hor) at (3, 3.5) {Hormonal / Chemical Methods};
\node[main] (iud) at (8, 3.5) {Intrauterine Devices (IUD)};
\node[main] (sur) at (13, 3.5) {Surgical (Permanent) Methods};

% Sub-nodes Natural
\node[sub] (abs) at (-7, 1.8) {Abstinence};
\node[sub] (rhythm) at (-7, 0.5) {Rhythm / Calendar (Safe Period)};
\node[sub] (with) at (-7, -0.8) {Withdrawal (Coitus Interruptus)};
\node[sub] (lam) at (-7, -2.1) {Lactational Amenorrhea (LAM)};

% Sub-nodes Barrier
\node[sub] (mcond) at (-2, 2.2) {Male Condom};
\node[sub] (fcond) at (-2, 0.8) {Female Condom};
\node[sub] (diaph) at (-2, -0.6) {Diaphragm / Cervical Cap (+ Spermicide)};

% Sub-nodes Hormonal
\node[sub] (comb) at (3, 2.2) {Combined Oral Pill (Oestrogen + Progestin)};
\node[sub] (mini) at (3, 0.8) {Mini-Pill (Progestin Only)};
\node[sub] (inj) at (3, -0.6) {Injectables (e.g. Depo-Provera)};
\node[sub] (impl) at (3, -2.0) {Implants (e.g. Jadelle, Implanon)};
\node[sub] (ecp) at (3, -3.4) {Emergency Pill (Morning-After)};

% Sub-nodes IUD
\node[sub] (cu) at (8, 1.8) {Copper IUD (Cu-T)};
\node[sub] (hormiud) at (8, 0.3) {Hormonal IUS (Levonorgestrel)};

% Sub-nodes Surgical
\node[sub] (vas) at (13, 1.8) {Vasectomy (Male)};
\node[sub] (tub) at (13, 0.3) {Tubal Ligation (Female)};

% Site of Action Labels
\node[labelstyle] at (-10.5, 1.8) {\textbf{Site/Action:} Behavioural; avoids sperm-egg contact.};
\node[labelstyle] at (-5.5, 2.2) {\textbf{Site/Action:} Physical block at vagina/cervix. \textit{Only method vs STIs/HIV.}};
\node[labelstyle] at (0.5, 2.2) {\textbf{Site/Action:} Hypothalamus/Pituitary $\to$ Blocks Ovulation; thickens cervical mucus; thins endometrium.};
\node[labelstyle] at (5.5, 1.8) {\textbf{Site/Action:} Uterus $\to$ Spermicidal (Cu) / Local hormonal effect $\to$ prevents fertilisation/implantation.};
\node[labelstyle] at (10.5, 1.8) {\textbf{Site/Action:} Blocks gamete transport permanently (Vas deferens / Fallopian tubes).};

% Arrows
\draw[arrow] (nat.south) -- (abs.north);
\draw[arrow] (nat.south) -- (rhythm.north);
\draw[arrow] (nat.south) -- (with.north);
\draw[arrow] (nat.south) -- (lam.north);

\draw[arrow] (bar.south) -- (mcond.north);
\draw[arrow] (bar.south) -- (fcond.north);
\draw[arrow] (bar.south) -- (diaph.north);

\draw[arrow] (hor.south) -- (comb.north);
\draw[arrow] (hor.south) -- (mini.north);
\draw[arrow] (hor.south) -- (inj.north);
\draw[arrow] (hor.south) -- (impl.north);
\draw[arrow] (hor.south) -- (ecp.north);

\draw[arrow] (iud.south) -- (cu.north);
\draw[arrow] (iud.south) -- (hormiud.north);

\draw[arrow] (sur.south) -- (vas.north);
\draw[arrow] (sur.south) -- (tub.north);
\end{tikzpicture}
\legende{Classification of contraceptive methods by mechanism, showing site of action.}
\end{popfigure}

\courssub{Natural Methods}

These methods rely on observing physiological signs or behavioural discipline. They involve no devices or drugs.

\begin{experience}[The Menstrual Cycle and the Fertile Window]
\textbf{Principle:} Fertilisation can only occur if intercourse happens during the \cle{fertile window} — approximately 5 days before ovulation (sperm survival) to 24 hours after ovulation (ovum lifespan).
\textbf{Calendar (Rhythm) Method:} The woman records cycle lengths for 6–12 months.
\begin{itemize}
    \item Shortest cycle $-$ 18 = First fertile day.
    \item Longest cycle $-$ 11 = Last fertile day.
    \item Abstain or use barrier during this window.
\end{itemize}
\textbf{Basal Body Temperature (BBT):} Progesterone raises body temperature by $\approx 0.3^\circ\mathrm{C}$ after ovulation. Daily morning temperature charting confirms ovulation \emph{has occurred} (retrospective).
\textbf{Cervical Mucus (Billings/ovulation method):} Oestrogen $\to$ clear, stretchy, slippery mucus (spinnbarkeit) $\to$ fertile. Progesterone $\to$ thick, sticky, opaque $\to$ infertile.
\textbf{Symptothermal Method:} Combines BBT + Mucus + Calendar calculation $\to$ higher reliability.
\end{experience}

\begin{aretenir}[Natural Methods — Key Points]
\begin{itemize}
    \item \textbf{Abstinence:} 100\% effective if total; requires high motivation.
    \item \textbf{Withdrawal (Coitus Interruptus):} Penis withdrawn before ejaculation. \textbf{Unreliable:} Pre-ejaculatory fluid (Cowper's gland secretion) may contain viable sperm. High failure rate ($\approx 20\%$ typical use).
    \item \textbf{Lactational Amenorrhea Method (LAM):} Exclusive breastfeeding (on demand, day \& night) $\to$ high prolactin $\to$ suppresses GnRH $\to$ no LH surge $\to$ no ovulation. Effective ($>98\%$) only if: (1) Baby $< 6$ months, (2) Amenorrhoeic, (3) Exclusive breastfeeding.
    \item \textbf{Advantages:} No cost, no side effects, religiously acceptable.
    \item \textbf{Disadvantages:} High user failure rate; requires training, discipline, regular cycles; \textbf{NO protection against STIs/HIV}.
\end{itemize}
\end{aretenir}

\courssub{Barrier Methods}

These physically block sperm from reaching the ovum.

\begin{definition}[Barrier Contraception]
A \cle{barrier method} is a device that creates a physical block preventing spermatozoa deposited in the vagina from entering the uterus and reaching the fallopian tubes.
\end{definition}

\begin{center}
\begin{tabularx}{\linewidth}{|l|X|X|X|}
\hline
\rowcolor{popblueL}
\textbf{Method} & \textbf{Mechanism / Site} & \textbf{Effectiveness (Typical Use)} & \textbf{STI Protection} \\
\hline
Male Condom (Latex/Polyurethane) & Sheath over erect penis; collects semen. & $\approx 85\%$ & \textbf{YES (Best)} \\
\hline
Female Condom (Polyurethane/Nitrile) & Loose sheath lining vagina; inner ring at cervix, outer ring at vulva. & $\approx 79\%$ & \textbf{YES} \\
\hline
Diaphragm / Cervical Cap (+ Spermicide) & Dome/cap fits over cervix; spermicide immobilises sperm. Requires fitting by clinician. & $\approx 83\%$ & Partial (Herpes/HPV limited) \\
\hline
\end{tabularx}
\end{center}

\begin{attention}[Common GCE Trap]
\textbf{Only condoms (male and female) offer significant protection against Sexually Transmitted Infections (STIs) including HIV/AIDS.}
The Pill, IUD, Implants, Injectables, Diaphragm, Withdrawal \textbf{DO NOT} protect against STIs. This is a favourite distinction in GCE Paper 2 essay questions.
\end{attention}

\courssub{Hormonal Methods}

These exploit the \cle{negative feedback} control of the hypothalamic-pituitary-ovarian axis.

\begin{propriete}[Mechanism of Hormonal Contraception]
Exogenous \cle{oestrogen} and/or \cle{progestin} (synthetic progesterone) suppress the secretion of \cle{GnRH} from the hypothalamus and \cle{FSH/LH} from the anterior pituitary.
\begin{itemize}
    \item \textbf{Oestrogen component:} Inhibits FSH $\to$ prevents follicular development $\to$ \textbf{no ovulation}.
    \item \textbf{Progestin component:} Inhibits LH surge $\to$ \textbf{no ovulation}; thickens cervical mucus $\to$ blocks sperm penetration; thins endometrium $\to$ unfavourable for implantation.
\end{itemize}
\end{propriete}

\begin{popfigure}
\centering
\begin{tikzpicture}[
    node distance=1.5cm and 2cm,
    box/.style={draw=popdark, thick, rounded corners, fill=white, text width=2.5cm, align=center, minimum height=1cm},
    boxact/.style={draw=poporange, thick, rounded corners, fill=poporange!10, text width=2.5cm, align=center, minimum height=1cm},
    arr/.style={-Stealth, thick, popdark},
    arrinh/.style={-Stealth, thick, popred, dashed},
    label/.style={font=\small, align=center}
]
% Normal Axis
\node[box] (hypo) at (-5, 3) {Hypothalamus};
\node[box] (pit) at (-5, 1) {Ant. Pituitary};
\node[box] (ova) at (-5, -1) {Ovary};
\node[box] (uter) at (-5, -3) {Uterus};

\draw[arr] (hypo) -- node[left,label] {GnRH} (pit);
\draw[arr] (pit) -- node[left,label] {FSH / LH} (ova);
\draw[arr] (ova) -- node[left,label] {Oestrogen / Progesterone} (uter);
\draw[arrinh, bend left=20] (ova) to node[right,label] {Negative Feedback} (hypo);
\draw[arrinh, bend left=20] (ova) to node[right,label] {Negative Feedback} (pit);

% With Pill
\node[boxact, align=center] (pill) at (5, 3){Contraceptive Pill\\(Oestrogen + Progestin)};
\node[boxact, align=center] (hypo2) at (5, 1){Hypothalamus\\(Suppressed)};
\node[boxact, align=center] (pit2) at (5, -1){Ant. Pituitary\\(Low FSH/LH)};
\node[boxact, align=center] (ova2) at (5, -3){Ovary\\(No Follicle Growth\\(No Ovulation)};

\draw[arr] (pill) -- node[left,label] {High Constant Levels} (hypo2);
\draw[arrinh] (hypo2) -- node[left,label] {No GnRH} (pit2);
\draw[arrinh] (pit2) -- node[left,label] {No LH Surge} (ova2);

% Cross out normal ovulation
\node[font=\bfseries\Large, text=popred] at (-5, -4) {OVULATION};
\node[font=\bfseries\Large, text=popgreen] at (5, -4.5) {NO OVULATION};

% Title
\node[font=\bfseries, text=popblue] at (0, 4) {Negative Feedback Mechanism of the Combined Oral Contraceptive Pill};
\end{tikzpicture}
\legende{The pill maintains high steroid levels, mimicking the luteal phase, suppressing GnRH and the LH surge, thereby preventing ovulation.}
\end{popfigure}

\begin{center}
\begin{tabularx}{\linewidth}{|l|X|X|}
\hline
\rowcolor{popblueL}
\textbf{Method} & \textbf{Composition / Regimen} & \textbf{Key Features / Side Effects} \\
\hline
Combined Oral Contraceptive Pill (COC) & Oestrogen (Ethinylestradiol) + Progestin. 21 days on, 7 days off (withdrawal bleed). & High efficacy ($\approx 91\%$ typical). Regulates cycles, reduces dysmenorrhoea, acne. \textbf{Risks:} Thromboembolism (DVT/PE), hypertension, contraindicated in smokers $>35$ yrs, breastfeeding mums (reduces milk). \\
\hline
Progestin-Only Pill (POP / Mini-Pill) & Progestin only (e.g. Desogestrel). Taken continuously, same time daily ($\pm 3$ hrs). & Safe for breastfeeding, no oestrogen risks. Mainly thickens mucus; may suppress ovulation. Irregular bleeding common. \\
\hline
Injectables (e.g. Depo-Provera / Sayana Press) & DMPA (Depot Medroxyprogesterone Acetate) IM/SC every 12--13 weeks. & High efficacy ($\approx 94\%$). Amenorrhoea common. \textbf{Delay in return to fertility} (avg 9--10 months). Bone density loss with long use. \\
\hline
Implants (e.g. Jadelle 2-rods, Implanon 1-rod) & Levonorgestrel / Etonogestrel rods under skin of upper arm. 3--5 years. & Highest efficacy ($\approx 99.9\%$). Long-acting, reversible (LARC). Insertion/removal minor surgery. Irregular bleeding. \\
\hline
Emergency Contraceptive Pill (ECP) & Levonorgestrel 1.5mg (single dose) or Ulipristal Acetate. Within 72--120h post coitus. & Delays/blocks ovulation. \textbf{NOT an abortifacient} (does not disrupt implanted pregnancy). Not for regular use. \\
\hline
\end{tabularx}
\end{center}

\courssub{Intrauterine Devices (IUD / IUCD)}

\begin{definition}[Intrauterine Device]
An \cle{IUD} is a small, T-shaped device inserted into the uterine cavity by a trained provider. It provides long-acting reversible contraception (LARC).
\end{definition}

\begin{center}
\begin{tabularx}{\linewidth}{|l|X|X|}
\hline
\rowcolor{popblueL}
\textbf{Type} & \textbf{Mechanism} & \textbf{Duration / Notes} \\
\hline
Copper IUD (Cu-T 380A) & Copper ions $\to$ toxic to sperm (motility/viability) and ova; sterile inflammatory reaction in endometrium $\to$ prevents fertilisation \& implantation. & 10--12 years. Non-hormonal. Can be used as emergency contraception (inserted $\le 5$ days post coitus). Heavier, more painful periods initially. \\
\hline
Hormonal IUS (Levonorgestrel-IUS, e.g. Mirena) & Releases levonorgestrel locally $\to$ thickens cervical mucus, thins endometrium (often amenorrhoea), may suppress ovulation. & 5--7 years. Treats heavy menstrual bleeding (menorrhagia). Lower systemic hormones. \\
\hline
\end{tabularx}
\end{center}

\begin{attention}[IUD Misconceptions]
\begin{itemize}
    \item \textbf{Mechanism:} Primary action is \textbf{prevention of fertilisation} (spermicidal/ovicidal), not abortion. It does not disrupt an established pregnancy.
    \item \textbf{Infection Risk:} Slightly increased risk of PID (Pelvic Inflammatory Disease) only in first 20 days post-insertion (aseptic technique crucial). Not contraindicated in nulliparous women (modern guidelines).
    \item \textbf{Expulsion:} Check strings monthly after menses.
\end{itemize}
\end{attention}

\courssub{Surgical (Permanent) Methods}

These are \cle{sterilisation} procedures intended to be irreversible. Counselling is mandatory.

\begin{methode}[Vasectomy (Male Sterilisation)]
\begin{enumerate}
    \item Under local anaesthesia, a small incision or puncture is made in the scrotum.
    \item The \cle{vas deferens} (sperm duct) is isolated, a segment excised, and ends ligated/sealed (cautery/clips).
    \item \textbf{Not immediately effective:} Residual sperm proximal to cut. Must use backup contraception until \textbf{azoospermia confirmed by semen analysis} (usually $\approx 20$ ejaculates / 3 months).
    \item Testosterone production (Leydig cells) and sexual function unchanged.
\end{enumerate}
\end{methode}

\begin{methode}[Tubal Ligation (Female Sterilisation)]
\begin{enumerate}
    \item Usually laparoscopic (minilaparotomy post-partum).
    \item \cle{Fallopian tubes} are occluded: clips (Filshie), rings, coagulation, or excision (Pomeroy).
    \item \textbf{Effective immediately.}
    \item Ovarian function (hormones, ovulation) preserved; ovum degenerates in peritoneal cavity.
    \item Higher surgical risk than vasectomy (anaesthesia, bowel/bladder injury).
\end{enumerate}
\end{methode}

\begin{aretenir}[Sterilisation — Key Exam Points]
\begin{itemize}
    \item Considered \textbf{permanent}. Reversal (vasovasostomy, tubal reanastomosis) is microsurgery, expensive, not always successful, and not guaranteed by health systems.
    \item \textbf{No protection against STIs.}
    \item Ethical/Legal: In Cameroon, requires informed written consent, minimum age (often 21 or 2 children), spousal consent often requested (though legally debated), cooling-off period.
\end{itemize}
\end{aretenir}

\courssub{Comparative Evaluation and Exam Strategy}

\begin{methode}[Evaluating a Contraceptive Method — Three Criteria]
For GCE essays/tables, evaluate every method against these three pillars:
\begin{enumerate}
    \item \textbf{Effectiveness (Pearl Index):} \cle{Perfect Use} (method efficacy) vs \cle{Typical Use} (real world). LARC (Implants, IUD, Sterilisation) $\approx$ Perfect Use. User-dependent methods (Pills, Condoms, Natural) have wide gap.
    \item \textbf{Reversibility:} How quickly does fertility return?
        \begin{itemize}
            \item Immediate: Condoms, Diaphragm, IUD removal, Implant removal, Pills (usually 1--3 cycles).
            \item Delayed: Injectables (months).
            \item Permanent/Difficult: Sterilisation.
        \end{itemize}
    \item \textbf{STI/HIV Protection:} Only Male/Female Condoms.
\end{enumerate}
\end{methode}

\begin{exoresolu}[GCE Style Question]
\textbf{Question:} A newly married couple in Buea wants to delay their first pregnancy for 3 years. The husband travels frequently. The wife has a history of migraine with aura. They want a method that is highly effective, reversible, and does not require daily action. Recommend the most suitable method and justify your choice by comparing it with two other methods.

\textbf{Solution:}
\textbf{Recommendation:} \textbf{Subdermal Implant (e.g. Jadelle / Implanon NXT)}.

\textbf{Justification (Comparison Table):}

\begin{center}
\begin{tabularx}{\linewidth}{|l|X|X|X|}
\hline
\rowcolor{popblueL}
\textbf{Criterion} & \textbf{Subdermal Implant} & \textbf{Combined Oral Pill (COC)} & \textbf{Copper IUD} \\
\hline
\textbf{Effectiveness} & $>99.9\%$ (LARC, user-independent) & $\approx 91\%$ (User-dependent, missed pills) & $>99\%$ (LARC) \\
\hline
\textbf{Duration / Action} & 3--5 years; insert \& forget & Daily pill at same time & 10--12 years; insert \& forget \\
\hline
\textbf{Migraine with Aura} & \textbf{SAFE} (Progestin only) & \textbf{CONTRAINDICATED} (Oestrogen $\uparrow$ stroke risk) & \textbf{SAFE} (Non-hormonal) \\
\hline
\textbf{Husband Travel} & No adherence needed during absence & Requires daily discipline; missed pills on travel & No adherence needed \\
\hline
\textbf{Reversibility} & Rapid return (weeks) & Rapid return (1--3 cycles) & Rapid return (immediate) \\
\hline
\textbf{STI Protection} & No & No & No \\
\hline
\textbf{Side Effects} & Irregular bleeding (manageable) & Nausea, weight, thromboembolism risk & Heavier/painful periods \\
\hline
\end{tabularx}
\end{center}

\textbf{Conclusion:} The Implant is optimal: safe for her medical history, unaffected by husband's travel, highly effective for the 3-year span, and easily removable when they desire pregnancy. The IUD is a strong second choice but may worsen dysmenorrhoea; the COC is medically unsafe.
\end{exoresolu}

\courssub{Ethical, Religious and Health Considerations}

\begin{savaistu}[Discussion Points for GCE Paper 2 Essays]
\begin{itemize}
    \item \textbf{Religious Views:} Catholic Church teaches against artificial contraception (Humanae Vitae), promoting Natural Family Planning (NFP). Many Protestant and Muslim leaders support spacing for maternal health but may oppose sterilisation/abortifacients.
    \item \textbf{Adolescent Access:} Cameroon law allows adolescents access to reproductive health services. Barrier: provider bias, stigma, cost. School-based programmes vs parental consent debates.
    \item \textbf{Gender Burden:} Most methods (Pill, Injectables, Implants, IUD, Tubal ligation) are female-controlled. Vasectomy and Male Condom are only male methods. Promoting \emph{shared responsibility} is a public health goal.
    \item \textbf{Abortion vs Contraception:} Emergency Contraception (ECP) prevents ovulation; it is \textbf{not} medical abortion (Mifepristone/Misoprostol). Clarifying this distinction reduces stigma and increases ECP uptake.
    \item \textbf{Health System Integration:} Contraception integrated into ANC (post-partum), PNC, HIV clinics (dual protection), youth corners. Task-shifting (nurses inserting implants/IUDs) expands access.
\end{itemize}
\end{savaistu}

\begin{exercice}[Practice Questions]
\begin{itemize}
\item Define contraception and distinguish between contraception and contragestion.
\item Explain the physiological basis of the Lactational Amenorrhea Method (LAM). State the three criteria for its effectiveness.
\item Describe the mechanism of action of the Combined Oral Contraceptive Pill with reference to negative feedback.
\item Compare and contrast the Copper IUD and the Levonorgestrel IUS (Hormonal IUD) in terms of mechanism, duration, and non-contraceptive benefits.
\item A 35-year-old woman, smoker (15 cigarettes/day), requests the Combined Pill. Advise her with reasons.
\item Outline the procedure for Vasectomy. Why is it not immediately effective?
\item Discuss the statement: "The most effective contraceptive methods are those that do not depend on user compliance."
\end{itemize}
\end{exercice}

\begin{corrige}[Exercise 1]
\textbf{Contraception} is the deliberate prevention of fertilisation or implantation. \textbf{Contragestion} (interception) refers to agents acting after fertilisation but before implantation (e.g. high-dose progestins, IUD insertion post-coitus) to prevent implantation. The distinction is subtle; most modern definitions classify IUD/ECP as contraception (primary action pre-fertilisation).
\end{corrige}

\begin{corrige}[Exercise 2]
LAM: Suckling $\to$ Prolactin $\uparrow$ $\to$ GnRH pulse frequency $\downarrow$ $\to$ LH surge blocked $\to$ No ovulation.
Criteria: (1) Infant $< 6$ months, (2) Amenorrhoea (no menses), (3) Exclusive/near-exclusive breastfeeding on demand (day \& night).
\end{corrige}

\begin{corrige}[Exercise 3]
COC provides constant exogenous oestrogen/progestin $\to$ Negative feedback on Hypothalamus (GnRH $\downarrow$) and Ant. Pituitary (FSH/LH $\downarrow$) $\to$ No follicular recruitment (FSH low) $\to$ No LH surge $\to$ \textbf{Anovulation}. Progestin also thickens cervical mucus (sperm block) and atrophies endometrium (implantation block).
\end{corrige}

\begin{corrige}[Exercise 4]
Cu-IUD: Copper toxic to gametes, sterile inflammation $\to$ prevents fertilisation. 10+ yrs. Non-hormonal. Emergency contraception. Side effect: heavier menses.
LNG-IUS: Local progestin $\to$ thick mucus, thin endometrium, maybe anovulation. 5--7 yrs. Treats menorrhagia, dysmenorrhoea, endometriosis. Side effect: irregular spotting $\to$ amenorrhoea.
\end{corrige}

\begin{corrige}[Exercise 5]
\textbf{Advise against COC.} Oestrogen + Smoking + Age $>35$ = \textbf{High risk of cardiovascular events} (Venous Thromboembolism VTE, Stroke, MI). Absolute contraindication (UKMEC Category 4). Offer Progestin-only methods (POP, Implant, Injectables, LNG-IUS) or Copper IUD.
\end{corrige}

\begin{corrige}[Exercise 6]
Local anaesthesia $\to$ access vas deferens via scrotal incision/puncture $\to$ excise segment $\to$ ligate/cauterise ends. Not immediately effective because sperm remain in distal vas deferens/ampulla. Requires $\approx 20$ ejaculates / 3 months + semen analysis confirming azoospermia before relying on it.
\end{corrige}

\begin{corrige}[Exercise 7]
\textbf{Agree.} LARC methods (Implants, IUDs, Sterilisation) have Pearl Indices $\approx 0.1-0.8$ (Perfect $\approx$ Typical). User-dependent methods (Pills, Condoms, Natural) have Typical Use failure rates 7--24\% due to human error (forgetting, incorrect use, breakage). LARC removes "user failure". However, "best" method includes acceptability, side effects, STI protection (condoms), and medical eligibility.
\end{corrige}

\coursec{Life Cycles: Alternation of Generations and Metagenesis}

In the previous section we examined how human beings control their own reproduction. But reproduction is not always a simple matter of two parents producing offspring of the same kind. Walk through the humid forest on the slopes of Mount Cameroon after the rains: the green carpet of mosses and the delicate fronds of ferns are reproducing too — yet they do it through a remarkable strategy in which \emph{two different kinds of plant}, one haploid and one diploid, take turns to produce each other. Even some animals, such as jellyfish, alternate between two body forms across their life cycle. Understanding these strategies is essential for the GCE, where life-cycle diagrams with \cle{n} and \cle{2n} labels are a favourite examination topic.

\courssub{The Concept of a Life Cycle: Haploid and Diploid Phases}

\begin{definition}[Life cycle]
The \cle{life cycle} of an organism is the complete sequence of stages through which an individual passes, from its formation (by fertilization or by asexual reproduction) until it in turn produces the next generation.
\end{definition}

Every sexually reproducing organism must solve a chromosome problem. If gametes contained the full diploid number of chromosomes ($2n$), the chromosome number would double at every generation. The solution is a rhythm of two key events:

\begin{itemize}
\item \cle{Meiosis} halves the chromosome number: a diploid cell ($2n$) produces haploid cells ($n$).
\item \cle{Fertilization} (syngamy) restores the diploid number: two haploid gametes ($n + n$) fuse to form a diploid zygote ($2n$).
\end{itemize}

In animals like ourselves, meiosis directly produces the gametes, and the only haploid cells in the whole life cycle are the sperm and the egg. Plants (and some algae) do something different and far more elaborate: meiosis produces not gametes but \cle{spores}, and each spore grows by mitosis into a whole haploid organism which \emph{then} produces gametes by mitosis. This creates two alternating multicellular generations.

\begin{aretenir}[The three key rules of any life cycle]
\begin{itemize}
\item \textbf{Rule 1:} Meiosis is the \textbf{only} event that changes $2n \rightarrow n$.
\item \textbf{Rule 2:} Fertilization is the \textbf{only} event that changes $n \rightarrow 2n$.
\item \textbf{Rule 3:} Everything else (growth of gametophyte or sporophyte, gamete formation, budding, larval development) occurs by \textbf{mitosis}, which conserves the chromosome number.
\end{itemize}
\end{aretenir}

\courssub{Alternation of Generations in Plants}

\begin{definition}[Alternation of generations]
\cle{Alternation of generations} is the regular alternation, in the life cycle of a plant, between a haploid ($n$) gamete-producing generation called the \cle{gametophyte} and a diploid ($2n$) spore-producing generation called the \cle{sporophyte}. The sporophyte produces haploid spores by meiosis; each spore germinates into a gametophyte, which produces gametes by mitosis; fusion of gametes restores the diploid sporophyte.
\end{definition}

Notice the logic: the \emph{sporophyte} is named after what it makes (spores, by meiosis), and the \emph{gametophyte} after what it makes (gametes, by mitosis). The two generations differ in ploidy, in appearance, and often in which one is dominant.

\textbf{The life cycle of a moss (e.g. \emph{Funaria})}\par
The green, leafy moss plant you see on damp walls and stones in Buea during the rainy season is the \textbf{gametophyte} ($n$) — it is the dominant, photosynthetic, independent generation. It develops from a filamentous \cle{protonema} that germinates from the spore. At its tip it bears the sex organs:

\begin{itemize}
\item \cle{Antheridia} (male organs) release many motile, biflagellate sperm cells.
\item \cle{Archegonia} (female organs), flask-shaped, each contain a single egg cell.
\end{itemize}

Fertilization \textbf{requires liquid water}: the sperm must swim through a film of rain or dew to reach the egg inside the archegonium. This dependence on water is why mosses are restricted to moist habitats. The resulting zygote ($2n$) is not released; it grows \emph{in place}, on top of the gametophyte, into the \textbf{sporophyte}: a thin stalk (seta) carrying a capsule. The sporophyte is \textbf{dependent} — it has no roots of its own and draws water and nutrients from the gametophyte. Inside the capsule, spore mother cells undergo \textbf{meiosis} to produce haploid spores, which are scattered by wind and germinate into a new protonema and then a green gametophyte.

\textbf{The life cycle of a fern (e.g. \emph{Dryopteris})}\par
In ferns the roles are reversed. The familiar fern plant with large divided fronds is the \textbf{sporophyte} ($2n$) — dominant, independent, with true roots, stems and leaves. On the underside of the fronds you can see brown dots called \cle{sori} (singular: sorus); each sorus is a cluster of \cle{sporangia} in which meiosis produces haploid \cle{spores}. A spore that lands on moist soil germinates into the gametophyte: a small, flat, heart-shaped green plate called the \cle{prothallus}, only a few millimetres across. The prothallus is independent and photosynthetic, but short-lived and delicate. It bears antheridia and archegonia on its lower surface; sperm swim through a water film to the egg; the zygote ($2n$) grows on the prothallus into a young fern, and the prothallus withers away.

\begin{popfigure}
\begin{center}
\begin{tikzpicture}[scale=1.1, every node/.style={font=\small}]
% Nodes
\node[draw, fill=popgreenL, rounded corners, align=center, text width=2.8cm] (sporo) at (90:3.2) {\textbf{Sporophyte} ($2n$)\\ Fern plant with fronds};
\node[draw, fill=popblueL, rounded corners, align=center, text width=2.4cm] (sori) at (30:3.6) {\textbf{Sori} on frond\\ (clusters of sporangia)};
\node[draw, fill=popblueL, rounded corners, align=center, text width=2.2cm] (spore) at (-30:3.6) {\textbf{Spores} ($n$)\\ wind-dispersed};
\node[draw, fill=popgreenL, rounded corners, align=center, text width=2.8cm] (proth) at (-90:3.6) {\textbf{Prothallus} ($n$)\\ Gametophyte (heart-shaped)};
\node[draw, fill=popblueL, rounded corners, align=center, text width=2.4cm] (gam) at (-150:3.6) {\textbf{Gametes} ($n$)\\ Antheridia \& Archegonia};
\node[draw, fill=popgreenL, rounded corners, align=center, text width=2.2cm] (zyg) at (150:3.6) {\textbf{Zygote} ($2n$)\\ Young sporophyte};
% Arrows
\draw[->, thick, popblue] (sporo) -- (sori) node[midway, above right, font=\footnotesize, popmuted] {mitosis};
\draw[->, thick, popblue] (sori) -- (spore) node[midway, right, font=\footnotesize, popmuted, align=center]{\textbf{MEIOSIS}\\ $2n \rightarrow n$};
\draw[->, thick, popblue] (spore) -- (proth) node[midway, left, font=\footnotesize, popmuted] {germination (mitosis)};
\draw[->, thick, popblue] (proth) -- (gam) node[midway, left, font=\footnotesize, popmuted] {mitosis};
\draw[->, thick, popblue] (gam) -- (zyg) node[midway, left, font=\footnotesize, popmuted, align=center]{\textbf{FERTILIZATION}\\ $n+n \rightarrow 2n$ (water needed)};
\draw[->, thick, popblue] (zyg) -- (sporo) node[midway, above left, font=\footnotesize, popmuted] {growth (mitosis)};
% Event labels
\node[fill=poppink, rounded corners, align=center, font=\footnotesize] at (0:4.8) {\textbf{MEIOSIS}\\ in sporangia};
\node[fill=popgold, rounded corners, align=center, font=\footnotesize] at (180:4.8) {\textbf{FERTILIZATION}\\ on prothallus};
\end{tikzpicture}
\end{center}
\legende{Life cycle of a fern (\emph{Dryopteris}). The diploid sporophyte is the dominant, conspicuous plant. Meiosis in sporangia (within sori) produces haploid spores. These germinate into the small, independent haploid prothallus (gametophyte) which produces gametes by mitosis. Fertilization requires water for sperm to swim to the egg.}
\end{popfigure}

\begin{propriete}[Dominance of the generations — GCE comparison]
In \textbf{mosses} (bryophytes), the \cle{gametophyte} ($n$) is the dominant, independent, photosynthetic generation, while the sporophyte is small, short-lived and nutritionally dependent on the gametophyte. In \textbf{ferns} (pteridophytes), the \cle{sporophyte} ($2n$) is dominant, independent and long-lived, while the gametophyte (prothallus) is small, short-lived but still independent and free-living. (This paired comparison is frequently examined; learn it as a set of contrasting statements.)
\end{propriete}

\courssub{The Evolutionary Trend: From Mosses to Flowering Plants}

Comparing mosses, ferns and flowering plants reveals one of the clearest trends in plant evolution:

\begin{center}
\begin{tabularx}{\linewidth}{|l|X|X|X|}
\hline
\rowcolor{popblueL} \textbf{Group} & \textbf{Dominant generation} & \textbf{Gametophyte} & \textbf{Water needed for fertilization?} \\
\hline
Mosses (Bryophytes) & Gametophyte ($n$) & Large, independent (the leafy plant + protonema) & Yes — flagellated sperm swim \\
\hline
Ferns (Pteridophytes) & Sporophyte ($2n$) & Small prothallus, independent, free-living & Yes — flagellated sperm swim \\
\hline
Flowering plants (Angiosperms) & Sporophyte ($2n$) & Extremely reduced: pollen grain (male) and embryo sac (female), both dependent on the sporophyte & No — pollen tube delivers male nuclei \\
\hline
\end{tabularx}
\end{center}

\begin{aretenir}[The trend to remember for GCE]
From mosses to flowering plants, the \textbf{sporophyte becomes increasingly dominant} while the \textbf{gametophyte becomes increasingly reduced and dependent}. In angiosperms, the male gametophyte is reduced to the \cle{pollen grain} (and the pollen tube it produces), and the female gametophyte to the \cle{embryo sac} inside the ovule. The conquest of dry land was possible because pollination by wind or insects replaced the need for liquid water at fertilization.
\end{aretenir}

\begin{savaistu}[Why this matters in Cameroon]
The mist forests of Mount Cameroon and the Bakossi Mountains are famous for their richness in ferns and mosses. These plants thrive there precisely because the permanent humidity provides the water film their swimming sperm need — a direct, visible consequence of alternation of generations. Flowering plants, freed from this constraint by the pollen tube, dominate the drier savanna regions of the North.
\end{savaistu}

\courssub{Metagenesis in Animals}

Some animals also alternate between two generations, but of a very different nature. Consider the colonial cnidarian \emph{Obelia}, a relative of jellyfish.

\begin{definition}[Metagenesis]
\cle{Metagenesis} is the alternation, in the life cycle of certain animals (especially cnidarians), between an \textbf{asexual generation} (typically a sessile \cle{polyp}) and a \textbf{sexual generation} (typically a free-swimming \cle{medusa}). Both generations are diploid; there is \textbf{no change in ploidy} between them.
\end{definition}

In \emph{Obelia}, the life cycle runs as follows:

\begin{enumerate}
\item The \textbf{polyp colony} ($2n$), fixed to rocks or seaweed, grows and produces new individuals by \textbf{budding} — asexual reproduction (mitosis). Specialised reproductive buds (blastostyles) produce \textbf{medusae} by budding, which detach and swim away.
\item The free-swimming \textbf{medusa} ($2n$) is the sexual stage: it produces gametes ($n$) by meiosis in its gonads (located on the manubrium).
\item \textbf{Fertilization} of the egg produces a zygote ($2n$), which develops into a ciliated, free-swimming larva called the \cle{planula}.
\item The planula settles on a solid surface and grows, by budding (mitosis), into a new polyp colony — completing the cycle.
\end{enumerate}

\begin{popfigure}
\begin{center}
\begin{tikzpicture}[scale=1.1, every node/.style={font=\small}]
% Nodes
\node[draw, fill=popgreenL, rounded corners, align=center, text width=2.8cm] (polyp) at (90:3.2) {\textbf{Polyp colony} ($2n$)\\ Sessile, asexual (budding)};
\node[draw, fill=popblueL, rounded corners, align=center, text width=2.4cm] (medusa) at (30:3.6) {\textbf{Medusa} ($2n$)\\ Free-swimming, sexual};
\node[draw, fill=popblueL, rounded corners, align=center, text width=2.2cm] (gam) at (-30:3.6) {\textbf{Gametes} ($n$)\\ Sperm + Egg};
\node[draw, fill=popgreenL, rounded corners, align=center, text width=2.2cm] (zyg) at (-90:3.6) {\textbf{Zygote} ($2n$)};
\node[draw, fill=popblueL, rounded corners, align=center, text width=2.4cm] (planula) at (-150:3.6) {\textbf{Planula larva} ($2n$)\\ Ciliated, swims};
% Arrows
\draw[->, thick, poppurple] (polyp) -- (medusa) node[midway, above right, font=\footnotesize, popmuted, align=center]{\textbf{BUDDING}\\ (asexual, mitosis)};
\draw[->, thick, poppurple] (medusa) -- (gam) node[midway, right, font=\footnotesize, popmuted, align=center]{\textbf{MEIOSIS}\\ $2n \rightarrow n$};
\draw[->, thick, poppurple] (gam) -- (zyg) node[midway, right, font=\footnotesize, popmuted, align=center]{\textbf{FERTILIZATION}\\ $n+n \rightarrow 2n$};
\draw[->, thick, poppurple] (zyg) -- (planula) node[midway, left, font=\footnotesize, popmuted] {cleavage (mitosis)};
\draw[->, thick, poppurple] (planula) -- (polyp) node[midway, left, font=\footnotesize, popmuted] {settling \& growth (mitosis)};
% Event labels
\node[fill=popgold, rounded corners, align=center, font=\footnotesize] at (0:4.8) {\textbf{BUDDING}\\ (mitosis)};
\node[fill=poppink, rounded corners, align=center, font=\footnotesize] at (-60:4.8) {\textbf{MEIOSIS}\\ in medusa gonads};
\node[fill=popgold, rounded corners, align=center, font=\footnotesize] at (180:4.8) {\textbf{FERTILIZATION}\\ in water};
\end{tikzpicture}
\end{center}
\legende{Life cycle of \emph{Obelia}, illustrating metagenesis: a sessile asexual polyp colony buds off sexual medusae; gametes fuse to give a zygote, which develops into a planula larva that settles and forms a new colony. Note that polyp, medusa, zygote and planula are all diploid ($2n$). Only the gametes are haploid ($n$).}
\end{popfigure}

\begin{attention}[Do not confuse metagenesis with alternation of generations — GCE TRAP]
This is one of the most frequent errors at the GCE. In \textbf{alternation of generations} (plants), the two generations differ in \textbf{ploidy}: haploid gametophyte versus diploid sporophyte. In \textbf{metagenesis} (animals), the two generations (polyp and medusa) are \textbf{both diploid}; they differ only in body form and in the \emph{mode} of reproduction (asexual budding versus sexual reproduction). The only haploid cells in metagenesis are the gametes, as in any animal. Writing ``the medusa is haploid'' or ``meiosis produces the polyp'' will cost you marks.
\end{attention}

\begin{center}
\begin{tabularx}{\linewidth}{|l|X|X|}
\hline
\rowcolor{popblueL} \textbf{Feature} & \textbf{Alternation of Generations (Plants)} & \textbf{Metagenesis (Animals, e.g. Cnidarians)} \\
\hline
Ploidy of the two generations & Haploid ($n$) gametophyte $\leftrightarrow$ Diploid ($2n$) sporophyte & Both generations diploid ($2n$) \\
\hline
Chromosome number change & Yes: Meiosis ($2n \rightarrow n$) and Fertilization ($n \rightarrow 2n$) & Only at gamete formation (meiosis) and fertilization \\
\hline
Nature of generations & Morphologically distinct, often one dominant & Morphologically distinct (polyp vs medusa) \\
\hline
Reproductive mode of each & Gametophyte: sexual (gametes); Sporophyte: asexual (spores) & Polyp: asexual (budding); Medusa: sexual (gametes) \\
\hline
Examples & Mosses, ferns, flowering plants & \emph{Obelia}, \emph{Aurelia} (jellyfish), some flatworms \\
\hline
\end{tabularx}
\end{center}

\courssub{Biological Significance of These Strategies}

Why should organisms bother with two generations at all? Both strategies confer real advantages:

\begin{itemize}
\item \textbf{Dispersal.} In plants, light spores (and later pollen and seeds) scatter the species far from the parent. In \emph{Obelia}, the sessile colony cannot move, but the swimming medusa and planula larva carry the species to new substrates.
\item \textbf{Multiplication.} The asexual phase (sporophyte producing many spores by meiosis; polyp budding many medusae by mitosis) multiplies the output of a single successful fertilization many times over.
\item \textbf{Survival in varying conditions.} Each generation is adapted to a different way of life: the resistant sporophyte of the fern withstands dry air, while the delicate prothallus exploits moist microhabitats; the polyp resists and colonises, the medusa disperses.
\item \textbf{Genetic mixing.} The sexual phase maintains genetic recombination, providing the variation on which natural selection acts.
\end{itemize}

\begin{exoresolu}[Analysing a life-cycle diagram (GCE style)]
A diagram of the fern life cycle shows the following structures: (A) frond with sori, (B) spore, (C) prothallus, (D) egg and sperm, (E) zygote. (1) State the ploidy ($n$ or $2n$) of each structure. (2) Indicate between which structures meiosis and fertilization occur. (3) Name the dominant generation.
\tcblower
\textbf{Solution.}
\begin{enumerate}
\item (A) Frond with sori: $2n$ (it belongs to the sporophyte). (B) Spore: $n$ (product of meiosis). (C) Prothallus: $n$ (gametophyte, grown by mitosis from the spore). (D) Egg and sperm: $n$ (gametes, produced by mitosis on the gametophyte). (E) Zygote: $2n$ (product of fertilization).
\item \textbf{Meiosis} occurs in the sporangia of the sori, i.e. in the transition from (A) to (B) ($2n \rightarrow n$). \textbf{Fertilization} occurs when sperm and egg fuse, i.e. from (D) to (E) ($n + n \rightarrow 2n$).
\item The dominant generation is the \textbf{sporophyte} (the fern plant bearing the fronds), which is independent and photosynthetic; the prothallus is small and short-lived.
\end{enumerate}
\textbf{Marking tip:} one mark per correct ploidy, one for each correctly placed event, one for the dominant generation — an easy six marks if the three key rules are applied methodically.
\end{exoresolu}

\begin{methode}[How to analyse any life-cycle diagram in the exam]
\begin{enumerate}
\item Locate \textbf{meiosis} and \textbf{fertilization} first — they are the only two points where ploidy changes.
\item Label everything between fertilization and meiosis as $2n$ (sporophyte side in plants; polyp/medusa/zygote/planula side in metagenesis); label everything between meiosis and fertilization as $n$ (gametophyte side in plants; gametes only in metagenesis).
\item Identify which generation is \textbf{dominant} (largest, longest-lived, independent) — for plants only.
\item For animals, check: if both generations are diploid and differ only in reproductive mode, it is \textbf{metagenesis}, not alternation of generations.
\end{enumerate}
\end{methode}

\begin{exercice}[Checking your understanding]
\begin{enumerate}
\item Define alternation of generations and state how it differs from metagenesis.
\item In the moss \emph{Funaria}, state the ploidy of: the leafy plant, the capsule, the spore, the sperm, the protonema.
\item Explain why mosses and ferns are restricted to moist habitats while flowering plants are not.
\item Draw and fully label the life cycle of \emph{Obelia}, marking the points of meiosis and fertilization, and indicating the ploidy of each stage.
\item A student claims: ``In the fern life cycle, the prothallus is diploid because it grows from the zygote.'' Explain why this statement is incorrect.
\end{enumerate}
\end{exercice}

\begin{corrige}[Exercise 1]
\begin{enumerate}
\item Alternation of generations: regular alternation of a haploid gametophyte (produces gametes by mitosis) and a diploid sporophyte (produces spores by meiosis) in a plant life cycle. Metagenesis (animals) alternates an asexual polyp and a sexual medusa, \textbf{both diploid} — no change of ploidy between the two generations; meiosis occurs only during gamete formation in the medusa.
\item Leafy plant (gametophyte): $n$; capsule (part of the sporophyte): $2n$; spore: $n$; sperm: $n$; protonema: $n$ (grows from spore by mitosis).
\item In mosses and ferns the sperm are flagellated and must \textbf{swim through liquid water} to reach the egg, so reproduction is only possible in moist conditions. In flowering plants the \textbf{pollen tube} carries the male gametes to the embryo sac, making fertilization independent of external water.
\item The diagram must show: polyp colony ($2n$) $\xrightarrow{\text{budding (mitosis)}}$ medusa ($2n$) $\xrightarrow{\text{meiosis}}$ gametes ($n$) $\xrightarrow{\text{fertilization}}$ zygote ($2n$) $\xrightarrow{\text{cleavage}}$ planula ($2n$) $\xrightarrow{\text{settling \& growth}}$ new polyp colony. Meiosis is placed at gamete formation in the medusa; fertilization at the fusion of gametes.
\item The prothallus grows from the \textbf{spore} (which is $n$, produced by meiosis), not from the zygote. The zygote grows into the young sporophyte ($2n$). The student has confused the direction of the cycle.
\end{enumerate}
\end{corrige}

\begin{aretenir}[Summary of the section — GCE Essentials]
\begin{itemize}
\item A life cycle alternates haploid ($n$) and diploid ($2n$) phases; only \textbf{meiosis} and \textbf{fertilization} change the ploidy; all other growth is by mitosis.
\item \textbf{Alternation of generations} (plants): haploid gametophyte $\rightleftharpoons$ diploid sporophyte. Mosses: gametophyte dominant, sporophyte dependent. Ferns: sporophyte dominant, prothallus small but independent.
\item Evolutionary trend: increasing dominance of the sporophyte and reduction of the gametophyte, ending in the pollen grain and embryo sac of flowering plants.
\item \textbf{Metagenesis} (animals, e.g. \emph{Obelia}): asexual polyp $\rightleftharpoons$ sexual medusa, both diploid — do not confuse with alternation of generations.
\item Both strategies improve \textbf{dispersal}, \textbf{multiplication} and \textbf{survival}.
\item \textbf{Exam technique:} On any life-cycle diagram, first mark meiosis ($2n \rightarrow n$) and fertilization ($n \rightarrow 2n$), then deduce ploidy of all other stages.
\end{itemize}
\end{aretenir}

\coursec{Growth: Curves, Patterns and Hormonal Control}

In the previous section we followed organisms through their \cle{life cycles}, seeing how a fern alternates between sporophyte and gametophyte, and how a jellyfish passes through a polyp stage. Every stage of every life cycle has one thing in common: the organism must \emph{grow} to reach it. A maize seedling in a field near Bafoussam becomes a two-metre plant in a few months; a newborn baby in Yaoundé becomes an adult of perhaps \SI{170}{cm} in about eighteen years; a caterpillar becomes a butterfly only after a series of moults. How do we measure this growth, what shape does it follow, and what controls it? These are the questions of this final section.

\courssub{What is Growth?}

\textbf{Intuition.} If you weigh a baby every month, the mass increases. But if you soak a dry sponge in water, its mass also increases --- yet the sponge has not grown. True growth is not just ``getting bigger''; it is a \emph{permanent} increase in the amount of living material, built from new cells and new cytoplasm.

\begin{definition}[Growth]
\cle{Growth} is the \textbf{permanent and irreversible increase in size and dry mass} of an organism (or of one of its parts), resulting from three cellular processes:
\begin{itemize}
\item \cle{cell division} (mitosis), which increases the \emph{number} of cells;
\item \cle{cell enlargement} (expansion), which increases the \emph{size} of cells --- in plants mainly by uptake of water into the vacuole (vacuolation);
\item \cle{cell differentiation}, by which cells become specialised in structure for particular functions.
\end{itemize}
\end{definition}

\begin{attention}[Growth is not just swelling]
A seed soaked in water increases in mass and volume, but this is \emph{imbibition}, a reversible physical uptake of water --- not growth. Growth implies the synthesis of new protoplasm and is irreversible: an adult who loses water through sweating has not ``un-grown''.
\end{attention}

\courssub{Measuring Growth}

To study growth scientifically we must measure a chosen parameter repeatedly over time. Three measures are commonly used.

\begin{center}
\begin{tabularx}{\linewidth}{|l|X|X|}
\hline
\rowcolor{popblueL} \textbf{Measure} & \textbf{Advantages} & \textbf{Limitations} \\
\hline
Length / height & Quick, cheap, non-destructive; ideal for seedlings and children & One dimension only; ignores changes in width and mass \\
\hline
Fresh mass & Easy (weigh the organism); non-destructive & Fluctuates with water content (meals, rain, transpiration); unreliable \\
\hline
Dry mass & The \textbf{true measure} of growth: mass of organic matter actually synthesised & Destructive (the organism must be killed and dried in an oven at about $\SI{100}{\ensuremath{{}^{\circ}\mathrm{C}}}$ to constant mass); needs several identical samples \\
\hline
\end{tabularx}
\end{center}

\begin{attention}[Fresh mass versus dry mass --- the classic trap]
A very common examination error is to use fresh mass as evidence of growth. Fresh mass rises after drinking or rainfall and falls during drought, \emph{without any change in the amount of living tissue}. \textbf{Dry mass is the only reliable measure of true growth}, because it measures the organic material actually built by the organism.
\end{attention}

\begin{methode}[Plotting and interpreting a growth curve]
To obtain and analyse a growth curve (e.g. height of maize seedlings):
\begin{enumerate}
\item Germinate a large batch of genetically similar seeds under identical conditions.
\item At regular intervals (e.g. every 3 days), measure the chosen parameter (height, or dry mass of a sample of seedlings).
\item Plot the parameter (y-axis) against time (x-axis).
\item Identify the phases: slow start, rapid middle, levelling off.
\item Calculate the \cle{growth rate} at any time as the gradient of the curve: $\text{rate} = \dfrac{\Delta \text{size}}{\Delta t}$. The steepest part of the curve corresponds to the maximum growth rate.
\end{enumerate}
\end{methode}

\courssub{The Sigmoid Growth Curve}

When the size of an organism (a maize plant, a yeast population, a child measured over years) is plotted against time, the curve obtained is usually \textbf{S-shaped}, called a \cle{sigmoid curve}. It shows four phases.

\begin{popfigure}
\begin{center}
\begin{tikzpicture}
\begin{axis}[
width=11cm, height=7cm,
xlabel={Time}, ylabel={Size (dry mass / height)},
xmin=0, xmax=10, ymin=0, ymax=10.5,
xtick=\empty, ytick=\empty,
axis lines=left, thick]
\addplot[popblue, very thick, smooth, domain=0:10, samples=100] {9/(1+exp(-1.1*(x-5)))};
\draw[dashed, popmuted] (axis cs:2.5,0) -- (axis cs:2.5,1.4);
\draw[dashed, popmuted] (axis cs:5,0) -- (axis cs:5,4.5);
\draw[dashed, popmuted] (axis cs:7.5,0) -- (axis cs:7.5,7.6);
\node[popdark, align=center, font=\small] at (axis cs:1.25,2.6) {Lag\\ phase};
\node[popdark, align=center, font=\small] at (axis cs:3.7,3.1) {Exponential\\ (log) phase};
\node[popdark, align=center, font=\small] at (axis cs:6.3,6.2) {Decelerating\\ phase};
\node[popdark, align=center, font=\small] at (axis cs:8.9,8.2) {Stationary\\ phase};
\end{axis}
\end{tikzpicture}
\end{center}
\legende{The sigmoid (S-shaped) growth curve with its four phases.}
\end{popfigure}

\begin{aretenir}[Phases of the sigmoid curve]
\begin{itemize}
\item \textbf{Lag phase:} slow growth; the organism is adjusting, cells are small and few, enzymes and organelles are being built up.
\item \textbf{Exponential (log) phase:} growth rate is maximal; cells divide rapidly and resources are abundant. Size increases by a constant \emph{factor} per unit time.
\item \textbf{Decelerating phase:} growth slows as limiting factors appear (nutrients, space, accumulation of waste; in animals, the approach of adult size).
\item \textbf{Stationary phase:} growth stops; in plants the organs are mature, in animals adult size is reached and cell production only balances cell loss.
\end{itemize}
The whole period from the start of growth to the stationary phase is called the \cle{grand period of growth}.
\end{aretenir}

\courssub{Intermittent Growth: the Stepped Curve of Arthropods}

Not all organisms grow smoothly. A grasshopper or a cockroach is enclosed in a rigid \cle{exoskeleton} (cuticle) made of chitin, which cannot stretch. Such an animal can only increase in size by shedding its cuticle at intervals.

\begin{definition}[Ecdysis]
\cle{Ecdysis} (moulting) is the periodic \textbf{shedding of the rigid exoskeleton} of an arthropod, allowing a rapid increase in body size before the new, larger cuticle hardens. The stage between two successive moults is called an \cle{instar}.
\end{definition}

Consequently, the growth curve of an insect is \textbf{stepped (intermittent)}: long plateaux where length hardly changes (the cuticle restricts growth), separated by sudden vertical jumps at each moult, when the soft new cuticle allows rapid expansion by uptake of air or water.

\begin{popfigure}
\begin{center}
\begin{tikzpicture}
\begin{axis}[
width=11cm, height=6.5cm,
xlabel={Time (days)}, ylabel={Body length (mm)},
xmin=0, xmax=50, ymin=0, ymax=34,
xtick={0,10,20,30,40,50},
axis lines=left, thick]
\addplot[poporange, very thick] coordinates {(0,4) (10,4) (10,9) (20,9) (20,16) (30,16) (30,24) (40,24) (40,30) (50,30)};
\pgfplotsinvokeforeach{10,20,30,40}{
\draw[->, poppurple, thick] (axis cs:#1,2) -- (axis cs:#1,6);
\node[poppurple, font=\scriptsize, rotate=90, anchor=south] at (axis cs:#1,1) {Moult #1};
}
\node[popdark, font=\small, align=center] at (axis cs:25,27) {instars (plateaux)};
\end{axis}
\end{tikzpicture}
\end{center}
\legende{Stepped growth curve of an insect: each vertical jump corresponds to a moult (ecdysis); the horizontal plateaux are the instars.}
\end{popfigure}

\begin{attention}[Exam tip --- stepped versus smooth curves]
A frequent GCE question asks: \emph{``Explain why the growth curve of an insect is stepped while that of a human is smooth.''} The key idea: the insect's \textbf{rigid, inextensible exoskeleton} prevents continuous expansion, so size can only increase suddenly at each \textbf{moult}; the human skeleton is \textbf{internal} and grows continuously with the body, so growth is smooth and continuous.
\end{attention}

\courssub{The Human Growth Curve}

Human growth, plotted as total height against age, is roughly sigmoid but with a special feature: \textbf{two periods of rapid growth} --- in \emph{infancy} (the first two years) and at \emph{puberty} (the \cle{adolescent growth spurt}). This is seen most clearly on a \textbf{height-velocity curve}, which plots the \emph{rate} of growth (cm per year) against age.

\begin{popfigure}
\begin{center}
\begin{tikzpicture}
\begin{axis}[
width=11.5cm, height=7cm,
xlabel={Age (years)}, ylabel={Height velocity (cm/year)},
xmin=0, xmax=19, ymin=0, ymax=14,
xtick={0,2,4,6,8,10,12,14,16,18},
axis lines=left, thick,
legend style={at={(0.97,0.97)}, anchor=north east, font=\small}]
\addplot[popblue, very thick, smooth] coordinates {(0,12) (1,9) (2,7) (4,6) (6,5.4) (8,5) (10,4.8) (11.5,5) (13,8.5) (14,9.5) (15,6.5) (16,3.5) (17,1.5) (18,0.5) (19,0.2)};
\addlegendentry{Boys}
\addplot[poppink, very thick, smooth] coordinates {(0,12) (1,9) (2,7) (4,6) (6,5.4) (8,5.2) (9.5,5.2) (11,7.8) (12,8.2) (13,5.5) (14,3) (15,1.2) (16,0.4) (17,0.2) (19,0.1)};
\addlegendentry{Girls}
\draw[<->, popdark, thick] (axis cs:11.2,10.5) -- (axis cs:14.2,10.5);
\node[popdark, font=\small] at (axis cs:12.7,11.3) {pubertal spurt};
\end{axis}
\end{tikzpicture}
\end{center}
\legende{Height-velocity curves for boys and girls. Note the high rate in infancy and the adolescent growth spurt, which occurs about two years earlier in girls.}
\end{popfigure}

\begin{aretenir}[Features of human growth]
\begin{itemize}
\item Growth is fastest in \textbf{infancy}, then slows during childhood.
\item A marked \textbf{adolescent growth spurt} occurs at puberty: earlier in \textbf{girls} (about 11--13 years) than in \textbf{boys} (about 13--15 years), which is why girls of 12 are often taller than boys of the same age.
\item Growth stops when the \textbf{epiphyseal (growth) plates} of the long bones fuse, under the action of sex hormones.
\item Body parts grow at different rates --- \cle{allometric growth}: the head is proportionally large in a newborn (about a quarter of body length) but grows relatively little afterwards, while the legs grow greatly, changing the body proportions.
\end{itemize}
\end{aretenir}

\courssub{Growth Patterns: Determinate and Indeterminate}

\begin{definition}[Determinate and indeterminate growth]
\cle{Determinate growth} is growth that \textbf{stops at a characteristic adult size} (or when the organ matures): it is shown by most animals, by leaves, flowers and fruits, and by annual plants such as maize and cowpea. \cle{Indeterminate growth} continues \textbf{throughout life}, because permanent growth zones (meristems, growth plates) remain active: it is shown by perennial plants (trees), fish, and reptiles.
\end{definition}

In plants, two kinds of growth are distinguished:
\begin{itemize}
\item \cle{Primary growth}: \emph{elongation} of roots and shoots from \textbf{apical meristems} (at the tips); it produces the sigmoid curve of a seedling.
\item \cle{Secondary growth}: \emph{increase in girth} (thickening) of stems and roots in woody dicotyledons, due to the \textbf{vascular cambium} producing secondary xylem (wood) and secondary phloem, and the cork cambium producing bark. This is why an old \emph{iroko} or baobab tree keeps thickening year after year.
\end{itemize}

\courssub{Hormonal Control of Growth in Humans}

Growth is coordinated by \cle{hormones}, chemical messengers secreted by endocrine glands and carried in the blood to their target organs.

\begin{center}
\begin{tabularx}{\linewidth}{|l|l|X|}
\hline
\rowcolor{popblueL} \textbf{Hormone} & \textbf{Gland} & \textbf{Role in growth} \\
\hline
Growth hormone (GH) & Anterior pituitary & Stimulates protein synthesis, cell division and elongation of long bones (via the epiphyseal plates) throughout childhood \\
\hline
Thyroxine & Thyroid & Controls the basal metabolic rate; essential for normal growth and for development of the brain and skeleton \\
\hline
Testosterone / Oestrogen & Testes / Ovaries & Trigger the pubertal growth spurt and development of secondary sexual characteristics; later cause fusion of the epiphyseal plates, ending growth \\
\hline
Insulin & Pancreas (\emph{islets of Langerhans}) & Promotes uptake of glucose and protein synthesis; needed for normal tissue growth \\
\hline
\end{tabularx}
\end{center}

\begin{propriete}[Disorders of growth hormone]
The effect of growth hormone depends on the \textbf{age} at which the imbalance occurs (this is examinable as applied knowledge, not as a formal proof):
\begin{itemize}
\item \textbf{Deficiency before puberty} $\rightarrow$ \cle{dwarfism} (pituitary dwarf): very short stature but normal body proportions and normal intelligence.
\item \textbf{Excess before puberty} $\rightarrow$ \cle{gigantism}: abnormal overall tallness, because the epiphyseal plates are still open.
\item \textbf{Excess in adults} $\rightarrow$ \cle{acromegaly}: the long bones can no longer lengthen, so only the bones of the hands, feet, jaw and brow thicken, giving coarse features.
\end{itemize}
\end{propriete}

\begin{savaistu}[Thyroxine and cretinism]
Severe thyroxine deficiency in childhood (for example where dietary iodine is lacking) causes \emph{cretinism}: stunted growth together with mental retardation. Iodised salt, distributed in Cameroon precisely to combat iodine-deficiency disorders, prevents this condition.
\end{savaistu}

\courssub{Hormonal Control of Growth in Plants (Brief)}

Plant growth is controlled by \cle{plant growth substances} (phytohormones), produced in small quantities and transported to their sites of action.

\begin{itemize}
\item \cle{Auxins} (e.g. indole-acetic acid, IAA), produced in shoot and root tips, promote \textbf{cell elongation} by softening cell walls; they also mediate phototropism and apical dominance.
\item \cle{Gibberellins} promote \textbf{stem elongation} (internode lengthening), stimulate seed germination and can induce flowering; dwarf varieties of maize and rice often lack active gibberellin.
\item \cle{Cytokinins}, produced mainly in roots, promote \textbf{cell division} (cytokinesis) and delay ageing (senescence) of leaves.
\end{itemize}

Auxins and cytokinins act together: the \emph{balance} between them determines whether cultured plant cells form shoots or roots --- the principle behind \emph{in vitro} propagation of plantain and banana plantlets used by Cameroonian farmers.

\begin{exoresolu}[Interpreting a growth experiment]
The dry mass of samples of maize seedlings was measured every week after germination:

\begin{center}
\begin{tabular}{|l|ccccccc|}
\hline
\rowcolor{popblueL} Week & 1 & 2 & 3 & 4 & 5 & 6 & 7 \\
\hline
Mean dry mass (g) & 0.4 & 0.6 & 1.1 & 2.0 & 3.0 & 3.6 & 3.8 \\
\hline
\end{tabular}
\end{center}

\textbf{Questions:} (a) Why was dry mass, not fresh mass, used? (b) Between which weeks was the growth rate greatest? (c) Identify the phases of growth shown and explain the slowing after week 5.
\tcblower
\textbf{Solution.}
\begin{enumerate}
\item[(a)] Dry mass measures the \textbf{organic matter actually synthesised} by the seedlings; fresh mass would fluctuate with water uptake and loss and would not prove true growth.
\item[(b)] Growth rate $= \dfrac{\Delta m}{\Delta t}$. Weekly increases: 0.2, 0.5, 0.9, 1.0, 0.6, 0.2 g. The greatest increase is \textbf{between weeks 4 and 5} (\SI{1.0}{g/week}).
\item[(c)] Weeks 1--2: \textbf{lag phase} (slow growth, seedlings establishing). Weeks 2--5: \textbf{exponential (log) phase} (rapid, accelerating growth). Weeks 5--7: \textbf{decelerating then stationary phase} --- growth slows and nearly stops as the plants approach maturity and limiting factors (nutrients, light, space) take effect. The curve is therefore \textbf{sigmoid}.
\end{enumerate}
\end{exoresolu}

\begin{exercice}[Growth curves and hormones]
\begin{enumerate}
\item Define growth and state the three cellular processes that bring it about.
\item Sketch and label the sigmoid growth curve, naming its four phases.
\item Explain, with a labelled sketch, why the growth curve of a locust is stepped.
\item A boy receives injections of growth hormone from age 8. Predict the consequence if the treatment is continued in excessive doses into adulthood, and name the resulting condition.
\item State one role of each of the following in plant growth: auxin, gibberellin, cytokinin.
\end{enumerate}
\end{exercice}

\begin{corrige}[Exercise --- Growth curves and hormones]
\begin{enumerate}
\item Growth is the permanent, irreversible increase in size and dry mass of an organism, brought about by cell division, cell enlargement and cell differentiation.
\item S-shaped curve with: lag phase (slow start), exponential/log phase (steepest), decelerating phase, stationary phase (plateau at adult size).
\item The locust's rigid chitinous exoskeleton cannot stretch, so length stays constant during each instar (plateau); at each moult (ecdysis) the old cuticle is shed and the body expands rapidly before the new cuticle hardens (vertical jump). The curve is therefore a series of steps.
\item Excess GH before puberty causes gigantism; if excess continues after the epiphyseal plates have fused, the long bones cannot lengthen but the bones of the hands, feet and face thicken --- this is \textbf{acromegaly}.
\item Auxin: cell elongation (and phototropism); gibberellin: stem/internode elongation and germination; cytokinin: cell division and delay of senescence.
\end{enumerate}
\end{corrige}

\begin{aretenir}[Chapter summary --- Growth]
\begin{itemize}
\item Growth = irreversible increase in size and \textbf{dry mass} (division, enlargement, differentiation).
\item The sigmoid curve (lag, log, decelerating, stationary) describes most organisms; the grand period of growth covers all four phases.
\item Arthropods grow in \textbf{steps} because of moulting (ecdysis); humans grow smoothly with spurts in infancy and at puberty.
\item Determinate growth stops at maturity; indeterminate growth continues for life (trees, fish).
\item Human growth is controlled by GH, thyroxine, sex hormones and insulin; plant growth by auxins, gibberellins and cytokinins.
\end{itemize}
\end{aretenir}

\newpage

\coursec{Integration activity}

\courssub{Aims of the Activity}
\begin{aretenir}[Learning Objectives]
This integration activity consolidates knowledge from the entire Reproduction chapter. By completing it, you will:
\begin{itemize}
    \item Demonstrate understanding of \cle{sexual reproduction} in humans and flowering plants, including hormonal control, gametogenesis, fertilisation and early development.
    \item Apply concepts of \cle{asexual reproduction} and \cle{vegetative propagation} to agricultural practice.
    \item Design a controlled \cle{germination} experiment and interpret \cle{growth curves} (sigmoid and pubertal spurt).
    \item Evaluate \cle{contraceptive methods} in relation to the menstrual cycle and make evidence‑based recommendations.
    \item Communicate scientific reasoning clearly in a written advisory report and an oral presentation.
\end{itemize}
\end{aretenir}

\courssub{The Mfoundi Family Scenario}
\begin{experience}[Case Study Context]
The Mfoundi family lives on the outskirts of Yaoundé, Cameroon. Mr. and Mrs. Mfoundi have three children aged 14, 10 and 6. They wish to space their next child by at least three years for health and economic reasons. Mrs. Mfoundi is 32 years old, has a regular 28‑day menstrual cycle, and has no medical contraindications to hormonal contraception. The family also runs a small diversified farm of 0.5 ha where they cultivate:
\begin{itemize}
    \item \textbf{Plantain (Musa spp.)} – they want to multiply an improved high‑yielding cultivar using suckers.
    \item \textbf{Mango (Mangifera indica)} – they plan to graft scions of a late‑fruiting variety onto established rootstocks.
    \item \textbf{Maize (Zea mays)} – they have just harvested seeds from their best plants and want to test germination before the next planting season.
\end{itemize}
Their 14‑year‑old son, Pierre, is in the middle of puberty. The family has recorded his height every six months since age 10. They also measured the height of their maize crop weekly from sowing to harvest (12 weeks). The data are summarised below.

\begin{center}
\begin{tabularx}{\linewidth}{|c|X|X|}\hline
\textbf{Week after sowing} & \textbf{Maize height (cm)} & \textbf{Notes} \\ \hline
1 & 2.1 & Emergence \\
2 & 4.8 & \\
3 & 9.5 & \\
4 & 16.2 & \\
5 & 24.7 & \\
6 & 34.1 & \\
7 & 43.0 & \\
8 & 50.5 & \\
9 & 55.8 & \\
10 & 58.9 & \\
11 & 60.2 & \\
12 & 60.5 & Harvest \\ \hline
\end{tabularx}
\end{center}

\begin{center}
\begin{tabularx}{\linewidth}{|c|X|X|}\hline
\textbf{Age (years)} & \textbf{Pierre’s height (cm)} & \textbf{Notes} \\ \hline
10.0 & 138.2 & Pre‑pubertal \\
10.5 & 140.5 & \\
11.0 & 143.0 & \\
11.5 & 146.2 & \\
12.0 & 150.8 & Onset of growth spurt \\
12.5 & 156.5 & \\
13.0 & 163.0 & Peak height velocity \\
13.5 & 168.2 & \\
14.0 & 171.5 & Deceleration \\ \hline
\end{tabularx}
\end{center}

The family asks you, as a team of Upper Sixth biology students, to act as their advisers. You must produce a one‑page advisory report (maximum 300 words) and defend it in a 5‑minute class presentation.
\end{experience}

\courssub{Tasks}
\begin{enumerate}
    \item \textbf{Contraceptive choice:} Select the most appropriate contraceptive method for Mrs. Mfoundi. Justify your choice by linking its mode of action to the hormonal events of the menstrual cycle (FSH, LH, oestrogen, progesterone). Discuss effectiveness, reversibility, side‑effects and cultural acceptability in the Cameroonian context.
    \item \textbf{Vegetative propagation plan:} Design a step‑by‑step plan to multiply the improved plantain cultivar using suckers and to graft the mango scions onto rootstocks. Include timing, selection criteria, sanitation, and after‑care. Explain why these methods are considered asexual reproduction and how they preserve the genotype.
    \item \textbf{Germination experiment:} Propose a controlled experiment to test the germination capacity of the maize seeds. State the hypothesis, variables (independent, dependent, controlled), experimental set‑up (number of replicates, substrate, temperature, moisture, light), data to be collected, and statistical analysis (e.g., germination percentage, mean germination time).
    \item \textbf{Growth curve analysis:}
    \begin{enumerate}
        \item Plot the maize height data and Pierre’s height data on separate graphs (time on x‑axis, height on y‑axis). Identify the type of growth curve for each (sigmoid vs. pubertal spurt).
        \item For the maize curve, label the lag, exponential, and stationary phases. Calculate the approximate maximum growth rate (cm week\(^{-1}\)) during the exponential phase.
        \item For Pierre’s curve, determine the age at peak height velocity (PHV) and the peak height velocity value (cm year\(^{-1}\)). Relate these to hormonal changes (testosterone, growth hormone, IGF‑1).
    \end{enumerate}
    \item \textbf{Integration and report writing:} Synthesize your answers into a concise advisory report (≤300 words) addressing all four points. Prepare a 5‑minute presentation highlighting the scientific reasoning and practical recommendations.
\end{enumerate}

\courssub{Model Answers}

\begin{exemplebox}[Task 1: Contraceptive Choice for Mrs. Mfoundi]
\textbf{Recommended method:} Combined Oral Contraceptive Pill (COCP) containing ethinylestradiol (oestrogen) and levonorgestrel (progestogen).

\textbf{Mode of action linked to menstrual cycle:}
\begin{itemize}
    \item The COCP maintains high circulating levels of synthetic oestrogen and progesterone.
    \item These exert \textbf{negative feedback} on the hypothalamus and anterior pituitary, suppressing the secretion of \cle{GnRH}, \cle{FSH} and \cle{LH}.
    \item Crucially, the \textbf{mid‑cycle LH surge is prevented}; without it, \cle{ovulation} does not occur.
    \item Additional effects: progestogen thickens \cle{cervical mucus} (blocking sperm penetration) and alters the \cle{endometrium} (making it unsuitable for implantation).
\end{itemize}

\textbf{Effectiveness:} Perfect use ≈ 99.7\% (Pearl Index ~0.3); typical use ≈ 91\% (Pearl Index ~9).
\textbf{Reversibility:} Full fertility returns within 1–3 cycles after stopping.
\textbf{Side‑effects:} Nausea, breast tenderness, weight change, mood changes (usually transient); rare but serious risk of thromboembolism (counselling required).
\textbf{Cameroonian context:} Widely available in public health centres and pharmacies at subsidised cost; culturally accepted among educated urban couples; daily adherence counselling is essential.
\textbf{Alternative:} If daily adherence is a concern, Long‑Acting Reversible Contraceptives (LARCs) such as the subdermal implant (etonogestrel) or hormonal IUD (levonorgestrel) offer >99\% effectiveness with user‑independent action.
\end{exemplebox}

\begin{exemplebox}[Task 2: Vegetative Propagation Plan]
\textbf{A. Plantain (Musa spp.) – Multiplication by Sword Suckers}
\begin{methode}[Step‑by‑step]
\begin{enumerate}
\item \textbf{Timing:} Start of rainy season (March–April in Yaoundé area) for optimal establishment.
\item \textbf{Mother plant selection:} Healthy, disease‑free (no Banana Bunchy Top Virus, no nematode damage), 6–8 months after planting, bearing the desired high‑yielding cultivar traits.
\item \textbf{Sucker selection:} Choose \textbf{sword suckers} (narrow, lanceolate leaves, strong conical corm) 30–50 cm tall. Avoid water suckers (broad leaves, weak corm).
\item \textbf{Detachment:} Use a clean, sharp machete. Cut the sucker from the mother corm, ensuring a portion of corm (5–10 cm) and roots remain attached.
\item \textbf{Sanitation:} Trim leaves to 2–3 cm from the corm to reduce transpiration. Dip corm in 1\% sodium hypochlorite (bleach) solution for 5 minutes, then rinse in clean water. Optional: hot water treatment (50 °C for 20 min) against nematodes.
\item \textbf{Planting:} Prepare holes 30 cm × 30 cm × 30 cm at 3 m × 3 m spacing. Add 5–10 kg compost or well‑decomposed manure per hole. Plant sucker upright, cover corm with soil, firm gently.
\item \textbf{After‑care:} Mulch heavily with dry grass or banana leaves. Water regularly until established. Apply mulch and organic fertiliser every 3–4 months. Monitor for pests (weevils, nematodes) and diseases.
\end{enumerate}
\end{methode}
\textbf{Genetic basis:} Suckers arise from meristematic tissue of the mother corm → \cle{mitotic divisions} only → \textbf{genetically identical clones}. The improved cultivar’s genotype is preserved exactly.

\textbf{B. Mango (Mangifera indica) – Grafting (Cleft Graft)}
\begin{methode}[Step‑by‑step]
\begin{enumerate}
\item \textbf{Rootstock:} 6–12 month‑old seedlings of a local, hardy, well‑adapted variety (e.g., ‘Kensington’ or local ‘Mangoro’), pencil‑thickness (≈1 cm diameter), actively growing.
\item \textbf{Scion collection:} During dry season dormancy (December–January). Select 15–20 cm mature shoots with 3–4 plump buds from the desired late‑fruiting variety. Wrap in damp cloth, keep cool.
\item \textbf{Grafting operation:} Under shade. Decapitate rootstock 20–30 cm above soil. Make a vertical cleft (3–4 cm deep) in the rootstock stem. Cut scion base into a wedge. Insert scion into cleft, aligning \cle{cambium} layers on at least one side. Secure tightly with grafting tape (polythene).
\item \textbf{Sealing:} Apply grafting wax or petroleum jelly over the union and cut surfaces to prevent desiccation.
\item \textbf{After‑care:} Place in high‑humidity propagator (or cover with clear polythene bag) for 3–4 weeks. Ventilate gradually. Remove any rootstock sprouts below the graft. Harden off by increasing exposure to ambient conditions over 2 weeks.
\item \textbf{Field planting:} Transplant at onset of rains (March), 10 m × 10 m spacing.
\end{enumerate}
\end{methode}
\textbf{Genetic basis:} Grafting joins two genotypes: the \textbf{scion} (fruiting variety) and the \textbf{rootstock} (adaptability). The union is \cle{asexual} – no meiosis, no gamete fusion. The scion’s genotype is preserved, ensuring true‑to‑type fruit quality and late‑fruiting trait.
\end{exemplebox}

\begin{exemplebox}[Task 3: Controlled Germination Experiment]
\begin{methode}[Experimental Design]
\begin{enumerate}
\item \textbf{Hypothesis:} Maize seeds selected from the family’s best plants have a germination percentage ≥ 90\% and a mean germination time ≤ 4 days under optimal conditions.
\item \textbf{Variables:}
\begin{itemize}
    \item \textbf{Independent:} Seed source (Family’s saved seeds vs. Certified commercial seeds of same variety – control).
    \item \textbf{Dependent:} Germination percentage (\%), Mean Germination Time (MGT, days).
    \item \textbf{Controlled:} Substrate (sterile river sand or moist filter paper), Temperature (25 °C ± 1 °C), Moisture (field capacity / constant moisture), Light (continuous darkness – maize germinates in dark), Seed depth (3 cm in sand / between papers), Seed viability (use seeds from same harvest year).
\end{itemize}
\item \textbf{Set‑up:} \textbf{Randomised Complete Block Design}. 4 replicates (blocks) × 2 treatments = 8 experimental units. Each unit: 50 seeds. Seeds placed in Petri dishes (filter paper) or sand trays. Incubate in a germination cabinet or dark room at 25 °C.
\item \textbf{Data collection:} Count germinated seeds (radicle ≥ 2 mm) daily for 10 days. Record number germinated per day (\(n_i\)) for each replicate.
\item \textbf{Statistical analysis:}
\begin{itemize}
    \item \textbf{Germination percentage:} \(GP = \frac{\text{total germinated}}{\text{total seeds}} \times 100\). Compare treatments using \textbf{Chi‑square test} (\(\chi^2\)) on contingency table (germinated vs. non‑germinated).
    \item \textbf{Mean Germination Time (MGT):} \(MGT = \frac{\sum (n_i \times d_i)}{\sum n_i}\) where \(d_i\) = day of count. Compare MGT between treatments using \textbf{unpaired t‑test} (check normality, equal variance) or Mann‑Whitney U test. Significance level \(\alpha = 0.05\).
\end{itemize}
\end{enumerate}
\end{methode}
\end{exemplebox}

\begin{exemplebox}[Task 4: Growth Curve Analysis]
\textbf{A. Maize Height Growth}
\begin{aretenir}[Sigmoid Growth Curve]
The maize data produce a classic \textbf{sigmoid (S‑shaped) curve} typical of annual plants.
\begin{itemize}
    \item \textbf{Lag phase (Weeks 1–2):} Slow growth; establishment of root system, limited photosynthetic area.
    \item \textbf{Exponential (log) phase (Weeks 3–7):} Rapid cell division and elongation; maximum photosynthetic rate; steepest slope.
    \item \textbf{Stationary phase (Weeks 8–12):} Growth slows and stops; resources diverted to grain filling; senescence begins.
\end{itemize}
\textbf{Maximum Growth Rate Calculation:}
Use the steepest interval (Weeks 3–7):
\[
\text{Rate} = \frac{\Delta \text{height}}{\Delta \text{time}} = \frac{43.0 - 9.5}{7 - 3} = \frac{33.5}{4} = 8.375 \approx \mathbf{8.4\ cm\ week^{-1}}
\]
\end{aretenir}

\textbf{B. Pierre’s Height Growth}
\begin{aretenir}[Pubertal Growth Spurt]
The data show a \textbf{pubertal growth spurt} – a rapid increase in height velocity during adolescence.
\begin{itemize}
    \item \textbf{Peak Height Velocity (PHV) age:} \textbf{13.0 years} (height = 163.0 cm).
    \item \textbf{PHV value:} Calculate over the interval centred on PHV (12.5–13.5 years):
    \[
    PHV = \frac{168.2 - 156.5}{13.5 - 12.5} = \frac{11.7}{1.0} = \mathbf{11.7\ cm\ year^{-1}}
    \]
\end{itemize}
\textbf{Hormonal basis:} Rising \cle{testosterone} (Leydig cells) stimulates \cle{GH} secretion from anterior pituitary → GH stimulates liver to produce \cle{IGF‑1} (Insulin‑like Growth Factor 1). GH and IGF‑1 act synergistically on \cle{epiphyseal plates} (growth plates) to increase chondrocyte proliferation and matrix synthesis → longitudinal bone growth. Testosterone also directly stimulates growth plate maturation (eventually causing epiphyseal closure).
\end{aretenir}
\end{exemplebox}

\begin{exemplebox}[Task 5: Advisory Report (Model – 298 words)]
\textit{To the Mfoundi Family,}

We recommend the \textbf{combined oral contraceptive pill (COCP)} for child spacing. It maintains high oestrogen/progestogen levels, suppressing the hypothalamic‑pituitary axis and preventing the mid‑cycle LH surge; thus ovulation is blocked. Additional effects: thickened cervical mucus and altered endometrium. Perfect‑use effectiveness is 99.7\%; typical use 91\%. It is fully reversible (fertility returns in 1–3 cycles), affordable, and widely available in Yaoundé health centres. Cultural acceptability is high; we advise daily intake at the same time and counselling on transient side‑effects (nausea, weight change) and rare thromboembolism risk. If adherence is difficult, a subdermal implant or hormonal IUD are excellent alternatives.

For the farm: \textbf{Plantain} – multiply using \textbf{sword suckers} (30–50 cm) from healthy mother plants at the start of the rains. Sanitise corms in 1\% bleach, plant at 3 m × 3 m with compost and mulch. This \cle{clonal propagation} preserves the improved genotype exactly. \textbf{Mango} – \textbf{cleft‑graft} scions of the late‑fruiting variety onto 6–12 month local rootstocks during the dry season. Align cambium, seal with wax, maintain high humidity until union forms (3–4 weeks). Grafting combines the scion’s fruiting genotype with the rootstock’s hardiness – an asexual method ensuring true‑to‑type fruit.

\textbf{Maize germination test:} Four replicates of 50 seeds each on sterile moist sand at 25 °C in darkness. Count daily for 10 days. Calculate germination percentage and Mean Germination Time \(MGT = \frac{\sum n_i d_i}{\sum n_i}\). Compare with certified seeds using \(\chi^2\) and t‑tests (\(\alpha=0.05\)). Expect ≥90\% germination if seeds are viable.

\textbf{Growth analysis:} Maize follows a \textbf{sigmoid curve}; maximum growth rate ≈ 8.4 cm week\(^{-1}\) (weeks 3–7). Pierre shows a \textbf{pubertal spurt}; PHV ≈ 11.7 cm year\(^{-1}\) at age 13, driven by testosterone, GH and IGF‑1 acting on epiphyseal plates. Both curves illustrate distinct biological growth patterns: deterministic in annual plants, hormone‑mediated in adolescents.

Implementing these recommendations will support family health, farm productivity, and informed monitoring of Pierre’s development. We are available for follow‑up guidance.

\textit{Upper Sixth Biology Advisory Team}
\end{exemplebox}

\begin{attention}[Common Pitfalls in the Integration Activity]
\begin{itemize}
\item Do not confuse the contraceptive mechanism: the COCP prevents the \textbf{LH surge} (not the FSH surge alone) via negative feedback.
\item Vegetative propagation must emphasise \textbf{genetic identity (clones)} – do not describe it as sexual reproduction or imply genetic recombination.
\item In the germination experiment, \textbf{control all abiotic factors} (temperature, moisture, light, substrate); omitting any invalidates the test.
\item When calculating growth rates, use the \textbf{steepest part of the curve}: exponential phase for maize (weeks 3–7), PHV interval (12.5–13.5 years) for Pierre.
\item The advisory report must be \textbf{concise (≤300 words)} and address \textbf{all four tasks}; marks are awarded for integration and application, not just isolated facts.
\item For mango grafting, specify \textbf{cambial alignment} – this is the key to a successful vascular union.
\item In growth analysis, distinguish between \textbf{absolute growth rate} (cm/time) and \textbf{relative growth rate}; the question asks for absolute rate.
\end{itemize}
\end{attention}

\newpage

\coursec{Methods and worked examples}

\courssub{Skill 1: Identifying the Type of Asexual Reproduction}

\begin{methode}[Recognising asexual reproduction mechanisms]
When faced with a description, a diagram or a micrograph of an organism reproducing asexually, proceed as follows.
\begin{enumerate}
\item \textbf{Count the parents.} If only \cle{one parent} is involved and no gametes are formed or fused, the reproduction is asexual.
\item \textbf{Look at what is produced.} Ask: is the offspring a whole new individual, an outgrowth, or a resistant cell?
\item \textbf{Match the observation to the mechanism:}
\begin{itemize}
\item \cle{Binary fission}: one cell splits into two equal daughter cells (Amoeba, Paramecium, bacteria).
\item \cle{Budding}: a small outgrowth (bud) forms on the parent, then detaches (Hydra, yeast).
\item \cle{Fragmentation / regeneration}: the body breaks into pieces, each regenerating a whole organism (planaria, starfish, Spirogyra).
\item \cle{Sporulation}: spores produced in a sporangium germinate into new individuals (Mucor, ferns, mosses).
\item \cle{Vegetative propagation}: new plants from vegetative organs — stems (runners, rhizomes, tubers, cuttings), roots or leaves — e.g. cassava and sugarcane cuttings, Irish potato tubers, ginger rhizomes, onion bulbs.
\item \cle{Parthenogenesis}: an unfertilised egg develops into an adult (aphids, drones in bees).
\end{itemize}
\item \textbf{State the genetic consequence}: offspring are \cle{genetically identical} (clones) to the parent, except for rare mutations.
\end{enumerate}
\textbf{Reflex:} ``One parent + no gamete fusion = asexual; offspring = clone.''
\end{methode}

\begin{exoresolu}[Identifying modes of asexual reproduction]
An Upper Sixth student observes the following in the school biology garden and laboratory: (a) a cassava stem cutting planted in the soil develops roots and shoots; (b) a slide of \emph{Mucor} shows black sporangia releasing spores; (c) a \emph{Hydra} with a small lateral outgrowth; (d) an \emph{Amoeba} constricting into two cells. For each case, name the type of asexual reproduction and justify your answer. State one advantage and one disadvantage of asexual reproduction.
\tcblower
\textbf{Solution.}
\begin{enumerate}
\item \textbf{(a) Cassava cutting:} \cle{Vegetative propagation} (propagation by stem cuttings). Justification: a new plant develops from a vegetative organ (stem), not from a seed or a gamete.
\item \textbf{(b) Mucor:} \cle{Sporulation}. Justification: spores are produced inside sporangia; each spore can germinate into a new mycelium.
\item \textbf{(c) Hydra:} \cle{Budding}. Justification: a small bud grows from the parent's body wall and will detach as an independent individual.
\item \textbf{(d) Amoeba:} \cle{Binary fission}. Justification: the single parent cell divides by mitosis into two genetically identical daughter cells.
\end{enumerate}
\textbf{Advantage:} rapid reproduction and colonisation (no mate needed, no energy spent on gametes); all offspring are well adapted to the parent's environment.\\
\textbf{Disadvantage:} no genetic variation, so the whole population may be wiped out by a single disease or environmental change (e.g. cassava mosaic disease spreading through infected cuttings).
\end{exoresolu}

\courssub{Skill 2: Accounting for Chromosome Numbers in Gamete Formation and Double Fertilization}

\begin{methode}[Tracking chromosome and DNA numbers in flowering plants]
\begin{enumerate}
\item \textbf{Fix the diploid number} $2n$ given in the question. Write it down before any calculation.
\item \textbf{Apply the key rule:} \cle{meiosis} halves the chromosome number ($2n \to n$); \cle{fertilization} restores it ($n + n \to 2n$).
\item \textbf{In the flower, remember the ploidy of each structure:}
\begin{itemize}
\item Pollen grain (male gametophyte): $n$; it contains a \cle{generative nucleus} ($n$) which divides by mitosis into \textbf{two male gamete nuclei} ($n$ each), plus a tube nucleus ($n$).
\item Embryo sac (female gametophyte): $n$; it contains the \cle{egg cell} ($n$) and two \cle{polar nuclei} ($n$ each, i.e. $2n$ when fused).
\end{itemize}
\item \textbf{Double fertilization:} one male nucleus ($n$) + egg ($n$) $\to$ \cle{zygote} ($2n$); the other male nucleus ($n$) + the two polar nuclei ($2n$) $\to$ \cle{endosperm nucleus} ($3n$, triploid).
\item \textbf{Check:} the zygote must be $2n$; the endosperm must be $3n$. If your answer differs, re-examine step 4.
\end{enumerate}
\textbf{Reflex:} ``Egg $n$, zygote $2n$, endosperm $3n$.''
\end{methode}

\begin{popfigure}
\begin{center}
\begin{tikzpicture}[scale=0.95, every node/.style={font=\small}]
% Pollen grain
\draw[thick, popblue] (0,0) circle (0.9);
\node[popblue] at (0,1.25) {Pollen grain ($n$)};
\fill[poporange] (-0.3,0.25) circle (0.14);
\fill[poporange] (0.3,0.25) circle (0.14);
\fill[popgreen] (0,-0.4) circle (0.14);
\node[align=center] at (0,-1.3) {2 male nuclei ($n$)\\ tube nucleus ($n$)};
% Pollen tube
\draw[thick, popblue] (0.75,-0.5) .. controls (1.6,-1.2) and (2.2,-2.2) .. (3.1,-3.0);
\draw[thick, popblue] (0.55,-0.75) .. controls (1.4,-1.5) and (2.0,-2.5) .. (2.85,-3.25);
\node[popblue, align=center] at (1.6,-2.6) {pollen\\ tube};
% Embryo sac
\draw[thick, poppurple] (5.2,-3.6) ellipse (1.5 and 2.1);
\node[poppurple] at (5.2,-1.2) {Embryo sac ($n$)};
\fill[popgreen] (4.4,-4.9) circle (0.16);
\node[align=left, anchor=west] at (1.1,-5.1) {egg cell ($n$)};
\fill[poporange] (5.0,-3.6) circle (0.14);
\fill[poporange] (5.45,-3.6) circle (0.14);
\node[align=center] at (6.9,-3.6) {2 polar\\ nuclei ($n$+$n$)};
% Arrows of double fertilization
\draw[-latex, thick, popdark] (3.1,-3.2) -- (4.15,-4.6);
\draw[-latex, thick, popdark] (3.2,-3.3) -- (4.6,-3.55);
% Products
\node[align=center, popdark] at (2.6,-6.6) {\textbf{zygote} ($2n$)\\ $\to$ embryo};
\node[align=center, popdark] at (7.6,-6.6) {\textbf{endosperm} ($3n$)\\ $\to$ food store};
\draw[-latex, thick, popdark] (4.4,-5.1) .. controls (3.8,-5.9) .. (3.1,-6.3);
\draw[-latex, thick, popdark] (5.3,-3.95) .. controls (6.4,-5.2) .. (7.2,-6.2);
\end{tikzpicture}
\end{center}
\legende{Double fertilization in a flowering plant: one male nucleus fuses with the egg cell (zygote, $2n$); the second fuses with the two polar nuclei (endosperm, $3n$).}
\end{popfigure}

\begin{exoresolu}[Chromosome numbers in maize reproduction]
Maize (\emph{Zea mays}) has a diploid number $2n = 20$. (a) State the number of chromosomes in: (i) a pollen mother cell; (ii) a pollen grain nucleus; (iii) the egg cell; (iv) the zygote; (v) the endosperm. (b) Explain how the endosperm comes to have its chromosome number. (c) How many chromosomes would the embryo of the seed contain?
\tcblower
\textbf{Solution.}\\
(a)
\begin{enumerate}
\item (i) \textbf{Pollen mother cell:} it is a diploid cell of the anther, so $2n = 20$ chromosomes.
\item (ii) \textbf{Pollen grain nucleus:} produced by meiosis of the pollen mother cell, so $n = 10$ chromosomes.
\item (iii) \textbf{Egg cell:} a haploid gamete, $n = 10$ chromosomes.
\item (iv) \textbf{Zygote:} egg ($n$) + male nucleus ($n$) $= 2n = 20$ chromosomes.
\item (v) \textbf{Endosperm:} triploid, $3n = 30$ chromosomes.
\end{enumerate}
(b) \textbf{Origin of the endosperm:} during \cle{double fertilization}, the second male gamete nucleus ($n = 10$) fuses with the two polar nuclei ($n + n = 20$) in the embryo sac. The fusion product therefore contains $10 + 20 = 30$ chromosomes ($3n$). It divides mitotically to form the endosperm, the nutritive tissue of the seed.\\
(c) \textbf{Embryo:} the embryo develops by repeated mitosis of the zygote, so every embryonic cell contains $2n = 20$ chromosomes.\\
\textbf{Check:} zygote $2n = 20$ \checkmark; endosperm $3n = 30$ \checkmark.
\end{exoresolu}

\courssub{Skill 3: Interpreting Germination Experiments and Calculating Germination Percentage}

\begin{methode}[Analysing a germination experiment]
\begin{enumerate}
\item \textbf{Identify the variables:} independent variable (the factor varied: water, oxygen, temperature, light), dependent variable (germination: number or percentage of seeds germinated), controlled variables (same seed batch, same number of seeds, same duration).
\item \textbf{Compute the germination percentage:}
\[ \text{Germination \%} = \frac{\text{number of seeds germinated}}{\text{total number of seeds sown}} \times 100 \]
\item \textbf{Compare set-ups:} a seed germinates only if \cle{water}, \cle{oxygen} and a \cle{suitable temperature} are all present. Light is needed only by some species (e.g. lettuce); most seeds germinate equally well in darkness.
\item \textbf{Draw the conclusion} by linking each result to the missing or present condition, and state the role of each factor (water: imbibition and enzyme activation; oxygen: aerobic respiration for energy; temperature: enzyme activity).
\item \textbf{Evaluate reliability:} large sample, replicates, control of other variables.
\end{enumerate}
\end{methode}

\begin{exoresolu}[Germination of bean seeds]
A student in Buea sets up four flasks, each with 50 bean seeds: Flask A — moist cotton wool, open to air, at $25\ ^\circ$C; Flask B — dry cotton wool, open to air, $25\ ^\circ$C; Flask C — seeds completely covered with boiled, cooled water, $25\ ^\circ$C; Flask D — moist cotton wool, open to air, at $4\ ^\circ$C in a refrigerator. After 5 days: A = 45 germinated, B = 0, C = 0, D = 2. (a) Calculate the germination percentage in each flask. (b) Explain the result of each flask. (c) State the purpose of boiling and cooling the water in flask C.
\tcblower
\textbf{Solution.}\\
(a) Germination \% $= \dfrac{\text{germinated}}{50}\times 100$:
\begin{itemize}
\item A: $\frac{45}{50}\times 100 = 90\ \%$
\item B: $\frac{0}{50}\times 100 = 0\ \%$
\item C: $0\ \%$
\item D: $\frac{2}{50}\times 100 = 4\ \%$
\end{itemize}
(b) \textbf{Flask A} (control): water, oxygen and a suitable temperature are all present, so enzymes are activated, food reserves are mobilised and respiration supplies energy — high germination ($90\ \%$). \textbf{Flask B}: no water, so no imbibition and no enzyme activation — no germination. \textbf{Flask C}: boiled water contains almost no dissolved oxygen; without oxygen, aerobic respiration cannot release the energy needed for growth — no germination. \textbf{Flask D}: the low temperature slows enzyme activity drastically, so germination is very poor ($4\ \%$).\\
(c) \textbf{Boiling} expels dissolved air (oxygen) from the water; \textbf{cooling} prevents the heat from killing the seeds, so that the only missing factor is oxygen. This makes flask C a valid test of the need for oxygen.\\
\textbf{Conclusion:} bean seeds need water, oxygen and a suitable temperature to germinate.
\end{exoresolu}

\courssub{Skill 4: Reading Hormone Graphs of the Menstrual Cycle}

\begin{methode}[Interpreting the menstrual cycle graph]
\begin{enumerate}
\item \textbf{Identify each curve} by its profile over the 28 days: \cle{FSH} rises early (follicle growth); \cle{oestrogen} rises to a peak just before mid-cycle; \cle{LH} shows a sharp \textbf{peak at about day 14} (this triggers ovulation); \cle{progesterone} dominates the second half (secreted by the corpus luteum).
\item \textbf{Locate ovulation:} it occurs just after the LH peak, around \cle{day 14}.
\item \textbf{Link hormones to events:} FSH $\to$ follicle development; oestrogen $\to$ repair and thickening of the endometrium + positive feedback causing the LH surge; LH $\to$ ovulation and formation of the corpus luteum; progesterone $\to$ maintenance of the endometrium and \textbf{negative feedback} inhibiting FSH and LH.
\item \textbf{Explain menstruation:} if no fertilization occurs, the corpus luteum degenerates, progesterone and oestrogen fall, the endometrium breaks down and is shed (days 1--5).
\item \textbf{Answer ``what if pregnant?'':} the corpus luteum is maintained (by HCG from the embryo), progesterone stays high, no menstruation.
\end{enumerate}
\textbf{Reflex:} ``LH peak = ovulation; progesterone fall = menstruation.''
\end{methode}

\begin{popfigure}
\begin{center}
\begin{tikzpicture}[scale=0.9]
\draw[-latex, thick] (0,0) -- (10.6,0) node[below left, font=\small]{days};
\draw[-latex, thick] (0,0) -- (0,4.4) node[above, font=\small]{blood level};
\foreach \x/\d in {0/0, 1.43/4, 2.5/7, 5/14, 7.5/21, 10/28}
  \draw (\x,0.06) -- (\x,-0.06) node[below, font=\footnotesize]{\d};
% FSH: early moderate rise
\draw[very thick, popgreen, smooth] plot coordinates {(0,0.9) (1.4,1.5) (2.6,1.2) (4.6,0.8) (5,1.1) (5.4,0.7) (7.5,0.5) (10,0.5)};
\node[popgreen, font=\small] at (1.7,1.85) {FSH};
% Oestrogen: peak day 12-13, smaller second rise
\draw[very thick, poporange, smooth] plot coordinates {(0,0.4) (2,0.7) (3.5,1.6) (4.4,2.9) (5,1.1) (6.5,1.5) (8,1.4) (10,0.5)};
\node[poporange, font=\small] at (3.6,3.15) {oestrogen};
% LH: sharp peak day 13-14
\draw[very thick, popblue, smooth] plot coordinates {(0,0.5) (2.5,0.6) (4.2,0.9) (4.8,3.9) (5.3,1.0) (7,0.6) (10,0.5)};
\node[popblue, font=\small] at (5.6,3.9) {LH};
% Progesterone: second half
\draw[very thick, poppurple, smooth] plot coordinates {(0,0.15) (5,0.2) (6,1.2) (7.3,2.4) (8.6,1.6) (10,0.25)};
\node[poppurple, font=\small] at (8.1,2.75) {progesterone};
% Ovulation marker
\draw[dashed, popdark] (5,0) -- (5,4.1);
\node[popdark, font=\footnotesize, align=center] at (5,4.35) {ovulation};
% Menstruation band
\fill[popink, opacity=0.12] (0,0) rectangle (1.43,4.1);
\node[font=\footnotesize, align=center] at (0.72,3.7) {menstr.};
\end{tikzpicture}
\end{center}
\legende{Typical profiles of FSH, LH, oestrogen and progesterone during a 28-day menstrual cycle (schematic).}
\end{popfigure}

\begin{exoresolu}[Analysing a hormone graph]
The graph of a woman's cycle shows: FSH rising from day 1 to day 6; oestrogen peaking on day 12; a sharp LH peak on day 13--14; progesterone rising from day 15 to day 22 then falling. (a) On which day did ovulation most likely occur? Justify. (b) Which structure secretes progesterone after day 15, and from what does it develop? (c) Explain why progesterone falls after day 22. (d) Explain how a pregnancy would change the progesterone curve. (e) Explain how the combined contraceptive pill exploits this cycle.
\tcblower
\textbf{Solution.}\\
(a) \textbf{Ovulation: about day 14.} The LH surge (peak on day 13--14) is the direct trigger of ovulation; the egg is released roughly a day after the LH peak, i.e. around day 14.\\
(b) Progesterone is secreted by the \cle{corpus luteum}, which develops from the \textbf{ruptured Graafian follicle} left in the ovary after ovulation, under the influence of LH.\\
(c) \textbf{Fall of progesterone:} in the absence of fertilization, the corpus luteum degenerates after about 10--12 days (it becomes the corpus albicans). With no corpus luteum, progesterone (and oestrogen) secretion collapses; the endometrium is no longer maintained and breaks down — menstruation begins.\\
(d) \textbf{If pregnant:} the embryo implants and secretes \cle{HCG} (human chorionic gonadotrophin), which maintains the corpus luteum. Progesterone therefore \textbf{remains high} instead of falling; the endometrium is preserved and menstruation does not occur. Later, the placenta takes over progesterone production.\\
(e) \textbf{The combined pill} contains synthetic oestrogen and progesterone. These hormones exert \textbf{negative feedback} on the hypothalamus and pituitary, inhibiting FSH and LH secretion. Without FSH, follicles do not mature; without the LH surge, \textbf{ovulation does not occur}, so no egg is available for fertilization.
\end{exoresolu}

\courssub{Skill 5: Comparing Spermatogenesis and Oogenesis}

\begin{methode}[Constructing a comparison of gametogenesis]
\begin{enumerate}
\item \textbf{Fix the framework:} both processes involve \cle{mitosis} of germ cells, growth, \cle{meiosis}, and maturation/differentiation.
\item \textbf{Compare point by point} (never mix the two columns of your reasoning): site, timing, number of functional gametes per meiosis, size of gametes, continuity of the process.
\item \textbf{Key figures to remember:}
\begin{itemize}
\item Spermatogenesis: 1 primary spermatocyte $\to$ \textbf{4 functional spermatozoa}; continuous from puberty; occurs in the seminiferous tubules; controlled by FSH and testosterone (LH acts on Leydig cells).
\item Oogenesis: 1 primary oocyte $\to$ \textbf{1 functional ovum} + polar bodies; begins before birth, meiosis I completed only just before ovulation, meiosis II completed only at fertilization; controlled by FSH, LH, oestrogen, progesterone.
\end{itemize}
\item \textbf{Hormonal control:} name the gland (hypothalamus/pituitary), the hormone, the target and the effect. GnRH $\to$ FSH + LH.
\end{enumerate}
\end{methode}

\begin{center}
\begin{tabularx}{\linewidth}{|l|X|X|}
\hline
\rowcolor{popblueL} \textbf{Feature} & \textbf{Spermatogenesis} & \textbf{Oogenesis} \\
\hline
Site & Seminiferous tubules of testes & Ovaries \\
\hline
Timing & Continuous from puberty & Begins before birth; meiosis completed only at ovulation/fertilization \\
\hline
Functional gametes per meiosis & 4 spermatozoa & 1 ovum + polar bodies \\
\hline
Cytokinesis & Equal & Unequal (cytoplasm kept in one large cell) \\
\hline
Gamete & Small, motile, no reserves & Large, non-motile, rich in yolk \\
\hline
Main hormones & FSH, LH, testosterone & FSH, LH, oestrogen, progesterone \\
\hline
\end{tabularx}
\end{center}

\begin{exoresolu}[Gamete yield and hormonal control]
(a) A man produces sperm continuously from puberty. Starting from 100 primary spermatocytes, how many functional spermatozoa are produced? Starting from 100 primary oocytes, how many functional ova are produced? Explain the difference. (b) Describe the hormonal control of the testes. (c) A tumour destroys a man's Leydig (interstitial) cells. Predict and explain the consequences.
\tcblower
\textbf{Solution.}\\
(a) Each \textbf{primary spermatocyte} undergoes meiosis I (two secondary spermatocytes) then meiosis II, giving \textbf{4 spermatids}, each of which differentiates into a functional spermatozoon. So 100 primary spermatocytes $\to 100 \times 4 = \mathbf{400}$ \textbf{spermatozoa}.\\
Each \textbf{primary oocyte} undergoes meiosis with \textbf{unequal cytokinesis}: meiosis I gives one large secondary oocyte and one small polar body; meiosis II gives one large ovum and further polar bodies. Only \textbf{1 functional ovum} results per primary oocyte. So 100 primary oocytes $\to \mathbf{100}$ \textbf{ova} (plus polar bodies that degenerate).\\
\textbf{Explanation:} in oogenesis the cytoplasm is conserved in a single large cell so that the ovum contains abundant food reserves for the early embryo; in spermatogenesis all four products are small, motile gametes, maximising the chance that one reaches the egg.\\
(b) \textbf{Hormonal control of the testes:} the hypothalamus secretes \cle{GnRH}, which stimulates the anterior pituitary to release \cle{FSH} and \cle{LH}. FSH acts on the \textbf{Sertoli cells} of the seminiferous tubules, promoting spermatogenesis and the nourishment of developing sperm. LH stimulates the \textbf{Leydig (interstitial) cells} to secrete \cle{testosterone}, which is essential for sperm maturation, the development of secondary sexual characteristics and libido. Testosterone exerts \textbf{negative feedback} on the hypothalamus and pituitary, keeping its level stable.\\
(c) \textbf{Destruction of Leydig cells:} testosterone secretion falls drastically. Consequences: spermatogenesis is impaired or stops (infertility), because testosterone is required for sperm maturation; secondary sexual characteristics may regress; libido decreases. Because the negative feedback is removed, blood levels of LH (and FSH) \textbf{rise} — a diagnostic sign that the failure is in the testis, not the pituitary.
\end{exoresolu}

\courssub{Skill 6: Analysing Growth Curves and Growth Patterns}

\begin{methode}[Interpreting a growth curve]
\begin{enumerate}
\item \textbf{Identify the type of curve:}
\begin{itemize}
\item \cle{Sigmoid (S-shaped)} curve: slow growth (lag phase) $\to$ rapid exponential growth $\to$ stationary phase — typical of annual plants and many organs.
\item \cle{Intermittent/stepped} curve: growth in steps after each moult (ecdysis) — typical of insects and arthropods with an exoskeleton.
\item Human growth curve: rapid in infancy, steady in childhood, a \textbf{pubertal growth spurt}, then a plateau at adulthood.
\end{itemize}
\item \textbf{Compute the growth rate} between two times: $\text{rate} = \dfrac{\text{size at } t_2 - \text{size at } t_1}{t_2 - t_1}$.
\item \textbf{Explain each phase} biologically (cell division, cell elongation, differentiation; moulting; hormonal spurts).
\item \textbf{Link to hormones when asked:} in plants, auxins, gibberellins, cytokinins; in humans, growth hormone (pituitary) and thyroid hormones, plus sex hormones at puberty.
\end{enumerate}
\end{methode}

\begin{popfigure}
\begin{center}
\begin{tikzpicture}[scale=0.9]
\begin{axis}[
  width=10.5cm, height=6.2cm,
  xlabel={time (weeks)}, ylabel={height (cm)},
  xmin=0, xmax=9, ymin=0, ymax=70,
  xtick={1,2,3,4,5,6,7,8},
  grid=both, grid style={popline, very thin},
  legend style={at={(0.03,0.97)}, anchor=north west, font=\small}
]
\addplot[very thick, popblue, mark=*, mark size=1.6pt] coordinates {(1,4) (2,8) (3,18) (4,34) (5,50) (6,58) (7,60) (8,60)};
\addlegendentry{maize seedling}
\addplot[very thick, poporange, domain=0.4:8.6, samples=60] {60/(1+exp(-(x-4)*1.3))};
\addlegendentry{ideal sigmoid}
\end{axis}
\end{tikzpicture}
\end{center}
\legende{Growth of a maize seedling: lag phase, exponential phase and stationary phase form a sigmoid (S-shaped) curve.}
\end{popfigure}

\begin{exoresolu}[Growth of a maize seedling and of a locust]
The height of a maize seedling was measured weekly: week 1: 4 cm; week 2: 8 cm; week 3: 18 cm; week 4: 34 cm; week 5: 50 cm; week 6: 58 cm; week 7: 60 cm; week 8: 60 cm. (a) Calculate the growth rate between weeks 3 and 5, and between weeks 6 and 8. (b) Describe and explain the shape of the curve. (c) A locust's growth curve is stepped, not smooth. Explain. (d) Name one hormone that promotes stem elongation in maize and state its action.
\tcblower
\textbf{Solution.}\\
(a) \textbf{Growth rate, weeks 3--5:} $\dfrac{50 - 18}{5 - 3} = \dfrac{32}{2} = 16\ \mathrm{cm/week}$.\\
\textbf{Growth rate, weeks 6--8:} $\dfrac{60 - 58}{8 - 6} = \dfrac{2}{2} = 1\ \mathrm{cm/week}$.\\
(b) \textbf{Shape:} the curve is \cle{sigmoid} (S-shaped). \textbf{Lag phase} (weeks 1--2): growth is slow because the seedling has few cells and is still establishing leaves and roots. \textbf{Exponential (log) phase} (weeks 2--5): rapid growth by active cell division in meristems and cell elongation, with photosynthesis supplying abundant food — the rate peaks at $16\ \mathrm{cm/week}$. \textbf{Stationary phase} (weeks 6--8): growth slows and stops as the plant reaches maturity; cells differentiate and resources are diverted to flowering and fruiting.\\
(c) \textbf{Stepped curve of the locust:} insects have a rigid \cle{exoskeleton} (cuticle) that cannot expand. Growth in size is only possible when the old cuticle is shed (\cle{ecdysis}/moulting) and the new, soft cuticle expands before hardening. Between moults, size is constant, so the curve rises in steps at each moult. Moulting is controlled by the hormone \textbf{ecdysone}.\\
(d) \textbf{Gibberellins} promote stem elongation: they stimulate both cell division and especially \textbf{cell elongation} in the internodes. (Auxins also promote cell elongation in shoots by stimulating cell wall loosening.)
\end{exoresolu}

\courssub{Skill 7: Evaluating Birth Control Methods and Life Cycles}

\begin{methode}[Answering questions on contraception and life cycles]
\begin{enumerate}
\item \textbf{Classify the method} before discussing it: natural (abstinence, rhythm/calendar method), barrier (condom, diaphragm), hormonal (pill, implant, injectable), intrauterine device (IUD), surgical/permanent (vasectomy, tubal ligation).
\item \textbf{State the mechanism precisely:} e.g. pill = negative feedback preventing ovulation; condom = physical barrier preventing sperm meeting egg (and reduces sexually transmitted infections); IUD = prevents implantation; vasectomy = cutting and tying the vas deferens.
\item \textbf{Evaluate:} reliability, advantages, disadvantages, protection (or not) against STIs such as HIV.
\item \textbf{For life cycles:} in \cle{alternation of generations} (ferns, mosses), identify the \textbf{sporophyte} ($2n$, produces spores by meiosis) and the \textbf{gametophyte} ($n$, produces gametes by mitosis); fertilization restores $2n$. In \cle{metagenesis} (e.g. Obelia, jellyfish), an asexual phase (polyp) alternates with a sexual phase (medusa).
\item \textbf{Always place meiosis correctly:} meiosis occurs at spore formation in plants, at gamete formation in animals.
\end{enumerate}
\end{methode}

\begin{popfigure}
\begin{center}
\begin{tikzpicture}[scale=1.0, every node/.style={font=\small}]
\node[draw, thick, popblue, rounded corners, align=center, minimum width=3.2cm] (sporo) at (0,2.6) {\textbf{Sporophyte} ($2n$)\\ fern plant};
\node[draw, thick, popgreen, rounded corners, align=center, minimum width=3.2cm] (spore) at (4.6,2.6) {spores ($n$)};
\node[draw, thick, popgreen, rounded corners, align=center, minimum width=3.2cm] (gameto) at (4.6,-1.2) {\textbf{Gametophyte} ($n$)\\ prothallus};
\node[draw, thick, poporange, rounded corners, align=center, minimum width=3.2cm] (gametes) at (0,-1.2) {gametes ($n$)\\ sperm + egg};
\draw[-latex, very thick, popdark] (sporo) -- (spore) node[midway, above]{meiosis};
\draw[-latex, very thick, popdark] (spore) -- (gameto) node[midway, right]{mitosis};
\draw[-latex, very thick, popdark] (gameto) -- (gametes) node[midway, above]{mitosis};
\draw[-latex, very thick, popdark] (gametes) -- (sporo) node[midway, left, align=center]{fertilization\\ zygote ($2n$)};
\end{tikzpicture}
\end{center}
\legende{Alternation of generations in a fern: a diploid sporophyte alternates with a haploid gametophyte; meiosis occurs at spore formation.}
\end{popfigure}

\begin{exoresolu}[Contraception and the fern life cycle]
(a) A couple in Yaoundé asks a health worker for advice. Compare the condom and the combined pill under the headings: mechanism, reliability, and protection against sexually transmitted infections. (b) Explain why vasectomy is considered a permanent method and why it does not affect a man's sexual characteristics. (c) In a fern, state the ploidy of the sporophyte, the spores, the gametophyte (prothallus) and the zygote, and state where meiosis occurs.
\tcblower
\textbf{Solution.}\\
(a) \textbf{Condom:} \emph{Mechanism} — a thin sheath worn over the penis acts as a \textbf{physical barrier}, preventing sperm from entering the vagina. \emph{Reliability} — high when used correctly every time, but it can fail through tearing or incorrect use. \emph{STI protection} — \textbf{yes}: it is the only common method that greatly reduces transmission of HIV and other sexually transmitted infections.\\
\textbf{Combined pill:} \emph{Mechanism} — synthetic oestrogen and progesterone cause \textbf{negative feedback} on the pituitary, inhibiting FSH and LH, so follicles do not mature and \textbf{ovulation is prevented}; the cervical mucus is also thickened. \emph{Reliability} — very high when taken daily without omission. \emph{STI protection} — \textbf{none}: it offers no barrier against infections.\\
\textbf{Advice:} for protection against both pregnancy and STIs, the condom (or both methods combined) is preferable.\\
(b) \textbf{Vasectomy:} the \cle{vas deferens} (sperm duct) on each side is cut and tied, so sperm can no longer leave the testes and appear in the semen. It is \textbf{permanent} because reconnection is difficult and usually not attempted. Sexual characteristics are \textbf{unaffected} because the Leydig cells still secrete \textbf{testosterone directly into the blood} (not through the vas deferens); erection, libido and ejaculation of seminal fluid (from the prostate and seminal vesicles) continue normally — only sperm are absent from the ejaculate.\\
(c) \textbf{Fern life cycle (alternation of generations):}
\begin{itemize}
\item Sporophyte (the familiar fern plant): \textbf{diploid} ($2n$).
\item Spores: produced in sporangia on the fronds \textbf{by meiosis}, so \textbf{haploid} ($n$).
\item Gametophyte (prothallus): develops from a spore by mitosis, so \textbf{haploid} ($n$); it bears antheridia and archegonia producing gametes ($n$) by mitosis.
\item Zygote: formed by fusion of sperm and egg, so \textbf{diploid} ($2n$); it grows into a new sporophyte.
\end{itemize}
\textbf{Meiosis occurs at spore formation} in the sporangia of the sporophyte — this is the key event that halves the chromosome number and switches the cycle from the diploid to the haploid generation.
\end{exoresolu}

\begin{attention}[Frequent errors in GCE answers]
\begin{itemize}
\item Confusing \textbf{pollination} (transfer of pollen to the stigma) with \textbf{fertilization} (fusion of gamete nuclei) — they are distinct events.
\item Forgetting that the \textbf{endosperm is triploid} ($3n$), not diploid.
\item Saying ``LH causes menstruation'' — in fact it is the \textbf{fall of progesterone} that triggers menstruation; the LH surge triggers \textbf{ovulation}.
\item Claiming oogenesis produces four ova per meiosis — it produces \textbf{one ovum} plus polar bodies.
\item Placing meiosis at gamete formation in plants — in plants meiosis produces \textbf{spores}, not gametes.
\item Writing percentages without showing the formula and substitution — always show $\frac{\text{germinated}}{\text{total}}\times 100$.
\end{itemize}
\end{attention}

\newpage

\coursec{Exercises}

\courssub{I check my knowledge}

\begin{exercice}[MCQ: Asexual reproduction mechanisms]
Which of the following is \emph{not} a form of asexual reproduction in animals?
\begin{enumerate}
\item Budding in Hydra
\item Binary fission in Amoeba
\item Parthenogenesis in aphids
\item Conjugation in Paramecium
\end{enumerate}
\end{exercice}

\begin{exercice}[True/False with justification]
State whether each statement is true or false. Justify your answer briefly.
\begin{enumerate}
\item The pollen grain of a flowering plant is the male gamete.
\item Double fertilization occurs in all vascular plants.
\item The corpus luteum secretes mainly oestrogen during the luteal phase.
\item In human spermatogenesis, one primary spermatocyte produces four functional spermatozoa.
\end{enumerate}
\end{exercice}

\begin{exercice}[Definitions]
Define the following terms precisely:
\begin{enumerate}
\item Alternation of generations
\item Metagenesis
\item Epigeal germination
\item Menarche
\end{enumerate}
\end{exercice}

\begin{exercice}[Direct calculation: Growth rate]
A locust nymph measures 8 mm at the beginning of an instar and 12 mm at the end of the same instar, which lasts 5 days. Calculate the average daily growth rate in mm/day. If the growth follows a sigmoid curve, explain why the rate would not be constant throughout the instar.
\end{exercice}

\courssub{I practise}

\begin{exercice}[Labelling and hormonal events]
Draw a large, labelled diagram of the human female reproductive system (anterior view). On your diagram, indicate the site of fertilization and the site of implantation. Then, describe the hormonal events that occur from day 1 to day 14 of a typical 28-day menstrual cycle, naming the hormones, their sources and their main actions.
\end{exercice}

\begin{exercice}[Data analysis: Hormone profiles]
The table below shows the concentrations of four hormones in the blood of a woman over a 28-day menstrual cycle.

\begin{center}
\begin{tabular}{|c|cccc|}
\hline
Day & FSH (mIU/mL) & LH (mIU/mL) & Oestrogen (pg/mL) & Progesterone (ng/mL) \\
\hline
1 & 5.2 & 4.8 & 30 & 0.5 \\
4 & 6.1 & 5.0 & 45 & 0.4 \\
7 & 7.5 & 6.2 & 80 & 0.3 \\
10 & 8.0 & 7.5 & 150 & 0.3 \\
13 & 9.5 & 25.0 & 280 & 0.4 \\
14 & 10.0 & 45.0 & 300 & 0.5 \\
17 & 4.0 & 8.0 & 180 & 12.0 \\
20 & 3.5 & 5.5 & 150 & 25.0 \\
23 & 3.0 & 4.0 & 120 & 30.0 \\
26 & 2.8 & 3.5 & 80 & 20.0 \\
28 & 4.5 & 4.0 & 40 & 2.0 \\
\hline
\end{tabular}
\end{center}

\begin{enumerate}
\item Plot the four hormone curves on the same axes (use a suitable scale).
\item Deduce the most probable day of ovulation. Justify your answer using the data.
\item Explain the role of the LH surge in ovulation and the subsequent formation of the corpus luteum.
\end{enumerate}
\end{exercice}

\begin{exercice}[Comparison: Spermatogenesis vs Oogenesis]
Copy and complete the table below to compare spermatogenesis and oogenesis in humans.

\begin{center}
\begin{tabular}{|l|c|c|}
\hline
Feature & Spermatogenesis & Oogenesis \\
\hline
Site & & \\
Timing (when does it start/end?) & & \\
Number of functional gametes per meiosis & & \\
Hormonal control (key hormones) & & \\
Presence of a resting phase & & \\
\hline
\end{tabular}
\end{center}
\end{exercice}

\begin{exercice}[Experiment design: Oxygen and germination]
Design an investigation to show that oxygen is necessary for the germination of bean seeds (\emph{Phaseolus vulgaris}). Your answer must include:
\begin{enumerate}
\item A clear hypothesis.
\item The independent, dependent and two controlled variables.
\item A description of the experimental set-up (including a control).
\item The expected results and how they support the hypothesis.
\item One precaution to ensure validity.
\end{enumerate}
\end{exercice}

\courssub{I prepare for the GCE}

\begin{exercice}[Structured essay: Female reproductive system and menstrual cycle] (20 marks)
\begin{enumerate}
\item Draw a fully labelled diagram of the human female reproductive system (anterior view). [6 marks]
\item Describe the hormonal changes that occur during the follicular phase (days 1–14) of the menstrual cycle. Explain how these changes lead to ovulation. [8 marks]
\item Outline the role of progesterone in preparing the uterus for implantation. [6 marks]
\end{enumerate}
\end{exercice}

\begin{exercice}[Life cycle interpretation: Fern] (15 marks)
The diagram below shows the life cycle of a fern (\emph{Dryopteris}) with stages labelled A to F. (Imagine a typical alternation of generations diagram: A = mature sporophyte, B = sporangium, C = spores, D = gametophyte/prothallus, E = antheridia/archegonia, F = zygote/young sporophyte.)

\begin{enumerate}
\item Identify which stages are haploid (n) and which are diploid (2n). [4 marks]
\item State where meiosis occurs and where fertilization occurs. [3 marks]
\item Explain the significance of the alternation of generations in ferns, particularly in relation to genetic variation and dispersal. [8 marks]
\end{enumerate}
\end{exercice}

\begin{exercice}[Problem question: Growth disorders and birth control] (20 marks)
\begin{enumerate}
\item A 10-year-old child shows stunted growth (height well below the 3rd percentile). Using your knowledge of growth hormone (GH) and thyroxine (T4), suggest \emph{three} possible endocrine causes for this condition. For each cause, describe a diagnostic test that could confirm it. [12 marks]
\item A couple wishes to delay pregnancy for two years. They are considering the following methods: combined oral contraceptive pill, copper intrauterine device (IUD), male condom, and subdermal implant. For each method, state:
\begin{itemize}
\item The precise stage of reproduction it blocks.
\item One advantage and one disadvantage.
\end{itemize}
[8 marks]
\end{enumerate}
\end{exercice}

\newpage

\coursec{Solutions to the exercises}

\courssub{I check my knowledge}

\begin{corrige}[Exercise 1 — MCQ: Asexual reproduction mechanisms]
The correct answer is \textbf{d) Conjugation in Paramecium}.

\textbf{Justification:}
\begin{itemize}
\item \textbf{Budding in Hydra} is asexual: a small outgrowth of the body wall develops tentacles, then detaches as a new individual.
\item \textbf{Binary fission in Amoeba} is asexual: the nucleus divides by mitosis and the cytoplasm splits into two genetically identical daughter cells.
\item \textbf{Parthenogenesis in aphids} is asexual: unfertilized eggs develop into new females (a form of reproduction without fertilization).
\item \textbf{Conjugation in Paramecium} is \emph{not} asexual reproduction: two individuals exchange micronuclear material (a sexual process of genetic recombination). It does not by itself increase the number of individuals, and it involves meiosis and exchange of gametic nuclei.
\end{itemize}
\end{corrige}

\begin{corrige}[Exercise 2 — True/False with justification]
\begin{enumerate}
\item \textbf{False.} The pollen grain is the \emph{male gametophyte}; it \emph{contains} the male gametes (two sperm nuclei formed from the generative nucleus). The male gamete itself is the sperm nucleus, not the whole grain.
\item \textbf{False.} Double fertilization (one sperm nucleus fusing with the egg cell to form the zygote, the other with the polar nuclei to form the endosperm) is characteristic of \emph{angiosperms} (flowering plants) only. Gymnosperms and pteridophytes (ferns) are vascular plants but do not show double fertilization.
\item \textbf{False.} The corpus luteum secretes mainly \emph{progesterone} during the luteal phase (with some oestrogen). Oestrogen is secreted mainly by the developing Graafian follicles during the follicular phase.
\item \textbf{True.} In spermatogenesis, one primary spermatocyte undergoes meiosis I to give two secondary spermatocytes, then meiosis II to give four spermatids, all of which differentiate into four functional spermatozoa (no polar bodies are formed).
\end{enumerate}
\end{corrige}

\begin{corrige}[Exercise 3 — Definitions]
\begin{enumerate}
\item \textbf{Alternation of generations:} a life cycle in which a multicellular \cle{diploid} generation (the \cle{sporophyte}, which produces spores by meiosis) alternates regularly with a multicellular \cle{haploid} generation (the \cle{gametophyte}, which produces gametes by mitosis). Example: ferns, mosses.
\item \textbf{Metagenesis:} a particular alternation of generations in which a \emph{sexual} generation alternates with an \emph{asexual} generation that reproduces by budding, as in \emph{Obelia} (the polyp colony buds off medusae; the medusae reproduce sexually).
\item \textbf{Epigeal germination:} germination in which the \cle{hypocotyl} elongates and curves, carrying the cotyledons \emph{above} the soil surface (e.g. bean, \emph{Phaseolus vulgaris}).
\item \textbf{Menarche:} the first menstrual period (first menstruation) of a girl, marking the beginning of reproductive life at puberty (typically around 11–14 years).
\end{enumerate}
\end{corrige}

\begin{corrige}[Exercise 4 — Direct calculation: Growth rate]
\textbf{Calculation of the average daily growth rate:}
\[
\text{Growth rate} = \frac{\text{final length} - \text{initial length}}{\text{time}} = \frac{12\ \mathrm{mm} - 8\ \mathrm{mm}}{5\ \mathrm{days}} = \frac{4}{5} = 0.8\ \mathrm{mm/day}.
\]

\textbf{Why the rate is not constant over a sigmoid curve:}
A sigmoid (S-shaped) growth curve has three phases:
\begin{itemize}
\item a \textbf{lag phase} at the start, when growth is slow (cells are few, the nymph is adapting and feeding little);
\item a \textbf{log (exponential) phase}, when growth is rapid and the rate is maximal;
\item a \textbf{stationary/deceleration phase}, when growth slows as the instar approaches its end (the cuticle hardens before ecdysis, food intake falls).
\end{itemize}
Because the instantaneous rate (slope of the curve) rises then falls, the value $0.8\ \mathrm{mm/day}$ is only an \emph{average}; the actual daily rate is lower than this at the beginning and end of the instar and higher in the middle.
\end{corrige}

\courssub{I practise}

\begin{corrige}[Exercise 5 — Labelling and hormonal events]
\textbf{Diagram} (anterior view) — the candidate's drawing must show and label: the two \textbf{ovaries}, the \textbf{Fallopian tubes (oviducts)} with fimbriae, the \textbf{uterus} (with endometrium and myometrium), the \textbf{cervix}, the \textbf{vagina}. The \textbf{site of fertilization} is the \emph{ampulla of the Fallopian tube} (outer third of the oviduct); the \textbf{site of implantation} is the \emph{endometrium of the uterus} (usually the upper posterior wall).

\textbf{Hormonal events, day 1 to day 14 (follicular phase):}
\begin{enumerate}
\item \textbf{Days 1–5 (menstruation):} the previous corpus luteum has degenerated, so progesterone and oestrogen levels are low. The endometrium breaks down and is shed. Low oestrogen/progesterone removes negative feedback on the hypothalamus and pituitary.
\item \textbf{From day 1:} the hypothalamus secretes \textbf{GnRH}, stimulating the anterior pituitary to release \textbf{FSH} (and some LH).
\item \textbf{FSH} (source: anterior pituitary) stimulates the growth of several ovarian follicles; the developing follicles secrete \textbf{oestrogen}.
\item \textbf{Oestrogen} (source: growing follicles) causes repair and thickening (proliferation) of the endometrium; at moderate levels it exerts negative feedback on FSH (so only one dominant follicle survives).
\item \textbf{Around days 12–13:} oestrogen reaches a high, sustained peak and switches to \emph{positive feedback} on the pituitary, triggering a surge of \textbf{LH} (and a smaller FSH surge).
\item \textbf{Day 14:} the LH surge causes \textbf{ovulation} — rupture of the mature Graafian follicle and release of the secondary oocyte.
\end{enumerate}
\end{corrige}

\begin{corrige}[Exercise 6 — Data analysis: Hormone profiles]
\textbf{1. Plotting.} Use days 1–28 on the x-axis. Because the units and magnitudes differ, plot LH and FSH on one scale (0–50 mIU/mL) and use a secondary y-axis (or a separate scaled plot) for oestrogen (0–300 pg/mL) and progesterone (0–30 ng/mL). Key features the graph must show: a sharp \textbf{LH peak at day 14} (45 mIU/mL), a smaller FSH peak around days 13–14, an \textbf{oestrogen peak at day 14} (300 pg/mL) just before/at ovulation, and a \textbf{progesterone peak around day 23} (30 ng/mL) in the luteal phase.

\textbf{2. Day of ovulation: day 14.} Justification from the data:
\begin{itemize}
\item LH rises sharply from 7.5 (day 10) to 45.0 mIU/mL (day 14) — the \cle{LH surge} directly triggers ovulation.
\item Oestrogen peaks at 300 pg/mL on day 14, then falls (180 by day 17) — the oestrogen peak precedes and triggers the LH surge.
\item Progesterone is low ($\leq 0.5$ ng/mL) until day 14, then rises steeply (12.0 by day 17) — progesterone is secreted by the corpus luteum, which can only form \emph{after} ovulation.
\end{itemize}

\textbf{3. Role of the LH surge:} the LH surge (a) completes the maturation of the Graafian follicle and the oocyte (resumption of meiosis I), (b) weakens the follicle wall (enzymatic digestion) causing its rupture and release of the oocyte — ovulation; (c) transforms the remaining granulosa and theca cells of the ruptured follicle into the \cle{corpus luteum} (luteinization), which then secretes progesterone (and oestrogen) to maintain the endometrium during the luteal phase.
\end{corrige}

\begin{corrige}[Exercise 7 — Comparison: Spermatogenesis vs Oogenesis]
\begin{center}
\begin{tabularx}{\linewidth}{|l|X|X|}
\hline
\rowcolor{popblueL} Feature & Spermatogenesis & Oogenesis \\
\hline
Site & Seminiferous tubules of the testes & Ovaries (follicles) \\
\hline
Timing & Starts at puberty; continuous throughout life & Begins before birth (oocytes arrest in prophase I); resumes at puberty; ends at menopause \\
\hline
Functional gametes per meiosis & 4 functional spermatozoa & 1 functional ovum + 2–3 polar bodies (which degenerate) \\
\hline
Hormonal control & FSH (Sertoli cells), LH (testosterone from Leydig cells), testosterone & FSH, LH, oestrogen, progesterone (cyclic) \\
\hline
Resting phase & Absent — continuous process & Present — long arrests in prophase I and metaphase II \\
\hline
\end{tabularx}
\end{center}
\end{corrige}

\begin{corrige}[Exercise 8 — Experiment design: Oxygen and germination]
\textbf{1. Hypothesis:} oxygen is necessary for the germination of bean seeds; seeds deprived of oxygen will not germinate.

\textbf{2. Variables:}
\begin{itemize}
\item \textbf{Independent variable:} presence/absence of oxygen.
\item \textbf{Dependent variable:} germination (percentage of seeds germinated, or radicle length after a fixed time).
\item \textbf{Controlled variables (any two):} same species/age/number of seeds; same temperature; same water availability; same light conditions; same duration of the experiment.
\end{itemize}

\textbf{3. Set-up:}
\begin{itemize}
\item \textbf{Control (Set-up A):} soak bean seeds, place them on moist cotton wool in an open flask exposed to air at room temperature.
\item \textbf{Experiment (Set-up B):} place identical soaked seeds in a flask of freshly \emph{boiled and cooled} water (boiling expels dissolved air), and cover the water surface with a layer of oil to prevent oxygen diffusing back in. (Alternatively, use alkaline pyrogallol in a sealed flask to absorb oxygen.)
\item Keep both flasks in the same place (same temperature and light) for 4–5 days.
\end{itemize}

\textbf{4. Expected results:} seeds in Set-up A (with oxygen) germinate — radicles emerge and grow; seeds in Set-up B (without oxygen) do not germinate. Since the only difference is the availability of oxygen, the result supports the hypothesis that oxygen is necessary for germination (needed for aerobic respiration, which releases the energy for growth of the embryo).

\textbf{5. Precaution:} use a large number of seeds (e.g. 20 per set-up) and repeat the experiment, to avoid the effect of dead or non-viable seeds; also check that seeds in B are still alive by transferring them to air afterwards — they should then germinate, proving lack of oxygen (not death) prevented germination.
\end{corrige}

\courssub{I prepare for the GCE}

\begin{corrige}[Exercise 9 — Structured essay: Female reproductive system and menstrual cycle (20 marks)]
\textbf{(a) Diagram [6 marks].} Examiner expectations: large clear drawing in pencil, ruled label lines, no shading. Award 1 mark for each correctly drawn and labelled structure, up to 6: ovary, Fallopian tube (oviduct) with fimbriae, uterus, endometrium/myometrium, cervix, vagina. Deduct for labels touching the wrong structure or for a diagram too small.

\textbf{(b) Hormonal changes in the follicular phase [8 marks] — indicative mark scheme (1 mark each):}
\begin{enumerate}
\item At the start of the cycle, progesterone and oestrogen levels are low because the previous corpus luteum has degenerated.
\item Low hormone levels remove negative feedback, so the hypothalamus secretes GnRH.
\item GnRH stimulates the anterior pituitary to secrete FSH (and LH).
\item FSH stimulates the growth and maturation of ovarian follicles.
\item The developing follicles secrete increasing amounts of oestrogen.
\item Oestrogen repairs and thickens the endometrium (proliferation) after menstruation.
\item High, sustained oestrogen exerts positive feedback, causing a surge of LH (and FSH) around days 13–14.
\item The LH surge triggers ovulation: rupture of the Graafian follicle and release of the secondary oocyte (day 14).
\end{enumerate}

\textbf{(c) Role of progesterone in preparing the uterus [6 marks] — 1 mark each:}
\begin{enumerate}
\item Progesterone is secreted by the corpus luteum after ovulation.
\item It maintains and further thickens the endometrium (secretory phase).
\item It stimulates the endometrial glands to secrete nutrients (uterine milk) for the early embryo.
\item It increases the blood supply/vascularization of the endometrium.
\item It inhibits uterine (myometrial) contractions, preventing expulsion of the embryo.
\item It inhibits FSH/LH release (negative feedback), preventing development of new follicles during a possible pregnancy; it also thickens cervical mucus, blocking sperm entry.
\end{enumerate}
\end{corrige}

\begin{corrige}[Exercise 10 — Life cycle interpretation: Fern (15 marks)]
\textbf{(a) Ploidy of the stages [4 marks] — 1 mark each:}
\begin{itemize}
\item \textbf{Diploid (2n):} A (mature sporophyte), B (sporangium — its tissues are part of the sporophyte), F (zygote and young sporophyte).
\item \textbf{Haploid (n):} C (spores), D (gametophyte/prothallus), E (antheridia and archegonia, and the gametes they produce).
\end{itemize}
(Award 4 marks for all six correct; 2–3 marks for partial correct classification.)

\textbf{(b) Sites of meiosis and fertilization [3 marks]:}
\begin{itemize}
\item \textbf{Meiosis} occurs in the \textbf{sporangium (B)}: spore mother cells divide by meiosis to produce haploid spores (C). [2 marks]
\item \textbf{Fertilization} occurs on/in the \textbf{gametophyte (D)}, when a sperm swims from an antheridium to the egg in an archegonium (E), forming the diploid zygote (F). [1 mark]
\end{itemize}

\textbf{(c) Significance of alternation of generations [8 marks] — indicative points:}
\begin{itemize}
\item \textbf{Genetic variation (4 marks):} meiosis in the sporangium produces genetically varied spores (independent assortment and crossing over); the gametophyte generation produces gametes, and \emph{sexual} fusion of gametes from different prothalli combines genes from two parents, increasing variation; variation provides raw material for natural selection and adaptation.
\item \textbf{Dispersal (3 marks):} the sporophyte produces large numbers of tiny, light spores that are dispersed by wind over long distances, allowing ferns to colonize new, moist habitats (e.g. the wet slopes of Mount Cameroon); spores are resistant and can survive unfavourable conditions.
\item \textbf{Life-cycle efficiency (1 mark):} the dominant sporophyte is adapted to land (vascular tissue, roots, leaves) while the small gametophyte ensures sexual reproduction where water is available for sperm to swim — each generation exploits the conditions suited to it.
\end{itemize}
\end{corrige}

\begin{corrige}[Exercise 11 — Problem question: Growth disorders and birth control (20 marks)]
\textbf{(a) Three possible endocrine causes of stunted growth [12 marks — 4 per cause: 2 for the cause, 2 for the test]:}
\begin{enumerate}
\item \textbf{Growth hormone (GH) deficiency} (hyposecretion by the anterior pituitary, e.g. pituitary dwarfism). \emph{Diagnostic test:} measure blood GH level, which will be abnormally low; since GH is secreted in pulses, use a \emph{stimulation test} (e.g. insulin or arginine provocation) — a deficient child shows no rise in GH; low IGF-1 level also confirms it.
\item \textbf{Primary hypothyroidism (thyroxine/T4 deficiency)} — thyroxine is needed for normal growth and for the action of GH; deficiency in childhood causes stunted growth (cretinism if congenital). \emph{Diagnostic test:} measure blood T4 (low) and TSH (high in primary hypothyroidism, since negative feedback is removed).
\item \textbf{Secondary hypothyroidism (TSH deficiency)} — the pituitary fails to stimulate the thyroid, leading to low thyroxine and stunted growth. \emph{Diagnostic test:} low T4 \emph{with} low or normal TSH (distinguishing it from primary hypothyroidism); a TRH stimulation test shows whether the pituitary can respond.
\end{enumerate}
(Any three coherent endocrine causes with a matching test are accepted; e.g. GH-receptor insensitivity (Laron syndrome) with a GH/IGF-1 assay showing high GH but low IGF-1.)

\textbf{(b) Birth control methods [8 marks — 2 per method]:}
\begin{itemize}
\item \textbf{Combined oral contraceptive pill:} blocks \emph{ovulation} (oestrogen + progesterone inhibit FSH/LH by negative feedback, so no follicle matures); also thickens cervical mucus. \emph{Advantage:} very effective when taken correctly and reversible. \emph{Disadvantage:} must be taken daily; possible side effects (nausea, weight gain, increased risk of blood clots); no protection against STIs.
\item \textbf{Copper IUD:} acts mainly \emph{before fertilization} — copper ions are spermicidal and the device causes a uterine reaction hostile to sperm and to implantation. \emph{Advantage:} long-acting (several years), no daily action needed, hormone-free. \emph{Disadvantage:} can cause heavier/more painful periods; requires insertion by trained personnel; no protection against STIs.
\item \textbf{Male condom:} blocks the \emph{deposition of sperm in the vagina} (physical barrier preventing sperm from reaching the egg). \emph{Advantage:} cheap, available without prescription, and protects against STIs including HIV. \emph{Disadvantage:} can tear or slip if misused; must be used correctly at every intercourse.
\item \textbf{Subdermal implant:} releases progestogen which blocks \emph{ovulation} and thickens cervical mucus (blocking sperm passage). \emph{Advantage:} very effective and long-acting (3–5 years), reversible on removal. \emph{Disadvantage:} requires minor surgery for insertion/removal; may cause irregular bleeding; no protection against STIs.
\end{itemize}
\textbf{Examiner expectations:} for full marks the candidate must state the \emph{precise stage blocked} (ovulation, sperm transport, fertilization or implantation), not merely "prevents pregnancy", and give one realistic advantage and disadvantage per method.
\end{corrige}

\newpage

\coursec{Summary sheet}

\begin{popfigure}
\begin{tikzpicture}[grow cyclic, mindmap, concept color=popblue,
level 1/.append style={level distance=3.6cm, sibling angle=51.4},
level 2/.append style={level distance=2.1cm, sibling angle=40},
every node/.style={concept, text=white, font=\small},
root concept/.append style={font=\bfseries}]
\node[root concept]{Reproduction}
child[concept color=popgreen]{node{Asexual reproduction}
  child{node{Fission, budding, fragmentation}}
  child{node{Spores, vegetative propagation}}
  child{node{Parthenogenesis}}}
child[concept color=poporange]{node{Sexual reproduction in plants}
  child{node{Pollen \& ovule formation}}
  child{node{Pollination \& double fertilization}}
  child{node{Seed, dispersal, germination}}}
child[concept color=poppurple]{node{Male system}
  child{node{Testes, ducts, glands}}
  child{node{Spermatogenesis}}
  child{node{FSH, LH, testosterone}}}
child[concept color=poppink]{node{Female system}
  child{node{Ovaries, oviducts, uterus}}
  child{node{Oogenesis \& menstrual cycle}}
  child{node{FSH, LH, oestrogen, progesterone}}}
child[concept color=popgold]{node{Fertilization to birth}
  child{node{Zygote, cleavage, implantation}}
  child{node{Placenta, hCG}}
  child{node{Birth process}}}
child[concept color=popdark]{node{Life cycles \& birth control}
  child{node{Alternation of generations}}
  child{node{Metagenesis}}
  child{node{Contraception methods}}}
child[concept color=popblue!70!black]{node{Growth}
  child{node{Growth curves \& patterns}}
  child{node{Hormonal control: GH, thyroxine}}};
\end{tikzpicture}
\legende{Mind map of the chapter ``Reproduction'': the nine major blocks of the GCE A-Level syllabus.}
\end{popfigure}

\begin{bilan}
\textbf{Key definitions}
\begin{itemize}
\item \cle{Reproduction}: biological process by which organisms produce new individuals, ensuring the continuity of the species.
\item \cle{Asexual reproduction}: production of offspring from a single parent, without gametes or fertilization; offspring are genetically identical to the parent (\cle{clones}). Mechanisms: binary fission (Amoeba, bacteria), budding (Hydra, yeast), fragmentation (Spirogyra, planarians), sporulation (ferns, fungi such as Mucor), vegetative propagation (runners, tubers, rhizomes, bulbs, cuttings -- e.g. cassava, cocoyam and plantain in Cameroon).
\item \cle{Parthenogenesis}: development of an unfertilized egg into a new individual (aphids, honeybees -- drones develop from unfertilized eggs).
\item \cle{Sexual reproduction}: fusion of two haploid gametes (\cle{fertilization}) to form a diploid \cle{zygote}; meiosis and fertilization together produce genetic variation.
\item \cle{Pollination}: transfer of pollen grains from the anther to the stigma (self-pollination vs cross-pollination; agents: wind, insects, water, birds).
\item \cle{Double fertilization} (flowering plants): one male nucleus fuses with the egg cell $\rightarrow$ diploid zygote ($2n$); the second male nucleus fuses with the two polar nuclei $\rightarrow$ triploid \cle{endosperm} nucleus ($3n$), which develops into the nutritive tissue of the seed.
\item \cle{Seed dispersal}: scattering of seeds away from the parent plant, reducing competition; agents: wind, water, animals, explosive (self-dispersal) mechanisms.
\item \cle{Germination}: resumption of growth of the embryo when water, oxygen and a suitable temperature are present; \cle{epigeal} germination (cotyledons carried above ground, e.g. bean) vs \cle{hypogeal} germination (cotyledons remain below ground, e.g. maize).
\item \cle{Spermatogenesis}: continuous production of spermatozoa in the seminiferous tubules of the testes from puberty onwards: spermatogonia ($2n$) $\rightarrow$ primary spermatocytes $\rightarrow$ (meiosis I) secondary spermatocytes ($n$) $\rightarrow$ (meiosis II) spermatids $\rightarrow$ spermatozoa.
\item \cle{Oogenesis}: production of ova in the ovaries; begins before birth (oocytes arrested in meiosis I), resumes at puberty; usually one secondary oocyte is released per cycle at \cle{ovulation}, and meiosis II is completed only if fertilization occurs.
\item \cle{Menstrual cycle}: approximately 28 days; phases: menstrual (days 1--5), follicular/proliferative (days 6--13), ovulation (about day 14), luteal/secretory (days 15--28).
\item \cle{Fertilization} (human): fusion of sperm and ovum nuclei, normally in the oviduct (Fallopian tube), restoring the diploid number.
\item \cle{Implantation}: embedding of the blastocyst in the endometrium of the uterus, about 6--7 days after fertilization.
\item \cle{Alternation of generations}: alternation of a haploid \cle{gametophyte} generation (produces gametes by mitosis) and a diploid \cle{sporophyte} generation (produces spores by meiosis) in the life cycle of a plant (mosses, ferns).
\item \cle{Metagenesis}: alternation of a sexual generation and an asexual generation in the life cycle of an animal (e.g. Obelia: polyp colony reproduces asexually to release medusae, which reproduce sexually).
\item \cle{Growth}: irreversible increase in size, cell number and dry mass, accompanied by cell division, cell enlargement and differentiation.
\end{itemize}

\textbf{Hormonal control -- the essentials}
\begin{itemize}
\item \textbf{Male}: hypothalamus secretes \cle{GnRH} $\rightarrow$ anterior pituitary: \cle{FSH} stimulates Sertoli cells (support of spermatogenesis); \cle{LH} (= ICSH) stimulates the interstitial (Leydig) cells $\rightarrow$ \cle{testosterone} (development of secondary sexual characteristics, maintenance of spermatogenesis; negative feedback on GnRH, FSH and LH).
\item \textbf{Female}: \cle{FSH} $\rightarrow$ growth of Graafian follicles $\rightarrow$ \cle{oestrogen} (repair and thickening of the endometrium; high level triggers the LH surge); \cle{LH} surge $\rightarrow$ \cle{ovulation} and formation of the corpus luteum $\rightarrow$ \cle{progesterone} (maintains the endometrium; negative feedback on FSH and LH). If there is no pregnancy, the corpus luteum degenerates $\rightarrow$ fall of oestrogen and progesterone $\rightarrow$ breakdown of the endometrium $\rightarrow$ menstruation.
\item \textbf{Pregnancy}: \cle{hCG} (secreted by the embryo, then the placenta) maintains the corpus luteum; later the placenta itself secretes oestrogen and progesterone, which maintain the uterus and inhibit FSH/LH (no further ovulation). \cle{Oxytocin} triggers uterine contractions at birth (positive feedback); \cle{prolactin} stimulates milk production after birth.
\item \textbf{Growth}: \cle{GH} (growth hormone, anterior pituitary) stimulates protein synthesis and growth of bones and soft tissues; \cle{thyroxine} (thyroid) is essential for normal growth, development and metamorphosis; sex hormones (testosterone, oestrogen) cause the pubertal growth spurt and later the closure of the epiphyses (end of growth in height).
\end{itemize}

\textbf{Development sequence (learn the order!)}
\[
\text{Zygote} \rightarrow \text{cleavage} \rightarrow \text{morula} \rightarrow \text{blastocyst} \rightarrow \text{implantation (day 6--7)} \rightarrow \text{embryo} \rightarrow \text{foetus (from week 8)} \rightarrow \text{birth}
\]

\textbf{Birth control methods}
\begin{itemize}
\item Natural: abstinence, rhythm/calendar method (avoiding intercourse around ovulation), withdrawal -- all relatively unreliable.
\item Barrier: condom (also protects against STIs including HIV), diaphragm.
\item Hormonal: contraceptive pill (oestrogen + progesterone inhibit FSH/LH and block ovulation), implants, injections.
\item Intra-uterine device (IUD): prevents implantation/fertilization.
\item Surgical (permanent): vasectomy (cutting of the sperm ducts) in the male, tubal ligation (cutting of the oviducts) in the female.
\end{itemize}

\textbf{Growth curves and patterns}
\begin{itemize}
\item Sigmoid (S-shaped) curve in annual plants and many organs: lag phase $\rightarrow$ exponential (log) phase $\rightarrow$ stationary phase.
\item Human growth curve: rapid growth in infancy, slow steady growth in childhood, a spurt at puberty, then growth stops in adulthood.
\item Intermittent (stepwise) growth in arthropods due to \cle{ecdysis} (moulting of the rigid exoskeleton).
\item Measurements: length/height, fresh mass, and \cle{dry mass} (most accurate -- it excludes variations in water content).
\end{itemize}

\textbf{Common mistakes to avoid}
\begin{itemize}
\item Confusing pollination (transfer of pollen to the stigma) with fertilization (fusion of gamete nuclei).
\item Forgetting that the endosperm is \textbf{triploid} ($3n$), not diploid.
\item Saying FSH/LH act on the uterus: they act on the \textbf{ovary}; oestrogen and progesterone act on the uterus.
\item Placing ovulation at the start of the cycle: it occurs around \textbf{day 14} of a 28-day cycle.
\item Confusing oestrogen (follicular phase) with progesterone (luteal phase and pregnancy).
\item Confusing gametophyte (produces gametes, $n$) with sporophyte (produces spores, $2n$).
\item Mixing up epigeal (cotyledons above ground) and hypogeal (cotyledons below ground) germination.
\item Writing that the sperm fertilizes the ovum ``in the uterus'': fertilization normally occurs in the \textbf{oviduct}.
\end{itemize}

\textbf{GCE tips}
\begin{itemize}
\item Draw and label diagrams: flower, seed, testis/seminiferous tubule, ovary, menstrual cycle hormone graph -- marks are awarded for correct labels, proportions and a title.
\item In hormonal essays, always state: \textbf{gland -- hormone -- target organ -- effect -- feedback}.
\item For graph questions (hormone levels, growth curves), describe each phase and link it to a precise biological event (e.g. LH peak $\rightarrow$ ovulation).
\item Use precise vocabulary: ``implantation'' not ``attachment''; ``ovulation'' not ``release of the ovum into the uterus''.
\item Manage time: structured questions on reproduction carry many small marks -- answer every sub-question.
\end{itemize}
\end{bilan}

\findecours

\end{document}
